% concatenated from iclr2026_conference.tex
\documentclass{article}

\usepackage{iclr2026_conference,times}
\iclrfinalcopy


% Recommended, but optional, packages for figures and better typesetting:
\usepackage{microtype}
\usepackage{comment}

\usepackage{paralist,enumitem}
%\usepackage[table]{xcolor}

\usepackage{graphicx}
\usepackage{subfigure}
\usepackage{booktabs} % for professional tables
\usepackage{amsmath}
\usepackage{amsfonts}
\usepackage{nccmath}
\usepackage{amsthm}
\newtheorem{lemma}{Lemma}
\usepackage{textcomp}
\usepackage{upgreek}
\usepackage{diagbox} 
\usepackage{makecell}

\usepackage{bbm}
\usepackage{mathtools}
\usepackage{multirow} 
\DeclarePairedDelimiter\ceil{\lceil}{\rceil}
\DeclarePairedDelimiter\floor{\lfloor}{\rfloor}



\newtheorem{theorem}{Theorem}
\newtheorem{definition}{Definition}
%\newtheorem{corollary}{Corollary}[theorem]
\newtheorem{corollary}{Corollary}
\newtheorem{observation}{Observation}
\newtheorem{remark}{Remark}


\newtheorem{proposition}{Proposition}
\newcommand{\mysubsection}[1]{\medskip\noindent\textbf{#1}}
\newcommand{\PComplexity}{PTIME}
\newcommand{\NPComplexity}{NP}
\newcommand{\NPCompleteComplexity}{NP-complete}
\newcommand{\coNPComplexity}{co-NP}
\newcommand{\coNPCompleteComplexity}{co-NP-complete}
\newcommand{\StoPComplexity}{$\Sigma_{2}^{P}$} % S2P
\newcommand{\prodComplexity}{$\Pi_{2}^{P}$} % S2P

\newcommand{\StoPCompleteComplexity}{$\Sigma_{2}^{P}$-complete} % S2P
\newcommand{\countPComplexity}{\#P}
\newcommand{\countPCompleteComplexity}{\#P-complete}



\usepackage{amssymb} % for \mathbb
\usepackage{amsmath} % for \mathbb
\usepackage[T1]{fontenc}
\newcommand{\z}{\textbf{z}}
\newcommand{\x}{\textbf{x}}
\newcommand{\y}{\textbf{y}}
\newcommand{\w}{\textbf{w}}
\newcommand{\X}{\{1,\dots,n\}}
\newcommand{\yes}{\textit{Yes}}
\newcommand{\no}{\textit{No}}

\newcommand{\fml}[1]{{\mathcal{#1}}}
\newcommand{\mbf}[1]{\ensuremath\mathbf{#1}}
\newcommand{\mbb}[1]{\ensuremath\mathbb{#1}}


\renewcommand\theproposition{\arabic{proposition}}

\newcommand{\guy}[1]{\marginpar{\textcolor{orange}{GuyK: #1}}}
\newcommand{\guyamir}[1]{\marginpar{\textcolor{blue}{GuyA: #1}}}
\newcommand{\shahaf}[1]{\marginpar{\textcolor{green}{Shahaf: #1}}}
\newcommand{\todo}[1]{{\color{red}{\textbf{[TODO]} #1}}}




% hyperref makes hyperlinks in the resulting PDF.
% If your build breaks (sometimes temporarily if a hyperlink spans a page)
% please comment out the following usepackage line and replace
% \usepackage{icml2021} with \usepackage[nohyperref]{icml2021} above.
\usepackage{hyperref}
\usepackage{xspace}
% Attempt to make hyperref and algorithmic work together better:
\newcommand{\theHalgorithm}{\arabic{algorithm}}

\usepackage{algorithm}

%\RequirePackage{algorithm}
%\RequirePackage{algorithmic}
\usepackage{algpseudocode}
% Use the following line for the initial blind version submitted for review:
%\usepackage{icml2024}
%\usepackage{paralist}

\newcommand{\mySuff}{\text{suff}\xspace}
\newcommand{\gq}{\text{$G$-$Q$}\xspace}


% If accepted, instead use the following line for the camera-ready submission:
%\usepackage[accepted]{icml2024}

% The \icmltitle you define below is probably too long as a header.
% Therefore, a short form for the running title is supplied here:

\DeclareMathOperator*{\argmax}{argmax} % thin space, limits underneath in displays
\DeclareMathOperator*{\maxunder}{max} % thin space, limits underneath in displays
\DeclareMathOperator*{\minunder}{min} % thin space, limits underneath in displays


\DeclareMathOperator*{\argmin}{argmin} % thin space, limits underneath in displays

% if you need to pass options to natbib, use, e.g.:
%     \PassOptionsToPackage{numbers, compress}{natbib}
% before loading neurips_2024


% ready for submission




% to compile a preprint version, e.g., for submission to arXiv, add add the
% [preprint] option:
%     \usepackage[preprint]{neurips_2024}


% to compile a camera-ready version, add the [final] option, e.g.:
%     \usepackage[final]{neurips_2024}


% to avoid loading the natbib package, add option nonatbib:
%    \usepackage[nonatbib]{neurips_2024}


\usepackage[utf8]{inputenc} % allow utf-8 input
\usepackage[T1]{fontenc}    % use 8-bit T1 fonts
\usepackage{hyperref}       % hyperlinks
\usepackage{url}            % simple URL typesetting
\usepackage{booktabs}       % professional-quality tables
\usepackage{amsfonts}       % blackboard math symbols
\usepackage{nicefrac}       % compact symbols for 1/2, etc.
\usepackage{microtype}      % microtypography
\usepackage{xcolor}         % colors

\usepackage{cleveref}
\usepackage{thmtools}


%\title{Tractable Probabilistic Global Explanations for Decision Trees and Neural Networks}

%\title{\centering Formal Explanations, Unified: \\ A Complexity-Theoretic Approach}

\title{\centering Unifying Formal Explanations: \\ A Complexity-Theoretic Perspective}

%\title{\centering Unifying Formal XAI Problems: \\ A Complexity-Theoretic Approach}

%\title{\centering Formal Explanations, Unified: \\ A Complexity-Theoretic Approach}

%A Complexity-Theoretic Approach



%\title{Unifying Complexity-Theoretic Perspectives on Provable Explanations}



%\title{On the Many Facets of Sufficient Explanations}

%\author{shahaf.bassan }
%\date{June 2024}
\author{%
  Shahaf Bassan$^{1}$\thanks{Equal contribution.}\ \ , Xuanxiang Huang$^{2*}$\ \ , Guy Katz$^{1}$ \\
  The Hebrew University of Jerusalem$^{1}$, Nanyang Technological University$^{2}$\\
{\texttt{shahaf.bassan@mail.huji.ac.il, xuanxiang.huang@ntu.edu.sg,}} \\
\texttt{g.katz@mail.huji.ac.il}
}
\begin{document}


\maketitle


\begin{abstract}
    Previous work has explored the computational complexity of deriving two fundamental types of explanations for ML model predictions:  
\begin{inparaenum}[(i)] \item \emph{sufficient reasons}, which are subsets of input features that, when fixed, determine a prediction, and  
\item \emph{contrastive reasons}, which are subsets of input features that, when modified, alter a prediction\end{inparaenum}. Prior studies have examined these explanations in different contexts, such as non-probabilistic versus probabilistic frameworks and local versus global settings. In this study, we introduce a unified framework for analyzing these explanations, demonstrating that they can all be characterized through the minimization of a unified probabilistic value function. We then prove that the complexity of these computations is influenced by three key properties of the value function: \begin{inparaenum}[(i)] \item \emph{monotonicity}, \item \emph{submodularity}, and \item \emph{supermodularity} \end{inparaenum} --- which are three fundamental properties in \emph{combinatorial optimization}. Our findings uncover some counterintuitive results regarding the nature of these properties within the explanation settings examined. For instance, although the \emph{local} value functions do not exhibit monotonicity or submodularity/supermodularity whatsoever, we demonstrate that the \emph{global} value functions do possess these properties. This distinction enables us to prove a series of novel polynomial-time results for computing various explanations with provable guarantees in the global explainability setting, across a range of ML models that span the interpretability spectrum, such as neural networks, decision trees, and tree ensembles. In contrast, we show that even highly simplified versions of these explanations become NP-hard to compute in the corresponding local explainability setting.\end{abstract}


%This distinction enables us to prove a line of novel polynomial-time results for computing various explanations with provable guarantees in the global explainability setting for a range of ML models, while demonstrating these computational complexity results over model types that cover the interpretability spectrum, including neural networks, decision trees, and tree ensembles. Interestingly, we show from the other end, that computing even highly simplified versions of these same explanations in the corresponding local explainability setting is NP-hard.



%\begin{abstract}
%    Previous work has explored the computational complexity of deriving two fundamental types of explanations for ML model predictions:  
%\begin{inparaenum}[(i)] \item \emph{sufficient reasons}, which are subsets of input features that, when fixed, determine a prediction, and  
%\item \emph{contrastive reasons}, which are subsets of input features that, when modified, alter a prediction\end{inparaenum}. Prior work has examined these explanations in different contexts, such as non-probabilistic versus probabilistic frameworks and local versus global settings. In this study, we introduce a unified framework for analyzing these explanations, demonstrating that they can all be characterized through the minimization of a unified probabilistic value function. We then prove that the complexity of these computations is influenced by three key properties of the value function:  
%(i) \emph{saturation}, (ii) \emph{monotonicity}, and (iii) \emph{submodularity/supermodularity}. We analyze the computational complexity of generating these explanations across various model types that cover the interpretability spectrum, including neural networks, decision trees, and tree ensembles. Our findings reveal surprising and counterintuitive results. For instance, although the local value functions do not exhibit monotonicity or sub/super-modularity, we demonstrate that the global value functions possess these properties. This distinction allows us to prove novel and surprising polynomial-time results for computing different explanations with provable guarantees for a variety of ML models in the global setting, while showing that computing these same explanations in the local setting remains NP-hard.
%\end{abstract}

%\begin{abstract}
%Many existing explainability techniques focus on identifying a minimal subset of input features that can sufficiently determine a prediction. Previous studies have explored the complexity of finding such subsets, highlighting their tractability or intractability in various scenarios. In this work, we propose a unifying framework for a comprehensive analysis of this problem, addressing different notions of sufficiency—namely baseline, probabilistic, and robust—as well as both local and global computations. We demonstrate that a wide range of diverse settings can be analyzed through a single greedy minimization algorithm over an objective function, where the complexity depends on the monotonicity of the objective function and its convexity, which varies across the different sufficiency definitions (baseline, probabilistic, robust) and computation scopes (local, global). These findings introduce a range of new complexity results, both positive and negative, for computing sufficiency-based explanations across various types of ML models.\end{abstract}

%\begin{abstract}
%A recent study has revealed an intriguing complexity aspect of computing explanations for both decision trees and neural networks —-- in these models, finding the smallest subset of input features sufficient to determine a prediction is surprisingly more feasible \emph{globally} (across all inputs $\mathbf{x}$) than \emph{locally} (for a specific $\mathbf{x}$). However, the stringent definition of sufficiency used in that work, which considers any non-redundant feature as sufficient, is not practical for real-world applications, where it is highly likely that all features are at least somewhat relevant. In this paper, we modify this definition to a probabilistic notion of sufficiency and demonstrate that under this broader framework, global explanations remain significantly more tractable than local ones. Particularly, we prove that obtaining cardinally minimal global sufficient subsets with probabilistic assurances is feasible in polynomial time for decision trees and within bounded-error probabilistic polynomial time for neural networks, rendering the task tractable for both types of models. This is in stark contrast to the corresponding \emph{local} form of the task, which remains intractable for both models, even under probabilistic notions.\end{abstract}

%\section{Introduction}

%In order to prove the correctness of the algorithm I suggested it is enough to prove the following. We need to prove that the function $c(S,f)$ is \emph{super-modular}. Specifically, a function $c:2^N\to \mathbb{R}$ is super-modular iff for any subsets $S\subseteq S'\in N$ and any $i\not\in S'$ it holds that:

%\begin{equation}
%    c(S\cup \{i\})-c(S)\leq c(S'\cup \{i\})-c(S')
%\end{equation}

%Note that this is exactly the opposite than the definition of \emph{sub-modular} where it holds that:

%\begin{equation}
%    c(S\cup \{i\})-c(S)\geq c(S'\cup \{i\})-c(S')
%\end{equation}

%However, in our case, we need to prove \emph{super}-modularity. Given the algorithm I suggested, let us understand why super-modularity helps. Let us assume that we start with the entire set $S:=\{1,\ldots, n\}$. We at the first point 


%\begin{equation}
%    c(S\cup \{i\})-c(S)\geq \frac{1}{2^{|S'|-|S|}}(c(S'\cup \{i\})-c(S'))
%\end{equation}


%\begin{equation}
%    \begin{aligned}
%        \frac{1}{2^{|\Bar{S}|}}\cdot \frac{1}{2} \cdot [c(\emptyset,f|_{x_S,x_i=1},z)-c(\emptyset,f|_{x_S,x_i=0},z)]^2
%    \end{aligned}
%\end{equation}

%Let us denote some $S \subseteq S'$. In case $S=S'$ the proof is trivial, so let us hence assume that $S \subsetneq S'$. Then of course we have that $|S'|>|S|$ and hence $|\Bar{S}|>|\Bar{S'}|$.


%\begin{equation}
%    \begin{aligned}
%        \frac{1}{2^{|\Bar{S}|}}\cdot
%        \frac{1}{2^{|\Bar{S'}|-|\Bar{S}|}}\cdot \frac{1}{2} \cdot [c(\emptyset,f|_{x_{S'},x_i=1},z)-c(\emptyset,f|_{x_{S'},x_i=0},z)]^2
%    \end{aligned}
%\end{equation}

%We have that:

%\begin{equation}
%    \begin{aligned}
%    c(\emptyset,f|_{x_{S'},x_i=1},z)-c(\emptyset,f|_{x_{S'},x_i=1},z)<       2^{|\Bar{S'}|-|\Bar{S}|}[c(\emptyset,f|_{x_{S},x_i=1},z)-c(\emptyset,f|_{x_{S},x_i=1},z)]
%    \end{aligned}
%\end{equation}


%\begin{equation}
%    \begin{aligned}
%        c(\emptyset,f|_{x_{S'},x_i=1},z)=     \frac{|\{z\in\{0,1\}^{\overline{S}},z_S=x_S,z_i=1,f(z)=f(x)\}|}{|\{z\in\{0,1\}^{\overline{S}}|}
%    \end{aligned}
%\end{equation}

%\begin{equation}
%    \begin{aligned}
%        \frac{1}{2^{|\Bar{S}|}}\cdot
%        \frac{1}{2^{|\Bar{S'}|-|\Bar{S}|}}\cdot \frac{1}{2} \cdot [c(\emptyset,f|_{x_{S'},x_i=1},z)-c(\emptyset,f|_{x_{S'},x_i=0},z)]^2
%    \end{aligned}
%\end{equation}



%\begin{equation}
%\begin{aligned}
%    P_{\mathbf{z}\sim \mathcal{D}}(f(\mathbf{z})=f(\mathbf{x})|\mathbf{z}_S=\mathbf{x}_S)\geq\delta\\
%    \frac{|\{z\in\{0,1\}^{\overline{S}},f(x_S;z_{\Bar{S}})=f(x)|}{|\{z\in\{0,1\}^{\overline{S}}|}=\\
%\frac{|\{z\in\{0,1\}^{\overline{S}},f(x_S;z_{\Bar{S}})=f(x)|}{2^{|\Bar{S}|}} 
%\end{aligned}
%\end{equation}

%For the larger set: $S'=S\cup \{i\}$ it holds that. $|S'|=|S|+1$ and $|\Bar{S'}|=|\Bar{S}|-1$. Hence it holds that:

%\begin{equation}
%\begin{aligned}
%    P_{z\sim \mathcal{D}}(f(x)=f(z)|x_S'=z_S')=\\
%    \frac{|\{z\in\{0,1\}^{\overline{S'}},f(x_S;z_{\Bar{S'}})=f(x)|}{|\{z\in\{0,1\}^{\overline{S'}}|}=\\
%\frac{|\{z\in\{0,1\}^{\overline{S'}},f(x_{S'};z_{\Bar{S'}})=f(x)|}{2^{|\Bar{S'}|}} = \\
%\frac{|\{z\in\{0,1\}^{\overline{S'}},f(x_{S'};z_{\Bar{S'}})=f(x)|}{2^{|\Bar{S}|-1}}
%\end{aligned}
%\end{equation}


%\begin{equation}
%\begin{aligned}
%    P_{z\sim \mathcal{D}}(f(x)=f(z)|x_S'=z_S')=\\
%    \frac{|\{z\in\{0,1\}^{\overline{S'}},f(x_S;z_{\Bar{S'}})=f(x)|}{|\{z\in\{0,1\}^{\overline{S'}}|}=\\
%\frac{|\{z\in\{0,1\}^{\overline{S'}},f(x_{S'};z_{\Bar{S'}})=f(x)|}{2^{|\Bar{S'}|}} = \\
%\frac{|\{z\in\{0,1\}^{\overline{S'}},f(x_{S'};z_{\Bar{S'}})=f(x)|}{2^{|\Bar{S}|-1}}
%\end{aligned}
%\end{equation}

%For the global case:

%\begin{equation}
%\begin{aligned}
%    P_{x,z\sim \mathcal{D}}(f(x)=f(z)|x_S'=z_S')=\\
 %   \frac{|\{x\in\{0,1\}^n,z\in\{0,1\}^{\overline{S}},f(x_S;z_{\Bar{S}})=f(x)|}{|\{z\in\{0,1\}^{\overline{S'}},x\in\{0,1\}^n|}=\\
%\frac{|\{x\in\{0,1\}^n,z\in\{0,1\}^{\overline{S}},f(x_{S};z_{\Bar{S}})=f(x)|}{2^{|\Bar{S}|+n}} = \\
%\frac{|\{x\in\{0,1\}^n,z\in\{0,1\}^{\overline{S}},f(x_{S};z_{\Bar{S}})=f(x)|}{2^{|\Bar{S}|+n}}
%\end{aligned}
%\end{equation}

%\begin{equation}
%\begin{aligned}
%\forall{\mathbf{y},\mathbf{z}\in B_p^{\epsilon_p}(\mathbf{x})}, \quad
% [f(\mathbf{y}_{S};\mathbf{z}_{\Bar{S}})= f(\mathbf{y})], \\
%B_p^{\epsilon_p}(\mathbf{x}):=\{\mathbf{z}\in \mathbb{R}^{n} | \ \ ||\mathbf{x}-\mathbf{z}||_p\leq \epsilon_p\}
%\end{aligned}
%\end{equation}

%\begin{equation}
%\begin{aligned}
%\forall{\mathbf{z}\in B_p^{\epsilon_p}(\mathbf{x})}, \quad
% [f(\mathbf{x}_{S};\mathbf{z}_{\Bar{S}})= f(\mathbf{x})] \\
%B_p^{\epsilon_p}(\mathbf{x}):=\{\mathbf{z}\in \mathbb{R}^{n} | \ \ ||\mathbf{x}-\mathbf{z}||_p\leq \epsilon_p\}
%\end{aligned}
%\end{equation}

%\begin{equation}
%\begin{aligned}
%   |\x\in\mathcal{D},\z\in\mathcal{D}(\z_S=\x_S), \z_S=\x_S; f(\x)=f(\z)|
%\end{aligned}
%\end{equation}


%\begin{equation}
%\begin{aligned}
%\textbf{Pr}_{\x.\z\sim\mathcal{D}_p}(f(\x)=f(\z) | \x_S=\z_S)=\\
%\mathbb{E}_{\x\sim\mathcal{D}_p}[\textbf{Pr}_{\z\sim\mathcal{D}_p}(f(\x)=f(\z) | \x_S=\z_S)]
%\end{aligned}
%\end{equation}
\section{Introduction}
Despite substantial progress in methods for explaining ML model decisions, the literature has consistently unfortunately found that many desirable explanation types with different provable guarantees are computationally hard to obtain~\citep{BaMoPeSu20, van2022tractability}, with the difficulty typically worsening in complex or highly non-linear models~\citep{BaMoPeSu20, adolfi2024computational, adolfi2024complexity}.
%Despite significant advancements in methods for explaining ML model decisions,
%the literature has repeatedly unfortunately found that computing many highly-desired explanation types pose substantial computational challenges~\citep{BaMoPeSu20, van2022tractability}, which are further amplified in complex or highly non-linear models~\citep{BaMoPeSu20, adolfi2024computational, adolfi2024complexity}, where explanations are often most critical. 
As a result, the \emph{computational complexity} of obtaining explanations has become a central theoretical focus, with many recent works aiming to chart which types of explanations can be efficiently obtained for different kinds of models, and which remain out of reach~\citep{BaMoPeSu20, waldchen2021computational, arenas2022computing, arenas2023complexity, marzouk2024tractability, ordyniak2023parameterized, LABER2024106437, bhattacharjee2024auditing, blanc2021provably, blanc2022query}.



%As a consequence, the computational complexity of obtaining different provable explanations has emerged as a central theoretical research focus, with recent work systematically analyzing the tractability of different guarantees across diverse model classes~\citep{BaMoPeSu20, waldchen2021computational, arenas2023complexity, marzouk2024tractability, ordyniak2023parameterized, LABER2024106437}, aiming to clarify which forms of explanation are practically attainable or unattainable.


%A substantial body of work has introduced methods for explaining ML model decisions~\citep{dwivedi2023explainable, lundberg2017unified, ribeiro2016should, sundararajan2017axiomatic}. However, numerous studies have observed that such explanations often lack provable guarantees, raising concerns about their reliability in safety-critical and high-stakes applications~\citep{darwiche2020reasons, dasgupta2022framework}. This observation has motivated a growing line of research on explanation methods equipped with formal guarantees such as sufficiency, stability, or minimality~\citep{blanc2021provably, ignatiev2019abduction, arenas2022computing, blanc2022query, marques2020explaining, xue2023stability, adolfi2024computational}. At the same time, the literature has repeatedly emphasized that computing such explanations poses substantial computational challenges~\citep{BaMoPeSu20, arenas2022computing, arenas2023complexity, bhattacharjee2024auditing, van2022tractability}, which are exacerbated in highly non-linear models~\citep{BaMoPeSu20, adolfi2024computational, adolfi2024complexity}. As a result, the computational complexity of provable explanation methods has become a central research focus, with recent works systematically analyzing the tractability of different guarantees across diverse model classes~\citep{BaMoPeSu20, waldchen2021computational, arenas2023complexity, marzouk2024tractability, ordyniak2023parameterized, LABER2024106437}.


%Despite significant advancements in methods for explaining ML model decisions~\citep{dwivedi2023explainable, lundberg2017unified, ribeiro2016should, sundararajan2017axiomatic}, many approaches lack provable guarantees for the explanations they provide. To address this, there has been growing interest within the ML community in developing methods that produce explanations with different provable guarantees (e.g.,~\citep{blanc2021provably, ignatiev2019abduction, adolfi2024computational, arenas2022computing, darwiche2020reasons, dasgupta2022framework, xue2023stability, blanc2022query, marques2020explaining}, among many others). However, a key challenge in computing such explanations is their high computational complexity~\citep{BaMoPeSu20, arenas2023complexity, arenas2022computing, bhattacharjee2024auditing, van2022tractability}, which is significantly amplified in complex or highly non-linear models~\citep{BaMoPeSu20, adolfi2024computational,adolfi2024complexity}. This has consequently sparked considerable interest in examining the computational complexity of generating various types of explanations with different guarantees across diverse ML models~\citep{BaMoPeSu20, waldchen2021computational, adolfi2024computational, arenas2022computing, arenas2023complexity, marzouk2024tractability, ordyniak2023parameterized, LABER2024106437}, aiming to clarify which forms of explanation are practically attainable or unattainable.


\textbf{From sufficient to contrastive reasons.} Among studies on the computational complexity of generating explanations, two fundamental types of explanations for ML models were extensively examined: \begin{inparaenum}[(i)]
\item \emph{sufficient reasons} and \item \emph{contrastive reasons}\end{inparaenum}~\citep{BaMoPeSu20, arenas2022computing, audemard2022preferred, ArBaBaPeSu21, barcelo2025explaining, marques2022delivering, ignatiev2020contrastive, darwiche2020reasons}. %\begin{inparaenum}[(i)]  
%\item 
A \emph{sufficient reason} is a subset of input features $S$ such that when these features are fixed to specific values, the model's prediction remains unchanged, regardless of the values assigned to the complementary set $ \overline{S}$. %\item 
A \emph{contrastive reason} is a subset of input features $S$ such that modifying these features %(while keeping those in $\overline{S}$ fixed)
leads to a change in the model's prediction.  
%\end{inparaenum}

Unlike additive attribution methods, which allocate importance scores across features but are often hard for humans to interpret~\citep{kumar2020problems} or lack actionability~\citep{bilodeau2024impossibility}, sufficient and contrastive reasons provide \emph{discrete, condition-like explanations} that directly answer ``what is enough to justify this prediction?'' or ``what must change to flip it?''. Their intuitive nature has given them a central role in many classic XAI methods~\citep{ribeiro2018anchors, carter2019made, ignatiev2019abduction, dhurandhar2018explanations}, shown greater effectiveness in improving human prediction over additive models~\citep{ribeiro2018anchors, yin2022interpreting, dhurandhar2018explanations}, and proved useful for downstream tasks such as bias detection~\citep{balkir2022necessity, carter2021overinterpretation, la2021guaranteed, muthukumar2018understanding}, model debugging~\citep{jacovi2021contrastive}, and anomaly detection~\citep{davidsoncxad}.
%This intuitive inherent nature brought them to possess a central role in many classic XAI methods~\citep{ribeiro2018anchors, carter2019made, ignatiev2019abduction, dhurandhar2018explanations}, and have demonstrated ability to better improve human prediction compared to additive models~\citep{ribeiro2018anchors, yin2022interpreting, dhurandhar2018explanations}, and their usefulness for downstream tasks such as bias detection~\citep{balkir2022necessity, carter2021overinterpretation, la2021guaranteed, muthukumar2018understanding}, model debugging~\citep{jacovi2021contrastive}, and anomaly detection~\citep{davidsoncxad}. 
A well-established principle further holds that \emph{smaller} sufficient and contrastive explanations enhance interpretability, making \emph{minimality} a central guarantee of interest~\citep{ribeiro2018anchors, lopardo2023sea, carter2019made, BaMoPeSu20, arenas2022computing, blanc2021provably, waldchen2021computational}.

%Unlike additive attribution methods, which distribute importance scores across all features but can be often be hard to interpret for humans~\citep{kumar2020problems}, or lack actionability~\citep{bilodeau2024impossibility}, sufficient and contrastive reasons provide \emph{discrete, condition-like explanations} that directly answer ``what is enough to justify this prediction?'' or ``what must change to flip it?''. This makes them highly intuitive and closer to human reasoning, which explains their central role in many classic XAI methods~\citep{ribeiro2018anchors, carter2019made, ignatiev2019abduction, dhurandhar2018explanations}, the fact that they were shown to be able to better improve human prediction capabilities compared to additive models~\citep{ribeiro2018anchors, yin2022interpreting,dhurandhar2018explanations}, and serve downstream tasks such as bias detection~\citep{balkir2022necessity, carter2021overinterpretation, la2021guaranteed, muthukumar2018understanding}, model debugging~\citep{jacovi2021contrastive}, and anomaly detection~\citep{davidsoncxad}. A well-established principle further suggests that \emph{smaller} explanations — both sufficient and contrastive — enhance interpretability, making \emph{minimality} an additional central guarantee of interest~\citep{ribeiro2018anchors, marques2022delivering, carter2019made, BaMoPeSu20, arenas2022computing, blanc2021provably, waldchen2021computational}.

%These types of explanations align with many widely used methods in the XAI community~\citep{ribeiro2018anchors, carter2019made, ignatiev2019abduction, dhurandhar2018explanations}. A well-established principle in the literature suggests that \emph{smaller} explanations --- both sufficient and contrastive --- enhance interpretability, making the \emph{minimality} of these subsets an additional central guarantee of interest~\citep{ribeiro2018anchors, marques2022delivering, carter2019made, BaMoPeSu20, arenas2022computing, blanc2021provably, waldchen2021computational}.% A well-established principle in the literature suggests that smaller explanations --- both sufficient and contrastive --- enhance interpretability~\citep{ribeiro2018anchors, marques2022delivering, carter2019made, BaMoPeSu20}. Consequently, ensuring the \emph{minimality} of these subsets is an additional key guarantee of interest~\citep{BaMoPeSu20, arenas2022computing, blanc2021provably, waldchen2021computational}.

%Within works that study the computational complexity of obtaining explanations, perhaps two of the most commonly studied forms of explanations consider two fundamental types of explanations for an ML model: a sufficient reason, and a contrastive reason: \begin{inparaenum}[(i)] \item A \emph{sufficient reason}, is a subset of input features $S$, that when fixed to some given set of values, determnies that a prediction of a model remains unchanged, regardless of the assignments to the complementray set $\overline{S}$. \item A \emph{contrastive reason}, is a subset of input features $S$, that when altered, (such that the features in $\overline{S}$ remain fixed) causes the prediction of a model to change. These forms of explanations are aligned with many methods that are used in practice within the XAI community~\citep{}. A well-established notion in the literature is that smaller explanations (both sufficient and contrastive) provide a better interpretation~\citep{}, and hence the \emph{minimality} of these subsets is also a highly studied formal guarantee of interest~\citep{}.

\cite{BaMoPeSu20} conducted one of the earliest studies on the complexity of deriving sufficient and contrastive reasons. Their work established that finding the minimal-sized sufficient reason for a decision tree is NP-hard, with the complexity further increasing to $\Sigma^P_2$-hard for neural networks. In the case of contrastive reasons, they demonstrated that computing the smallest possible explanation is solvable in polynomial time for decision trees but becomes NP-Hard for neural networks. Similar hardness results were later shown to hold for tree ensembles as well~\citep{izyamajo21, ordyniak2024explaining, audemard2022trading}.

\textbf{From non-probabilistic to probabilistic explanations.} 
A common criticism of the classic definition of sufficient and contrastive reasons is their rigidity and lack of flexibility, as they hold in an absolute sense over entire domains and can thus lead to excessively large or uninformative explanations~\citep{ignatiev2019abduction, ignatiev2020towards, arenas2022computing, waldchen2021computational}. %A common criticism of the definition of sufficient and contrastive reasons is their rigidity and lack of flexibility~\citep{arenas2022computing, waldchen2021computational, ignatiev2020towards}. Specifically, these explanations hold in an absolute sense over entire domains, which can lead to excessively large or uninformative explanations~\citep{ignatiev2019abduction, arenas2022computing, waldchen2021computational, izza2023computing}. 
To address these limitations, the literature has shifted towards a more general definition that incorporates a \emph{probabilistic} perspective on sufficient and contrastive reasons~\citep{waldchen2021computational, arenas2022computing, blanc2021provably, izza2023computing, xue2023stability}. Under this framework, the goal is to identify subsets of input features that influence a prediction with a probability exceeding a given threshold $\delta$.

%Following works on the computation of sufficient and contrastive explanations~\citep{ignatiev2019abduction, wu2024verix, bassan2023towards, darwiche2020reasons, huang2021efficiently}, A widespread critisism that was brought over the definitions of these explanations they lack flexability, and are too strict as they hold in an absolute sens over complete domains. This, for example may result in extensively large or unmeaningful explanations~\citep{ignatiev2019abduction, bassan2023towards}. To tackle these limitations, there has been a shift towards studying a generalization of this definition that considers of a \emph{probabilistic} definition of sufficient and contrastive reasons~\citep{waldchen2021computational, arenas2022computing, blanc2021provably, izza2023computing}. Under this new definition, one seeks to obtain subsets of input features that determine or alter a prediction with some probability above $\delta$. 

\cite{waldchen2021computational} were the first to study the complexity of probabilistic sufficient reasons, showing that for CNF classifiers the problem is $\text{NP}^{\textbf{PP}}$-hard. This hardness extends to tree ensembles and neural networks~\citep{BaMoPeSu20, ordyniak2023parameterized}. A central theoretical insight in the probabilistic setting is the \emph{lack of monotonicity} in the probability function~\citep{arenas2022computing, izza2023computing}, which makes even \emph{subset minimal} explanations computationally hard. Strikingly,~\citep{arenas2022computing} show that finding a subset minimal probabilistic sufficient reason is NP-hard even for decision trees --- unlike in the non-probabilistic case, where such explanations are computable in polynomial time~\citep{huang2021efficiently, BaMoPeSu20}.

%\cite{waldchen2021computational} were the first to examine the complexity of obtaining probabilistic sufficient reasons, demonstrating that for CNF classifiers, this problem is $\text{NP}^{\textbf{PP}}$-hard. Consequently, this hardness result extends to tree ensembles and neural networks~\citep{BaMoPeSu20, ordyniak2023parameterized}. A particularly important and intriguing theoretical insight that was found in the probabilistic setting is the \emph{lack of monotonicity} in the probability function~\citep{arenas2022computing, izza2023computing}. This observation is crucial in establishing that even a \emph{subset minimal} (or locally minimal) explanation can be computationally challenging to obtain. Notably,~\citep{arenas2022computing} present the surprising result that finding a subset minimal probabilistic sufficient reason is NP-hard even for decision trees. This contrasts with the non-probabilistic setting, where such an explanation can be computed in polynomial time~\citep{huang2021efficiently, BaMoPeSu20}.



\textbf{From local to global explanations.} In a more recent study,~\cite{bassanlocal} extend the complexity analysis from the \emph{local} (non-probabilistic) setting --- where explanations are tied to individual predictions --- to the \emph{global} (non-probabilistic) setting, which seeks sufficient or contrastive reasons over entire domains. However, as with other non-probabilistic methods~\citep{ignatiev2019abduction, BaMoPeSu20, ArBaBaPeSu21, darwiche2020reasons}, the criteria are extremely strict --- arguably even more so in the global setting than in the local one. For instance, any feature excluded from a global sufficient reason is deemed strictly redundant~\citep{bassanlocal}, often making the explanation span nearly all input features and thus less informative.


%In a more recent study,~\cite{bassanlocal} shift the focus of complexity analysis from the \emph{local} (non-probabilistic) setting --- where explanations are computed for individual predictions --- to the \emph{global} (non-probabilistic) setting, which aims to derive sufficient or contrastive reasons over entire global domains. However, similar to other non-probabilistic explanation methods~\citep{ignatiev2019abduction, BaMoPeSu20, ArBaBaPeSu21, darwiche2020reasons}, this approach is extremely stringent --- arguably even more so in the global setting than in the local one. For instance, any feature not included in a global non-probabilistic sufficient reason is deemed strictly redundant~\citep{bassanlocal}. In practice, this could lead to the global non-probabilistic sufficient reason encompassing all, or nearly all, of the input features, thereby diminishing its usefulness as an explanation.



\subsection*{Our Contributions}

\begin{enumerate}
    \item We unify previous explanation computation problems --- including sufficient, contrastive, probabilistic, non-probabilistic, as well as local and global --- into one framework, described as a minimization task over a unified value function. We then identify three fundamental properties of the value function that significantly impact the complexity of this task: \begin{inparaenum} \item \emph{monotonicity}, \item \emph{supermodularity}. and \item \emph{submodularity.} \end{inparaenum}
    \item Interestingly, we show that these properties behave in \emph{strikingly different manners} depending on the structure of the value function. In particular, we demonstrate the surprising result that while the \emph{local} value functions for both sufficient and contrastive reasons are \emph{non-monotonic}, their \emph{global} counterparts are \emph{monotonic non-decreasing}. Moreover, we identify additional intriguing properties unique to the global setting: the global sufficient value function is \emph{supermodular}, whereas the global contrastive value function is \emph{submodular} --- in contrast to the local setting, where neither property holds.
    \item We leverage these properties to derive new complexity results for explanation computation, revealing the intriguing finding that global explanations with guarantees can be computed efficiently, even though computing their local counterparts remains computationally hard. We demonstrate these findings across three widely used model types that span the interpretability spectrum: \begin{inparaenum}[(i)] \item
     neural networks, \item decision trees, and \item tree ensembles. \end{inparaenum} First, we prove that while computing a subset-minimal local sufficient/contrastive probabilistic explanation is NP-hard even for decision trees~\citep{arenas2022computing}, its global counterpart can be computed in polynomial time. We further extend this result to any black-box model (including complex models such as neural networks and tree ensembles) when using empirical distributions. Specifically, we show that obtaining a subset-minimal global sufficient/contrastive explanation is achievable in polynomial time, whereas the local version remains NP-hard for these models.
    \item Finally, we present an even stronger complexity result by exploiting the submodular and supermodular properties of the value functions in the global setting --- properties that do not hold in the local case. Specifically, we show that it is possible to achieve provable \emph{constant-factor} approximation guarantees for computing cardinally minimal global explanations --- even for complex models like neural networks or tree ensembles --- when the empirical distribution is fixed. In sharp contrast, we establish strong inapproximability results for the local setting, demonstrating that no \emph{bounded} approximation is possible, even under very simple assumptions.


\end{enumerate}

Owing to space constraints, we provide an overview of our main theorems and corollaries in the main text, with full proofs deferred to the appendix.


%we present only a brief outline of the proofs for our claims in the paper, with the complete proofs provided in the appendix.


%\textbf{Our contributions.}

%A primary limitation of sufficient and contrastive reasons, is that they hold for 
    
%\end{inparaenum}

\section{Preliminaries}

\textbf{Setting.} We consider an input space of dimension $n\in\mathbb{N}$, with input vectors $\x:=(\x_1,\x_2,\ldots,\x_n)$. Each coordinate $\x_i$ may take values from its corresponding domain $\mathcal{X}_i$, which can be either discrete or continuous. The full input feature space is therefore $\mathbb{F}:=\mathcal{X}_1\times \mathcal{X}_2 \times \ldots \times \mathcal{X}_n$. We consider classification models $f:\mathbb{F}\to [c]$ where $c\in\mathbb{N}$ is the number of classes. Moreover, we consider a \emph{generic} distribution $\mathcal{D} : \mathbb{F} \to [0,1]$ over the input space. In many settings, however, accurately approximating $\mathcal{D}$ is computationally infeasible. A natural alternative is to instead work with a fixed empirical dataset $\mathbf{D} := \{\mathbf{z}^1, \mathbf{z}^2, \ldots, \mathbf{z}^{|\mathbf{D}|}\}$, which serves as a practical proxy for the ``true'' distribution, for instance, by taking $\mathbf{D}$ as a sampled subset from the available training data.

%We define a \emph{generic} distribution $\mathcal{D}:\mathbb{F}\to [0,1]$ over the input features. In some cases, approximating $\mathcal{D}$ is computationally infeasible, so a natural assumption would be to instead lookg at some fixed dataset $\mathbf{D}$ comprising of vectors $\z^1,\z^2,\ldots,\z^{|\mathbf{D}|}$ that represents the ``true'' distribution, for example by being taken from the training data.


The explanations that we study are either local or global. In the local case, the explanations target a specific prediction $\x\in\mathbb{F}$, providing a form of reasoning for why the model predicted $f(\x)$ for that instance. In the global case, explanations aim to reflect the general reasoning behind the behavior of $f$ across a wider region of the input space, independent of any specific $\x$, and to characterize its overall decision-making logic.


\textbf{Models.} While many results presented in this work are general (e.g., inherent properties of value functions), we also provide some model-specific complexity results for widely-used ML models. To broadly address the interpretability spectrum, we chose to analyze models ranging from those typically considered ``black-box'' to those commonly regarded as ``interpretable''. Specifically, we focus on: \begin{inparaenum}[(i)] \item decision trees, \item neural networks (all architectures at least as expressive as feed-forward ReLU networks), and \item tree ensembles, including majority-voting ensembles (e.g., random forests) and weighted-voting ensembles (e.g., XGBoost).
\end{inparaenum} Formal definitions are provided in the Appendix. 

\textbf{Distributions.} We emphasize that in probabilistic explanation settings, the complexity can vary significantly with the input distribution $\mathcal{D}$. Particularly, we focus on three distribution types:
\begin{inparaenum}[(i)] \item \emph{general distributions}, which make no specific assumptions over $\mathcal{D}$ and thus encompass all possible distributions. We use these distributions mainly in proofs of properties that hold universally; \item \emph{empirical distributions}, which include all distributions derived from the finite dataset $\mathbf{D}$ --- an approach commonly employed in XAI~\citep{lundberg2017unified, van2022tractability}; and
\item \emph{independent distributions}, which assume that features in $\mathcal{D}$ are mutually independent --- another widely adopted assumption in XAI literature~\citep{arenas2022computing, arenas2023complexity, lundberg2017unified, ribeiro2018anchors}.
\end{inparaenum} We note that empirical distributions do not necessarily imply feature independence; rather, they can represent complex dependencies extracted from finite datasets~\citep{van2022tractability}. The complete formal definitions of these distributions are provided in the Appendix.

%While many of the results presented in this work are general (e.g., inherent properties of the value functions), we will present some model-specific computational complexity results over popular ML models. We aim for this anlysis to capture a broad range of the interpretability spectrum, and hence we will analyze both models that are considered ``black-box'' models, and those that are regarded as interpretable. Parritularly, our analysis will focus on: \begin{inparaenum}[(i)]
%    \item decision trees, \item neural networks with ReLU activations, and \item tree ensembles, where our results apply to both majority-voting based ensembles such as random forests, as well as to weighted-voting based ensembles, which includes boosted trees (e.g., XGBoost).
%\end{inparaenum} The full formalization of all of these models is attached within Appendix~\ref{}.




\section{Forms of Reasons}

In this section, we introduce the explanation types studied in this work, starting with the strict non-probabilistic definitions of (local/global) sufficient and contrastive reasons, and then extending to their more flexible and generalizable probabilistic counterparts.


%In this section, we introduce the types of explanations examined in this work, beginning with the strict non-probabilistic definitions of (local/global) sufficient and contrastive reasons. We then extend our discussion to the more flexible and generalized probabilistic definitions of (local/global) sufficient and contrastive reasons.

\subsection{\emph{Non-Probabilistic} Sufficient and Contrastive Reasons}

\textbf{Sufficient reasons.} In the context of \emph{feature selection}, users often select the top $k$ features that contribute to a model's decision. We examine the well-established \emph{sufficiency} criterion for this selection, which aligns with commonly used explainability methods~\citep{ribeiro2018anchors, carter2019made, ignatiev2019abduction, dasgupta2022framework}. This feature selection can be carried out either locally --- focusing on a single prediction --- or globally --- across the entire input domain. Following standard conventions, we define a \emph{local sufficient reason} as a subset of features $S\subseteq \{1,\ldots,n\}$ such that when the features in $S$ are fixed to their corresponding values in $\x$, the model's prediction remains $f(\x)$ regardless of the values assigned to the remaining features $\overline{S}$. Formally, $S$ is a local sufficient reason for $\langle f,\x\rangle$ iff the following condition holds:

\begin{equation}
    \forall \z\in\mathbb{F}, \quad f(\x_S;\z_{\Bar{S}})=f(\x).
\end{equation}

Here, $(\x_S;\z_{\Bar{S}})$ denotes a vector where features in $S$ take their values from $\x$, and those in $\overline{S}$ from $\z$. Local sufficient reasons are closely related to \emph{semi-factual} explanations~\citep{kenny2021generating, alfano2025even}, which search for alternative inputs $\mathbf{z}'$ that keep the prediction unchanged (i.e., $f(\mathbf{z}') = f(\mathbf{x})$) while being as ``close'' as possible to the original point. When we instantiate $\mathbf{z}'$ as $(\x_S; \z_{\bar{S}})$ and measure proximity via Hamming distance, the two notions coincide.

A \emph{global sufficient reason}~\citep{bassanlocal} is a subset of input features $S \subseteq \{1, \ldots, n\}$ that serves as a local sufficient reason for \emph{every} possible input $\x \in \mathbb{F}$:

\begin{equation}
    \forall \x,\z\in\mathbb{F}, \quad f(\x_S;\z_{\Bar{S}})=f(\x).
\end{equation}

\textbf{Contrastive reasons.} Another prevalent approach to providing explanations is by pinpointing subsets of input features that \emph{modify} a prediction~\citep{dhurandhar2018explanations, mothilal2020explaining, guidotti2024counterfactual}. This type of explanation aims to determine the minimal changes necessary to alter a prediction. Formally, a subset $S\subseteq \{1,\ldots, n\}$ is defined as a \emph{local contrastive reason} concerning $\langle f,\x\rangle$ iff:

\begin{equation}
     \exists\z\in\mathbb{F}, \quad f(\z_S;\x_{\Bar{S}})\neq f(\x).
\end{equation}

Contrastive reasons are also closely connected to \emph{counterfactual} explanations~\citep{guidotti2024counterfactual, mothilal2020explaining}, which seek a nearby assignment $\x'$ for which the prediction flips, i.e., $f(\x') \neq f(\x)$. When the distance metric is taken to be the Hamming weight, the two notions coincide.

Similarly to sufficient reasons, one can determine a subset of input features that, when altered, changes the prediction for all inputs within the domain of interest. This form of explanation is also closely connected to approaches for identifying bias~\citep{ArBaBaPeSu21, bassanlocal, darwiche2020reasons}, as well as to \emph{group} counterfactual explanation methods that search for counterfactuals over multiple data instances~\citep{carrizosa2024mathematical, warren2024explaining}. Formally, we define a subset $S$ as a \emph{global contrastive reason} with respect to $f$ iff:

\begin{equation}
    \forall \x\in\mathbb{F}, \quad \exists \z\in\mathbb{F}, \quad f(\z_S;\x_{\Bar{S}})\neq f(\x).
\end{equation}

\subsection{\emph{Probabilistic} Sufficient and Contrastive Reasons}

As discussed in the introduction, non-probabilistic sufficient and contrastive reasons are often criticized for imposing significantly overly strict conditions, motivating a shift to \emph{probabilistic} definitions~\citep{arenas2022computing, ribeiro2018anchors, izza2023computing, waldchen2021computational, bounia2023approximating, subercaseaux2024probabilistic, wang2021probabilistic, blanc2021provably}, which generalize these requirements by demanding that the guarantees hold with probability at least $\delta \in [0,1]$. The special case $\delta=1$ recovers the original non-probabilistic definitions.

%As mentioned within the introduction, a common criticism of non-probabilistic sufficient and contrastive reasons is that they impose overly strict conditions on the underlying subset, thereby motivating a shift to \emph{probabilistic} definitions~\citep{arenas2022computing, ribeiro2018anchors, izza2023computing, waldchen2021computational, bounia2023approximating, subercaseaux2025probabilistic, wang2021probabilistic, blanc2021provably}, which offer greater flexibility by requiring the conditions to hold with probability at least $\delta \in [0,1]$. A special case of this formulation is taking $\delta = 1$ exactly, which reverts to the original, non-probabilistic definitions.


%A common criticism of the previous definitions of sufficient and contrastive reasons is that they impose overly strict conditions on the underlying subset~\citep{arenas2022computing, izza2023computing, waldchen2021computational, bounia2023approximating, subercaseaux2025probabilistic}. In particular, non-probabilistic sufficient reasons require that the classification remains unchanged for \emph{every} value in the complementary set, making this criterion harder to satisfy and potentially leading to excessively large and less meaningful subsets~\citep{arenas2022computing, waldchen2021computational, ignatiev2019abduction, marques2022delivering}. In the global setting, this issue is further amplified, as we evaluate a much stronger sufficiency condition across the entire domain. A similar and inverse argument can be applied additionally to the definitions of contrastive reasons. This motivates redefining these concepts in a \emph{probabilistic} manner~\citep{arenas2022computing, ribeiro2018anchors, izza2023computing, waldchen2021computational, bounia2023approximating, subercaseaux2025probabilistic, wang2021probabilistic, blanc2021provably}, offering significantly greater flexibility. Specifically, we require that all previously established conditions hold with a probability above a certain threshold $\delta \in [0,1]$. A special case of this formulation is taking $\delta = 1$ exactly, which reverts to the original, non-probabilistic definitions.

%A common criticisim of the previous definition of sufficient and contrastive reasons is being too strict in the conditions required on the underlying subset. Particularly, non-probabilistic sufficient reasons which require that the classification remains the same for \emph{any} value in the complementary set, make this criterion less likely, and may provide substantially large and less meaningful subsets~\citep{}. For the global setting, this procedure is substantially enhanced since we are evaluating a much stronger sufficiency condition over the entire domain. For contrastive reasons, a similar reasoing can be made - as we would want to theoretically identify subsets that alter a prediction with some probability exceeding a threshold, and not only validating one particular instance. This brings forth the option of defining the previous definitions in a \emph{probabilistic} fashion~\citep{}, which is signficantly more flexible. We partifularly will want that all of the previously established conditions, will hold for some probability above a certain threshod $\delta\in[0,1]$. A private case of all of these definitions will be the case where $\delta=1$ which will yield the previous non-probabilistic definitions. We will formalize all probabilistic definitions in the following section.

\textbf{Probabilistic sufficient reasons.} We define $S\subseteq \{1,\ldots,n\}$ as a \emph{local $\delta$-sufficient reason} with respect to  $\langle f,\x\rangle$ if, when the features in $S$ are fixed to their corresponding values in $\x$ and the remaining features are sampled from a distribution $\mathcal{D}$, the classification remains unchanged with probability at least $\delta$. In other words:



%We say that $S\subseteq \{1,\ldots, n\}$ is a \emph{local probabilistic $\delta$-sufficient reason} concerning $\langle f,\x\rangle$ if when given that we fix the features of $S$ to the there corresponding values within $\x$, given that we sample the complementry features from some distribution $\mathcal{D}_p$, the classification remains the same with probability greater or equal to $\delta$, or in other words that:

\begin{equation}
\label{eq_suff_local}
    \textbf{Pr}_{\z\sim \mathcal{D}}(f(\z)=f(\x) \ | \  \z_S=\x_S)\geq \delta.
\end{equation}

where $\z_S = \x_S$ denotes that the features in $S$ of vector $\z$ are fixed to their corresponding values in $\x$. We adopt the standard notion of global explanations --- by averaging over all inputs in the global domain --- and define a subset $S \subseteq \{1, \ldots, n\}$ as a \emph{global $\delta$-sufficient reason} with respect to $\langle f \rangle$ if, when taking the expectation of the local sufficiency probability over samples drawn from the distribution $\mathcal{D}$, the expectation remains with value at least $\delta$. In other words:


%\overset{\text{FI}}{=} 
%    \textbf{Pr}_{\x,\z\sim \mathcal{D}}(f(\x)=f(\z) \ | \  \x_S=\z_S)\geq \delta

\begin{equation}
\label{eq_suff_global}
\mathbb{E}_{\x\sim\mathcal{D}}[\textbf{Pr}_{\z\sim \mathcal{D}}(f(\z)=f(\x) \ | \  \z_S=\x_S)] \geq \delta.
\end{equation}

\textbf{Probabilistic contrastive reasons.} Similar to the non-probabilistic case, we define a \emph{local $\delta$-contrastive reason} as a subset of input features that changes a prediction with some probability. Here, unlike sufficient reasons, we set the features of the \emph{complementary} set $\overline{S}$ to their respective values in $\x$, and when allowing the features in $S$ to vary, we require the prediction to differ from the original prediction $f(\x)$ with a probability exceeding $\delta$. Formally:

\begin{equation}
\label{eq_con_local}
    \textbf{Pr}_{\z\sim \mathcal{D}}(f(\z) \neq f(\x) \ | \  \z_{\Bar{S}}=\x_{\Bar{S}})\geq \delta.
\end{equation}

For the global setting, we define a subset $S$ to be a \emph{global $\delta$-contrastive reason}, analogous to global sufficient reasons, by computing the expectation over all local contrastive reasons sampled from the distribution $\mathcal{D}$, and requiring that this expected value exceeds $\delta$:


%For the global setting, we define a subset $S$ as a global $\delta$-contrastive reason if when fixing the features in $\Bar{S}$, the classification is altered with probability exceeding $\delta$:


%\overset{\text{FI}}{=} 
%    \textbf{Pr}_{\x,\z\sim \mathcal{D}}(f(\x) \neq f(\z) \ | \  \x_{\Bar{S}}

\begin{equation}
\label{eq_con_global}
\mathbb{E}_{\x\sim\mathcal{D}}[\textbf{Pr}_{\z\sim \mathcal{D}}(f(\z) \neq f(\x) \ | \  \z_{\Bar{S}}= \x_{\Bar{S}})]\geq \delta.
\end{equation}

%\textbf{Note.} We observe that in the specific cases where $\delta=1$, all probabilistic and non-probabilistic notions converge. This applies across all settings of local/global and sufficient/contrastive reasons. However, allowing probabilistic computations to be performed with respect to a given $\delta$ provides significantly greater flexibility in selecting the underlying sufficiency conditions.

\section{Forms of Minimality}

As discussed in the introduction, across all the explanation forms discussed so far --- whether sufficient or contrastive, local or global --- a common assumption in the literature is that explanations of \emph{smaller size} are more meaningful, thereby making their \emph{minimality} a particularly important provable guarantee. In this study, we explore two central notions of minimality across all our explanation types:%, including sufficient and contrastive, local and global, as well as probabilistic and non-probabilistic explanations: 

\begin{definition}
Assuming a subset $S\subseteq \{1,\ldots n\}$ is an explanation, then:
    \begin{enumerate}
        \item $S$ is a \emph{cardinally-minimal} explanation~\citep{BaMoPeSu20, bassanlocal} iff $S$ has the smallest explanation cardinality $|S|$ (i.e., there is no explanation $S'$ such that $|S'|<|S|$).
        \item $S\subseteq \{1,\ldots, n\}$ is a \emph{subset-minimal} explanation~\citep{arenas2022computing, ignatiev2019abduction} iff $S$ is an explanation, and any $S'\subseteq S$ is not an explanation.
    \end{enumerate}
\end{definition}


%\begin{inparaenum}[(i)] \item The first is a \emph{cardinally minimal} explanation~\citep{BaMoPeSu20, bassanlocal} (also referred to as globally minimal), meaning that the subset $S$ has the smallest possible cardinality $|S|$. \item The second is a \emph{subset minimal} explanation~\citep{arenas2022computing, ignatiev2019abduction} (sometimes called locally minimal), which is a subset $S$ that serves as an explanation, and any $S'\subseteq S$ is not an explanation.

We note that cardinal-minimality is strictly stronger than subset-minimality: every cardinally minimal $S$ is subset-minimal, but not vice versa (see Appendix~\ref{subset_cardinal_minimality} for an example). We use the terms subset and cardinally minimal, rather than local and global minima, to avoid confusion with local vs. global explanations (input- vs. domain-level reasoning). Both notions apply to \emph{all} explanation forms we analyze. For instance, a cardinally minimal local probabilistic $\delta$ sufficient reason is the one with the smallest $|S|$, while a subset-minimal one is any $S$ where no proper subset $S'\subseteq S$ also qualifies. 

%We note that cardinal-minimality is substantially stronger than subset-minimality: every cardinally minimal $S$ is subset-minimal, but not vice versa. An example illustrating this is given in Appendix~\ref{}. We also note that both definitions apply to \emph{all} explanation forms we analyze. For example, a cardinally minimal local probabilistic $\delta$ sufficient reason is the one with the smallest $|S|$, while a subset minimal one is any such $S$ where no proper subset $S' \subseteq S$ also qualifies. We use the terms subset minimal and cardinally minimal, rather than local and global minima, to avoid confusion with local vs. global explanations (input- vs. domain-level reasoning).

%We also note that both definitions apply to \emph{all} explanation forms we analyze. For instance, a cardinally minimal local probabilistic $\delta$ sufficient reason refers to the local probabilistic $\delta$ sufficient reason with the smallest cardinality $|S|$, while a subset minimal one is a subset $S$, which is a local probabilistic $\delta$ sufficient reason, where no proper subset $S' \subseteq S$ qualifies as a local probabilistic $\delta$ sufficient reason. We use the terms subset minimal and cardinally minimal instead of local and global minima to avoid confusion with local vs. global explanations, which refer to input-level versus domain-level reasoning.


%In this work, we discuss two main forms of minimality: the first is a \emph{cardinally minimal} sufficient reason (sometimes refered to as a globally-minimal), which means that the subset $S$ has minimal cardinality $|S|$. The second is a \emph{subset minimal} (sometimes refered to as a locally-minimal) which is a subset $S$ which is a local/global sufficient/contrastive reason such that for any $S'\subseteq S$ it holds that $S'$ is not a local/global sufficient/contrastive reason (where the same definition holds for either the probabilistic or non-probabilistic scenario). We use the terms subset minimal and cardinally minimal (and not local and global) to not confuse this with local and global explanations (explanations over individual inputs, or over entire global domains).

\section{A Unified Combinatorial Optimization Task}
%\guy{It takes slightly too long  to reach the new contribution - 4.5 pages in, in a 9-page paper. There is a bit of redundnacy in the text in previous sections - if we need to shorten the paper, we could remove some of that, and then section 5 would arrive more quickly. However, this is not crucial, I think.}

Interestingly, all previously discussed computational problems --- local or global, sufficient or contrastive, probabilistic or not --- can be cast as finding a minimal-size subset $S$ such that $v(S) \geq \delta$, where in non-probabilistic settings we set $\delta := 1$. We now introduce notation for the value functions: let $v^{\ell}_{\text{suff}}$ denote the local sufficiency probability from equation~\ref{eq_suff_local}, i.e., $\textbf{Pr}_{\z\sim\mathcal{D}p}(f(\x)=f(\z) \mid \x_S=\z_S)$, and define the global variant as $v^{g}_{\text{suff}}$ (equation~\ref{eq_suff_global}). Similarly, let $v^{\ell}_{\text{con}}$ and $v^{g}_{\text{con}}$ denote the local and global contrastive probabilities (equations~\ref{eq_con_local} and~\ref{eq_con_global}, respectively). Using this notation, we now formally define the unified task of finding a cardinally minimal $\delta$-local/global sufficient/contrastive reason:



%All the previously mentioned computational problems --- whether local or global, sufficient or contrastive, probabilistic or non-probabilistic --- can be framed as the task of finding a subset $S$ of minimal possible size such that $v(S) \geq \delta$. In the non-probabilistic cases, this corresponds specifically to setting $ \delta := 1$. We begin by introducing the following notation for the value function. Let $ v^{\ell}_{\text{suff}} $ represent the local sufficiency probability function from equation~\ref{eq_suff_local}, i.e., $ \textbf{Pr}_{\z\sim\mathcal{D}_p}(f(\x)=f(\z) \mid \x_S=\z_S) $. Similarly, we define the global sufficiency probability function from equation~\ref{eq_suff_global} as $ v^{g}_{\text{suff}}$. Additionally, we denote the local contrastive probability function from equation~\ref{eq_con_local} as $ v^{\ell}_{\text{con}} $ and its global counterpart from equation~\ref{eq_con_global} as $ v^{g}_{\text{con}} $. With these definitions, we can now formulate the unified problem of finding a cardinally minimal $\delta$-local/global sufficient/contrastive reason as follows:

\noindent\fbox{%
	\parbox{\columnwidth}{%
		\mysubsection{Cardinally Minimal $\delta$-Explanation}:
		
		\textbf{Input}: Model $f$, a distribution $\mathcal{D}$, (possibly, an input $\x$), a \emph{general} value function $v:2^n\to [0,1]$ (defined using $f$, $\mathcal{D}$, and possibly $\x$), and some $\delta\in [0,1]$.
		
		\textbf{Output}: 
		A subset $S\subseteq [n]$ such that $v(S)\geq \delta$ and $|S|$ is minimal.
	}%
}

Similarly, for the relaxed condition where the goal is to obtain a subset-minimal rather than a cardinally-minimal local or global sufficient/contrastive reason, we define the following relaxed optimization objective:


\noindent\fbox{%
	\parbox{\columnwidth}{%
		\mysubsection{Subset Minimal $\delta$-Explanation}:
		
		\textbf{Input}: Model $f$, a distribution $\mathcal{D}$, (possibly, an input $\x$), a \emph{general} value function $v:2^n\to [0,1]$ (defined using $f$, $\mathcal{D}$, and possibly $\x$), and some $\delta\in [0,1]$.
		
		\textbf{Output}: 
		A subset $S\subseteq [n]$ such that $v(S)\geq \delta$ and for any $S'\subseteq S$ it holds that $v(S')<\delta$.
	}%
}


\textbf{Properties of $v$ that affect the complexity.} We now outline several key properties of the value function, which we later show play a crucial role in determining the complexity of generating explanations. The first property is \emph{non-decreasing monotonicity}, which ensures that the marginal contribution $v(S \cup \{i\}) - v(S)$ is consistently non-negative. Formally:


\begin{definition}
    We say that a value function $v$ maintains \emph{non-decreasing monotonicity} if for any $S\subseteq\{1,\ldots,n\}$ and any $i\in\{1,\ldots,n\}$ it holds that: $v(S\cup\{i\})\geq v(S)$. 
\end{definition}

%\begin{definition}


%The first property is \emph{saturation}, which intuitively means that once the value function $v(S)$ reaches its optimal value, adding more features will not further increase it. Formally:

%\begin{definition}
%    We say that a value function $v$ maintains \emph{saturation} if given that $v_{max}:=\max_{S\subseteq \{1,\ldots,n\}}v(S)$ and we have some $S'$ for which $v(S')=v_{max}$, then for all $S''$ for which $S'\subseteq S''$,  it holds that: $v(S'')=v_{max}$.
%\end{definition}

%The second key property, \emph{monotonicity}, ensures that the marginal contribution $v(S \cup \{i\}) - v(S)$ is consistently non-negative for non-decreasing monotone functions or consistently non-positive for non-increasing monotone functions. In our work, we primarily focus on the non-increasing monotone case, which aligns with certain configurations of the value function. Formally:
%\begin{definition}
%    We say that a value function $v$ maintains \emph{non-decreasing (resp. non-increasing) monotonicity} if for any $S\subseteq\{1,\ldots,n\}$ and any $i\in\{1,\ldots,n\}$ it holds that: $v(S\cup\{i\})\geq v(S)$ (resp. $v(S\cup\{i\})\leq v(S)$). 
%\end{definition}

The other key properties of \emph{supermodularity} and \emph{submodularity} pertain to the behavior of the marginal contribution $v(S \cup \{i\}) - v(S)$. Specifically, in the supermodular case, this contribution forms a monotone non-decreasing function. In contrast, under the dual definition of the \emph{submodular} case, the marginal contribution $v(S \cup \{i\}) - v(S)$ is a monotone non-increasing function. Formally:


\begin{definition}
    Let there be some value function $v$, some $S\subseteq S'\subseteq\{1,\ldots,n\}$, and $i\not\in S'$. Then: 
    \begin{enumerate}
        \item $v$ maintains \emph{supermodularity} iff it holds that: $v(S\cup\{i\})-v(S)\leq v(S'\cup\{i\})-v(S')$.
        \item $v$ maintains \emph{submodularity} iff it holds that: $v(S\cup\{i\})-v(S)\geq v(S'\cup\{i\})-v(S')$.
    \end{enumerate}
\end{definition}

%\begin{definition}
%    We say that a value function $v$ maintains \emph{supermodularity} if for any $S, S'\subseteq\{1,\ldots,n\}$ such that $S\subseteq S'$ and any $i\not\in S'$ it holds that: $v(S\cup\{i\})-v(S)\leq v(S'\cup\{i\})-v(S')$. 
%\end{definition}

%In contrast, under the dual definition of the \emph{submodular} case, the marginal contribution $v(S \cup \{i\}) - v(S)$ is a monotone non-increasing function:


%\begin{definition}
%    We say that a value function $v$ maintains \emph{submodularity} if for any $S, S'\subseteq\{1,\ldots,n\}$ such that $S\subseteq S'$ and any $i\not\in S'$ it holds that: $v(S\cup\{i\})-v(S)\geq v(S'\cup\{i\})-v(S')$. 
%\end{definition}

\section{Unraveling New Properties of the Global Value Functions}

Prior work shows that in the local setting of non-probabilistic explanations, subset-minimal sufficient or contrastive reasons can be computed thanks to monotonicity, which holds only in the restricted case $\delta=1$. This assumption, however, is highly limiting and lacks practical flexibility. While one might hope to generalize to probabilistic guarantees for arbitrary $\delta$, prior results demonstrate that monotonicity breaks down in this setting, rendering the computation of explanations computationally harder~\citep{arenas2022computing, kozachinskiy2023inapproximability, izza2023computing, subercaseaux2024probabilistic, izza2024locally}.

%Prior art has shown that in the \emph{local} setting of \emph{non-probabilistic} explanations, a key property that enables the computation of subset-minimal sufficient or contrastive reasons is \emph{monotonicity}, which holds in this very restricted case of $\delta=1$~\citep{ignatiev2019abduction, BaMoPeSu20, wu2024verix, arenas2022computing}. However, this setting is significantly limited, as it assumes a stringent framework and requires $\delta = 1$, which significantly reduces its practical flexibility. While one might hope to overcome this limitation by adopting probabilistic guarantees for \emph{any} $\delta$, a key theoretical insight shown by previous works is that, unfortunately, \emph{monotonicity breaks down} in the local probabilistic setting~\citep{arenas2022computing, izza2023computing, subercaseaux2025probabilistic, izza2024locally} --- making the computation of explanations in this general setting computationally harder~\citep{arenas2022computing, kozachinskiy2023inapproximability}.

At the even more extreme case of the \emph{global} and \emph{non-probabilistic} setting,~\cite{bassanlocal} demonstrated the stringent \emph{uniqueness} property --- i.e., there is exactly \emph{one} subset-minimal explanation. However, requiring $\delta=1$ makes this setting highly restrictive, especially in the global case, where the explanation conditions must hold for \emph{all} inputs. In fact,~\cite{bassanlocal} proves that this unique minimal subset is equivalent to the subset of all features that are not strictly redundant (which may, in practice, be all of them). We prove that in the general global \emph{probabilistic} setting --- for any $\delta$ --- this uniqueness property actually does \emph{not} hold, and the number of subsets can even be \emph{exponential}.%we prove that in the more general global probabilistic setting --- valid for \emph{any} $\delta$ --- this uniqueness property actually does \emph{not} hold, and there can be actually be an \emph{exponential} number of such subsets.

\begin{restatable}{proposition}{nonuniqueness} \label{nonuniqueness_prop}
While the non-probabilistic case ($\delta = 1$) admits a unique subset-minimal global sufficient reason~\citep{bassanlocal}, in the general probabilistic setting (for arbitrary $\delta$), there exist functions $f$ that have $\Theta (\frac{2^n}{\sqrt{n}})$ subset-minimal global sufficient reasons.
\end{restatable}

Interestingly, although the uniqueness property fails to hold in the general case for arbitrary $\delta$ (and is restricted to the special case of $\delta=1$), we show that the crucial \emph{monotonicity} property holds for the global value function across all values of $\delta$. This applies to both the sufficient ($v^g_{\text{suff}}$) and contrastive ($v^g_{\text{con}}$) value functions. This finding is surprising as it stands in sharp contrast to the local setting, where the corresponding value functions ($v^{\ell}_{\text{suff}}$ and $v^{\ell}_{\text{con}}$) do not satisfy this property:

\begin{restatable}{proposition}{globalmonotonicity} \label{globalmonotonicity_prop}
While the local probabilistic setting (for any $\delta$) lacks monotonicity --- i.e., the value functions $v^{\ell}_{\text{con}}$ and $v^{\ell}_{\text{suff}}$ are non-monotonic~\citep{arenas2022computing, izza2023computing, subercaseaux2024probabilistic, izza2024locally} --- in the global probabilistic setting (also for any $\delta$), both value functions $v^g_{\text{con}}$ and $v^g_{\text{suff}}$ are monotonic non-decreasing.
\end{restatable}



%In the \emph{local} setting of \emph{non-probabilistic} explanations, a key property that enables one to obtain subset minimal sufficient or contrastive reasons is the \emph{monotonicity} property that holds under this setting~\citep{}. The intuition for convergence to subset-minimal explanations is that any simple greedy algorithm for a monotonic function will always converge to a subset minimal explanation. However, this very limited scenario which is non-probabiilistic and assumes explicitly that it must hold that $\delta:=1$ substantially effects the flexibility of this definition. This could be mitigated by alternating to probabilistic guarantees, but a key theoretical insight that was brought regarding the local versions of these explanations is that \emph{monotonicity does not hold} for the local probabilistic setting, which renders the computation of explanations in this setting computationally hard~\citep{}.

%Beyond the surprising insight that monotonicity holds for the global value function --- but not for the local one --- we identify additional structural properties unique to the global setting, such as submodularity and supermodularity. In particular, under the common assumption of feature independence, the global sufficient value function $v^{g}{\text{suff}}$ is supermodular. In contrast, the local counterpart $v^{\ell}{\text{suff}}$ lacks this property even under the much stronger assumption of a uniform (and hence independent) input distribution. Specifically:


Beyond the surprising insight that monotonicity holds for the global value functions --- but fails to hold in the local one --- we identify additional structural properties unique to the global setting, including submodularity or supermodularity. In particular, we show that under the common assumption of feature independence, the global sufficient value function $v^{g}_{\text{suff}}$ exhibits supermodularity. In contrast, its local counterpart $v^{\ell}_{\text{suff}}$ fails to exhibit this property even under the much more restrictive assumption of a uniform (and hence independent) input distribution. Specifically:

\begin{restatable}{proposition}{globalsupermodular} \label{globalsupermod_prop}
While the local probabilistic sufficient setting (for any $\delta$) lacks supermodularity --- even when $\mathcal{D}$ is uniform, i.e., the value function $v^{\ell}_{\text{suff}}$ is not supermodular --- in the global probabilistic setting (also for any $\delta$), when $\mathcal{D}$ exhibits feature independence, the value function $v^g_{\text{suff}}$ is supermodular.
\end{restatable}


Interestingly, for the second family of explanation settings --- specifically, that of obtaining a global probabilistic \emph{contrastive} reason --- we show that the corresponding value function is not supermodular, but rather \emph{submodular}. This result is particularly surprising when contrasted with the local setting, where the value function is neither submodular nor supermodular.

\begin{restatable}{proposition}{globalsubmodular} \label{globalsubmodular_prop}
While the local probabilistic contrastive setting (for any $\delta$) lacks submodularity --- even when $\mathcal{D}$ is uniform, i.e., the value function $v^{\ell}_{\text{con}}$ is not submodular --- in the global probabilistic setting (also for any $\delta$), when $\mathcal{D}$ exhibits feature independence, the value function $v^g_{\text{con}}$ is submodular.
\end{restatable}



%\begin{proposition}
%        All of the probabilistic value functions: $v^{\ell}_{suff}$, $v^{\ell}_{con}$, $v^g_{suff}$ and $v^{g}_{con}$ satisfy saturation.
%\end{proposition}

%It is interesting to note that the convergence 

%\begin{observation}
%    The convergence of simple greedy algorithms to subset-minimal explanations in the \textbf{non-probabilistic} are a direct result of the saturation properties.
%\end{observation}

%A very significant proof that we show is that although the \emph{local} value functions where shown to be non-monotonic by previous works~\citep{}, which is the key in showing substantial hardness results in obtaining even subset-minimal explanations, the \emph{global} variants of these value functions \emph{are} monotonic:





%\begin{proposition}
%    While both local value functions: $v^{\ell}_{suff}$, $v^{\ell}_{con}$ do not satisfy monotonicity, both global value functions $v^g_{suff}$ and $v^{g}_{con}$ satisfy non-decreasing monotonicity.
%\end{proposition}

%\begin{proposition}
%    While both local value functions: $v^{\ell}_{suff}$, $v^{\ell}_{con}$ do not satisfy sub-modularity or super-modularity, the following behavior holds for the global value functions: \begin{inparaenum}[(i)] \item
%    $v^g_{suff}$ satisfies super-modularity, and \item $v^{g}_{con}$ satisfies sub-modularity.\end{inparaenum}
%\end{proposition}

%\begin{theorem}
%    Both the local and global probability functions $v_l$ and $v_g$ satisfy saturation.
%\end{theorem}

\section{Computational Complexity Results}

\subsection{Subset Minimal Explanations}
\label{subsection_subset_minimal}

In this section, we examine the complexity of obtaining \emph{subset minimal} explanations (local/global, sufficient/contrastive) across the different model types analyzed. The key property at play here is the \emph{monotonicity} of the value function. The previous section established that monotonicity holds for both global value functions, $v^g_{\text{con}}$ and $v^g_{\text{suff}}$, but does not hold for the local value functions, $v^{\ell}_{\text{con}}$ and $v^{\ell}_{\text{suff}}$. This distinction is crucial in showing that a greedy algorithm converges to a subset minimal explanation in the global setting but fails in the local setting. As a result, we will showcase the surprising finding that computing various local subset-minimal explanation forms is hard, whereas computing subset-minimal global explanation forms is tractable (polynomial-time solvable). We will begin by introducing the following generalized greedy algorithm:


\begin{algorithm}
	\algnewcommand\algorithmicforeach{\textbf{for each}}
	\algdef{S}[FOR]{ForEach}[1]{\algorithmicforeach\ #1\ \algorithmicdo}
	\textbf{Input} Value function $v$, and some $\delta\in[0,1]$
	\caption{Subset Minimal Explanation Search}\label{alg:subset-minimal}
	\begin{algorithmic}[1]
		\State $S \gets \{1,\ldots, n\}$
        \While {$\min_{i\in S}v(S\setminus\{i\})\geq\delta$}
        \State $j\gets \argmax_{i\in S}v(S\setminus\{i\})$
        \State $S\gets S\setminus \{j\}$


\EndWhile		
		\State \Return $S$ \Comment{$S$ is a (subset minimal?) $\delta$-explanation}
	\end{algorithmic}
\end{algorithm}

Algorithm~\ref{alg:subset-minimal} aims to obtain a subset-minimal $\delta$-explanation. We begin the algorithm with the subset $S$ initialized as the entire input space $\{1,\ldots, n\}$. Iteratively, we check whether the minimal value that $v(S\setminus \{i\})$ can attain exceeds $\delta$. In each iteration, we remove a feature $j$ from $S$ that minimizes the decrease in the value function, selecting the feature $j$ that maximizes $v(S\setminus \{j\})$. Once this iterative process concludes, we return $S$.

The key determinant of whether Algorithm~\ref{alg:subset-minimal} yields a subset-minimal explanation is the \emph{monotonicity} property of the value function $v$. The algorithm concludes with a phase in which removing any individual feature from $S$ results in $v(S\setminus \{i\})$ being smaller than $\delta$. However, monotonicity ensures that this holds for any $v(S\setminus S')$, providing a significantly stronger guarantee. Given the monotonicity property of the value functions established in the previous sections, we derive the following proposition:










%Algorithm~\ref{alg:subset-minimal} opts to obtain a subset minimal $\delta$-explanation. It starts with the subset $S$ fixed to the entire input space $\{1,\ldots, n\}$. We iteratively check whether the minimal value that $v(S\setminus \{i\})$ can obtain is larger than $\delta$, and at each iteration reduce some feature $j$ from $S$ that causes the value function to decrease as least as possible (choosing the feature $j$ that gives us the maximal possible value of $v(S\setminus \{j\})$. After we finish with this iterative process, we return $S$.

%The main factor that will determine whether Algorithm~\ref{alg:subset-minimal} produces a subset minimal explanation or not is the \emph{monotonicity} property of the value function $v$. The algorithm ends in a phase where removing any individual feature from $S$ causes $v(S\setminus \{i\})$ to be smaller than $\delta$, however montonicity will imply that this is true for any $v(S\setminus S')$, which is a significantly stronger guarantee. Due to the montonicity proeprty of the value functions that we proved in the previous sections, we derive in the following theorem:

\begin{restatable}{proposition}{subsetminimalalgconv}\label{subset_minimal_alg_conv_prop}
    Computing Algorithm~\ref{alg:subset-minimal} with the \emph{local} value functions $v^{\ell}_{\text{con}}$ and $v^{\ell}_{\text{suff}}$ does not always converge to a subset minimal $\delta$-sufficient/contrastive reason. However, computing it with the \emph{global} value functions $v^{g}_{\text{con}}$ or $v^{g}_{\text{suff}}$ necessarily produces subset minimal $\delta$-sufficient/contrastive reasons. 
\end{restatable}



Building on this result, we proceed to establish new complexity findings for deriving various forms of subset-minimal explanations within our framework, considering the different analyzed model types.










%\begin{algorithm}
	%\algnewcommand\algorithmicforeach{\textbf{for each}}
	%\algdef{S}[FOR]{ForEach}[1]{\algorithmicforeach\ #1\ \algorithmicdo}
	%\textbf{Input} Value function $v$, and some $\delta\in[0,1]$
	%\caption{Subset Minimal Explanation Search}\label{alg:subset-minimal-local}
	%\begin{algorithmic}[1]
	%	\State $S \gets \emptyset$
     %   \State $F \gets \{1,\ldots, n\}$
     %   \While {\textbf{True}}
     %   \If{$v(S)\geq \delta$}\label{lst:line:sufficientline}
	%	\State \textbf{Break}
	%	\EndIf
    %    \State $j\gets \argmax_{i\in \Bar{S}}v(S\cup\{i\})$
    %    \State $S\gets S\cup \{j\}$


%\EndWhi5le		
		%\State \Return $S$ \Comment{$S$ is a %(subset minimal?) %$\delta$-reason}
%	\end{algorithmic}
%\end{algorithm}


%\begin{theorem}
%    Running Algorithm~\ref{} with the global value functions
%\end{theorem}




%\begin{algorithm}
	%\algnewcommand\algorithmicforeach{\textbf{for each}}
	%\algdef{S}[FOR]{ForEach}[1]{\algorithmicforeach\ #1\ \algorithmicdo}
	%\textbf{Input} $f$, $\x$
	%\caption{Local Subset Minimal Sufficient Reason}\label{alg:subset-minimal-local}
	%\begin{algorithmic}[1]
%		\State $S \gets \{1,\ldots,n\}$
%		\ForEach {$i \in \{1,...,n\}$ by some arbitrary ordering}\label{lst:line:orderinglinelocal}
%		\ForEach {$i \in \{1,...,n\}$ by some arbitrary ordering}\label{lst:line:orderinglinelocal}
        
		%\If{$\mySuff(f,\x,S\setminus \{i\})=1$}\label{lst:line:sufficientline}
		%\State $S \gets S\setminus \{i\}$
		%\EndIf
		%\EndFor\label{lst:line:endregularupper}
		%\State \Return $S$ \Comment{$S$ is a subset minimal \emph{local} sufficient reason}
	%\end{algorithmic}
%\end{algorithm}

\textbf{Decision trees.} We begin by examining a highly simplified and ostensibly ``interpretable'' scenario. Specifically, we assume that the model $f$ is a decision tree and that the distribution $\mathcal{D}$ is independent. Within this simplified setting, we demonstrate a strict separation:~\cite{arenas2022computing} established the surprising intractability result that, unless PTIME$=$NP, no polynomial-time algorithm exists for computing a subset-minimal local $\delta$-sufficient reason for decision trees under independent distributions (even under the uniform distribution). In contrast, we demonstrate the unexpected result that this exact problem can be solved efficiently in the global setting, meaning that a subset-minimal global $\delta$-sufficient reason for decision trees can indeed be computed in polynomial time. Formally:


\begin{restatable}
{theorem}{subsetminimaldt}\label{subset_minimal_dt_theor}
If $f$ is a decision tree and the probability term $v^{g}_{\text{suff}}$ can be computed in polynomial time given the distribution $\mathcal{D}$ (which holds for independent distributions, among others), then obtaining a subset-minimal \emph{global} $ \delta $-sufficient reason can be obtained in \emph{polynomial time}. However, unless PTIME$=$NP, \emph{no polynomial-time} algorithm exists for computing a \emph{local}  $\delta$-sufficient reason for decision trees even under independent distributions.
\end{restatable}

\textbf{Extension to other tractable models.} We additionally note that a similar tractability guarantee for computing global explanations also holds for \emph{orthogonal DNFs}~\citep{crama2011boolean}, which we briefly recall are Disjunctive Normal Form formulas whose terms are pairwise mutually exclusive. Establishing the result for this class is useful because orthogonal DNFs \emph{generalize} decision trees while preserving their clean structural properties, allowing our guarantees to extend beyond tree-structured models. For completeness, we provide the full argument in Appendix~\ref{app:odnf}.


%\begin{remark}
%This result extends to Orthogonal DNFs~\cite{crama2011boolean}, as their terms are pairwise inconsistent and analogous to paths in a decision tree.
%\end{remark}

%\begin{theorem}
%If $f$ is a decision tree and the probability term $v^{g}_{\text{suff}}$ can be computed in polynomial time given the distribution $\mathcal{D}_p$ (which holds for independent distributions, among others), then obtaining a subset-minimal \textit{{\textbf{global}}} $ \delta $-sufficient reason \textit{{\textbf{can be obtained in polynomial time}}}. However, unless P$=$NP, \textit{{\textbf{no polynomial-time algorithm exists}}} for computing a \textit{{\textbf{local}}}  $\delta$-sufficient reason for decision trees even under independent distributions.
%\end{theorem}


%\begin{theorem}
%If $ f$ is a decision tree and the probability terms $v^{g}_{suff}$ or $ v^{g}_{con}$ can be computed in polynomial time given the distribution $\mathcal{D}_p$ (which holds for independent distributions, among others), then obtaining a subset-minimal \textit{{\textbf{global}}} $ \delta $-sufficient or $\delta$-contrastive reason \textit{{\textbf{can be done in polynomial time}}}. However, unless P$=$NP, \textit{{\textbf{no polynomial-time algorithm exists}}} for computing a \textit{{\textbf{local}}}  $\delta$-contrastive or $ \delta$-sufficient reason for decision trees even under independent distributions.
%\end{theorem}

\textbf{Neural networks, tree ensembles, and other complex models.} We now extend our previous results to more complex models beyond decision trees. Specifically, we will demonstrate that when the distribution $\mathcal{D}$ is derived from empirical distributions, and under the fundamental assumption that the model $f$ allows polynomial-time inference, it follows that computing a subset-minimal global $\delta$-sufficient and contrastive reason can be done in polynomial time.

\begin{restatable}{proposition}{empiricalsubsetminimal}\label{empirical_subset_minimal_prop}
For any model $f$, and empirical distribution $\mathcal{D}$ --- computing a subset-minimal \emph{global} $\delta$-sufficient or $\delta$-contrastive reason for $f$ can be done in polynomial time. 
\end{restatable}

This strong complexity outcome, which holds for \emph{any} model under an empirical data distribution assumption, allows us to further differentiate the complexity of local and global explanation settings. Specifically, for certain complex models, computing subset-minimal \emph{local} explanations remains intractable even when restricted to empirical distributions. We establish this fact for both neural networks and tree ensembles, leading to the following theorem on a strict complexity separation:


\begin{restatable}{theorem}{seperationsubsetminemp}\label{seperation_subset_min_emp_theor}
    Assuming $f$ is a neural network or a tree ensemble, and $\mathcal{D}$ is an empirical distribution --- there exist \emph{polynomial-time} algorithms for obtaining subset minimal \emph{global} $\delta$-sufficient and contrastive reasons. However, unless PTIME$=$NP, there is \emph{no polynomial time} algorithm for computing a subset minimal \emph{local} $\delta$-sufficient reason or a subset minimal \emph{local} $\delta$-contrastive reason.
\end{restatable}

%This allows us to provide additional complexity seperations between the local and global settings of explanations for some complex models such as neural networks and tree ensembles. Particularly, we will prove that obtaining a subset minimal $\delta$ sufficient or contrastive reason, even under empirical distributions is  coNP and NP Hard. Hence, give our result in Theorem~\citep{}, we establish the following result: 


\subsection{Approximate Cardinally Minimal Explanations}

In this subsection, we shift our focus from subset-minimal sufficient/contrastive reasons to the even more challenging task of finding a \emph{cardinally minimal} $\delta$-sufficient/contrastive reason. We will demonstrate that in the global setting, the interplay between supermodularity/submodularity and monotonicity of the value function enables us to establish novel \emph{provable approximations} for computing explanations. In contrast, we will show that computing these explanations in the local setting remains intractable. This result further strengthens the surprising distinction between the tractability of computing global explanations versus local ones.

\textbf{A unified greedy approximation algorithm.} When working with a non-decreasing monotonic submodular function, the problem of identifying a cardinally minimal explanation closely aligns with the \emph{submodular set cover} problem~\citep{wolsey1982analysis}. In contrast, employing a \emph{supermodular} function leads to a non-submodular variation of this problem~\citep{shi2021minimum}. These problems have garnered significant interest due to their strong approximation guarantees~\citep{wolsey1982analysis, iyer2013submodular, chen2023bicriteria}. In the context of submodular set cover optimization, a standard approach involves using a classic greedy algorithm, which we will first outline (Algorithm~\ref{alg:cardinally-minimal}). This algorithm serves as the foundation for approximating a cardinally minimal sufficient or contrastive $\delta$-reason, and we will later examine its specific guarantees in both cases.

\begin{algorithm}
	\algnewcommand\algorithmicforeach{\textbf{for each}}
	\algdef{S}[FOR]{ForEach}[1]{\algorithmicforeach\ #1\ \algorithmicdo}
	\textbf{Input} Value function $v$, and some $\delta\in[0,1]$
	\caption{Cardinally Minimal Explanation Approximation Search}\label{alg:cardinally-minimal}
	\begin{algorithmic}[1]
		\State $S \gets \emptyset$
        \While {$\max_{i\in S}v(S\cup\{i\})<\delta$}
        \State $j\gets \argmax_{i\in S}v(S\cup\{i\})$
        \State $S\gets S\cup \{j\}$
\EndWhile		
		\State \Return $S$ \Comment{$|S|$ is a (provable approximation?) of a cardinally minimal $\delta$-explanation}
	\end{algorithmic}
\end{algorithm}

Algorithm~\ref{alg:cardinally-minimal} closely resembles Algorithm~\ref{alg:subset-minimal}, but works bottom-up rather than top-down. It starts with an empty subset and incrementally adds features, each time selecting the one that minimizes the increase in $v(S \cup {i})$, stopping when adding any feature would push the value over $\delta$.


%Algorithm~\ref{alg:cardinally-minimal} closely resembles Algorithm~\ref{alg:subset-minimal}, but runs bottom-up rather than top-down. It starts with an empty subset and iteratively adds the feature that yields the smallest increase in $v(S \cup {i})$, stopping once any addition would push the value past $\delta$.



%Algorithm~\ref{alg:cardinally-minimal} closely resembles Algorithm~\ref{alg:subset-minimal}, but it operates in a bottom-up manner rather than top-down. It begins with an empty explanation subset and incrementally adds features, selecting at each step the feature that results in the smallest increase in $v(S \cup \{i\})$. The process continues until adding any additional feature would cause the threshold to exceed $ \delta$.


\textbf{Cardinally minimal contrastive reasons.} We begin with global contrastive reasons, where monotonicity and submodularity hold, reducing the task to the classic \emph{submodular set cover} problem and allowing us to apply~\cite{wolsey1982analysis}'s classic guarantee via Algorithm~\ref{alg:cardinally-minimal}. For integer-valued objectives, the algorithm achieves a Harmonic-based factor, and more generally an $\mathcal{O}\!\left(\ln\!\left(\tfrac{v([n])}{\min_{i \in [n]} v(i)}\right)\right)$-approximation. Under an \emph{empirical} distribution $\mathcal{D}$ with fixed sample size, this becomes a \emph{constant approximation}, only \emph{logarithmic} in the sample size, yielding a substantially strong guarantee. By contrast, in the \emph{local} setting --- even for a \emph{single} sample point --- no \emph{bounded} approximation exists, marking a sharp gap between global and local cases.


%\textbf{Cardinally minimal contrastive reasons.} We begin with contrastive reasons. In the global setting, both monotonicity and submodularity hold, reducing the task to the well-studied \emph{submodular set cover} problem. Thus, finding cardinality-minimal $\delta$-contrastive reasons reduces to submodular set cover, letting us apply~\cite{wolsey1982analysis}'s classic guarantee via Algorithm~\ref{alg:cardinally-minimal}. For integer-valued objectives, the algorithm achieves a Harmonic-based factor, and more generally an $\mathcal{O}!\left(\ln!\left(\tfrac{v([n])}{\min_{i \in [n]} v(i)}\right)\right)$-approximation. Under an \emph{empirical} distribution $\mathcal{D}_p$ with fixed sample size, this becomes a \emph{constant approximation}, only \emph{logarithmic} in the sample size, yielding a substantially strong approximation factor. By contrast, in the \emph{local} setting, even for a \emph{single} point, no \emph{bounded} approximation exists, highlighting a significantly sharp gap between global and local settings.

%\textbf{Cardinally minimal contrastive reasons.} We begin with contrastive reasons. In the global setting, both monotonicity and submodularity hold, connecting the problem to the well-studied \emph{submodular set cover} problem. Thus, finding cardinality-minimal $\delta$-contrastive reasons reduces to submodular set cover, letting us apply~\cite{wolsey1982analysis}'s classic approximation guarantee via Algorithm~\ref{alg:cardinally-minimal}. For integer-valued objectives, the algorithm achieves a Harmonic-based factor, and more generally an $\mathcal{O}\left(\ln\left(\frac{v([n])}{\min_{i \in [n]} v({i})}\right)\right)$-approximation. Under an \emph{empirical} distribution $\mathcal{D}_p$ with fixed sample size, this becomes a \emph{constant} approximation, where the factor is only \emph{logarithmic} in the sample size, providing a substantially strong approximation factor. In sharp contrast, in the \emph{local} setting — even for a \emph{single} point — no \emph{bounded} approximation exists.


%We begin with contrastive reasons. Our results show that both monotonicity and submodularity hold in the global setting, enabling a connection to the well-studied \emph{submodular set cover} problem~\citep{wolsey1982analysis}. Specifically, finding cardinality-minimal $\delta$-contrastive reasons aligns with submodular set cover, allowing us to apply~\cite{wolsey1982analysis}'s classic approximation guarantee, achieved via Algorithm~\ref{alg:cardinally-minimal}. When the objective is integer-valued, the algorithm ensures a Harmonic-based approximation factor, and more generally yields an $\mathcal{O}\left(\ln\left(\frac{v([n])}{\min_{i \in [n]} v({i})}\right)\right)$-approximation. Moreover, under an \emph{empirical} distribution $\mathcal{D}_p$ with a fixed sample size, this becomes a \emph{constant} approximation. In sharp contrast, we show that in the \emph{local} setting --- even for a \emph{single} point --- no \emph{bounded} approximation guarantee exists.






%We begin our analysis with contrastive reasons. Notably, our new results establish that both monotonicity and submodularity hold in the global setting, allowing us to connect the problem to the well-studied \emph{submodular set cover} problem~\citep{wolsey1982analysis}. Specifically, since the task of finding cardinality-minimal $\delta$-contrastive reasons aligns with submodular set cover, we can apply the classic approximation guarantee by Wolsey et al.~\citep{wolsey1982analysis}, achieved via Algorithm~\ref{alg:cardinally-minimal}. When the objective function is integer-valued, this algorithm guarantees an approximation factor based on the Harmonic series, and more generally yields an $\mathcal{O}\left(\ln\left(\frac{v([n])}{\min_{i \in [n]} v(\{i\})}\right)\right)$-approximation. Furthermore, if we assume an \emph{empirical} distribution $\mathcal{D}_p$ with a fixed number of samples, we can show that this approximation becomes \emph{constant}. In sharp contrast, we show that in the \emph{local} setting --- even with just a \emph{single} data point --- the problem lacks any \emph{bounded} approximation guarantee.




%\begin{theorem}
%    Assuming that $\mathcal{D}_p$ is an empirical distribution over a finite dataset $D\subseteq \mathbb{F}$, then, algorithm~\ref{} provides an $O(ln(|\text{D}|))$ approximation for obtaining a global cardinally minimal $\delta$-contrastive reason for $f$.
%\end{theorem}


\begin{restatable}{theorem}{approxcardcon}\label{approx_card_con_prop}
Given a neural network or tree ensemble $f$ and an empirical distribution $\mathcal{D}$ over a fixed dataset $\text{D}$, Algorithm~\ref{alg:cardinally-minimal} yields a \emph{constant} $\mathcal{O}\left(\ln\left(\frac{v^g_{con}([n])}{\min_{i \in [n]} v^g_{con}(\{i\})}\right)\right)$-approximation, bounded by $\mathcal{O}(\ln(|\text{D}|))$, for computing a global cardinally minimal $\delta$-contrastive reason for $f$, assuming feature independence. In contrast, unless PTIME$=$NP, \emph{no bounded approximation} exists for computing a local cardinally minimal $\delta$-contrastive reason for any $\langle f,\x\rangle$, even when $|\text{D}| = 1$.
\end{restatable}

We also provide a matching \emph{lower bound} for the bound in Theorem~\ref{approx_card_con_prop} in the specific case of an empirical distribution without any independence assumption (see Appendix~\ref{section_hardness_of_approximation}).

\textbf{Cardinally minimal sufficient reasons.} Unlike the submodular set cover problem --- linked to cardinally minimal global \emph{contrastive} reasons and admitting strong approximations --- the supermodular variant, tied to global \emph{sufficient} reasons, is harder to approximate. Still, it offers guarantees when the function has bounded curvature~\citep{shi2021minimum}. The \emph{total curvature} of a function $v:2^{n}\to \mathbb{R}$ is defined as $k^f := 1 - \min_{i \in [n]} \frac{v([n]) - v([n] \setminus {i})}{v({i}) - v(\emptyset)}$. Leveraging results from~\cite{shi2021minimum}, we show that Algorithm~\ref{alg:cardinally-minimal} achieves an $\mathcal{O}\left(\frac{1}{1 - k^f} + \ln\left(\frac{v([n])}{\min_{i \in [n]} v({i})}\right)\right)$-approximation.


%In contrast to the submodular set cover problem --- which corresponds to finding cardinally minimal global \emph{contrastive} reasons and admits strong approximation guarantees --- the supermodular set cover problem, aligned with cardinally minimal global \emph{sufficient} reasons, is more challenging to approximate. However, it still offers approximation guarantees when the underlying function has bounded curvature~\citep{shi2021minimum}. We define the \emph{total curvature} of some function $v:2^{n}\to \mathbb{R}$ as $k^f := 1 - \min_{i \in [n]} \frac{v([n]) - v([n] \setminus {i})}{v({i}) - v(\emptyset)}$. Using the results by Shi et al.~\citep{shi2021minimum}, we show that Algorithm~\ref{alg:cardinally-minimal} achieves an approximation factor of $\mathcal{O}\left(\frac{1}{1 - k^f} + \ln\left(\frac{v([n])}{\min_{i \in [n]} v({i})}\right)\right)$.


Notably, under a fixed empirical distribution, the approximation becomes constant. While contrastive reasons admit a tighter $\mathcal{O}(\ln(|\text{D}|))$ bound, sufficient reasons incur an extra $\frac{1}{1 - k^f}$ factor --- yet still yield a constant-factor approximation. In sharp contrast, the local variant remains inapproximable, lacking any \emph{bounded} approximation even when $|\text{D}| = 1$.



%Notably, under an empirical distribution with a fixed number of data points, the approximation becomes constant. While the contrastive case achieves a tighter $\mathcal{O}(\ln(|\text{D}|))$ bound, the sufficient reason case incurs an extra $\frac{1}{1 - k^f}$ factor—yet still yields a constant-factor approximation. In sharp contrast, the local variant remains inapproximable, with no \emph{bounded} approximation, even under simple conditions like $|\text{D}| = 1$.

%Notably, under an empirical distribution with a fixed number of data points, this approximation becomes constant. However, unlike the contrastive setting which attains a tighter $\mathcal{O}(\ln(|\text{D}|))$ bound, the sufficient reason setting includes an additional $\frac{1}{1 - k^f}$ factor. Nonetheless, we still prove that it achieves a constant-factor approximation. In stark contrast, we show that the local version of this problem is significantly harder: it remains inapproximable --- even without a bounded approximation --- under highly simplified conditions such as when $|\text{D}| = 1$.




%In contrast to the submodular set-cover problem (which aligns with the task of obtaining cardinally minimal global contrastive reasons), the supermodular set cover problem (which aligns with cardinally minimal global sufficient reasons) is harder to approximate but also contains approximation guarantees when the \emph{curvature} of the underlying function is bounded~\citep{shi2021minimum}. We define the \emph{total curvature}: $k^f:=1-\min_{i\in [n]}\frac{v([n])-v([n]\setminus \{i\})}{v(\{i\})-v(\emptyset)}$, and can use the approximation guarantee obtained by Shi et al.~\citep{shi2021minimum} to show that Algorithm~\ref{alg:cardinally-minimal} provides a $O\left(\frac{1}{1-k^f}+\ln\left(\frac{v([n])}{\min_{i \in [n]} v(\{i\})}\right)\right)$ approximation. We particularly show that when assuming an \emph{empirical} distribution with a fixed number of data points this approximation is \emph{constant} (though here it is bounded by $O(|\text{D}|^2)$ where $|D|$ is the number of data-points, in contrast to the much better $O(\ln(|D|))$ in the contrastive explanation setting. However, we prove, once again that in stark contrast, the \emph{local} congfiguration of this task obtains substantial intractability - and does not even maintain a \emph{bounded} approximation, even in a highly simplified setting of $|\text{D}|=1$.


\begin{restatable}{theorem}{approxcardsuff}\label{approx_card_suff_prop}

%Given a neural network or tree ensemble $f$ and an empirical distribution $\mathcal{D}p$ over a fixed dataset $\text{D}$, Algorithm~\ref{alg:cardinally-minimal} achieves a \textbf{constant} $\mathcal{O}\left(\frac{1}{1-k^f}+\ln\left(\frac{v^g{\text{suff}}([n])}{\min_{i \in [n]} v^g_{\text{suff}}({i})}\right)\right)$-approximation for computing a global cardinally minimal $\delta$-sufficient reason for $f$. In contrast, unless PTIME$=$NP, \textbf{no bounded approximation} exists for computing a local cardinally minimal $\delta$-sufficient reason for any $\langle f,\x\rangle$, even when $|\text{D}| = 1$.
    Given a neural network or tree ensemble $f$ and an empirical distribution $\mathcal{D}$ over a fixed dataset $\text{D}$, Algorithm~\ref{alg:cardinally-minimal} yields a \emph{constant} $\mathcal{O}\left(\frac{1}{1-k^f}+\ln\left(\frac{v^g_{\text{suff}}([n])}{\min_{i \in [n]} v^g_{\text{suff}}(\{i\})}\right)\right)$-approximation for computing a global cardinally minimal $\delta$-sufficient reason for $f$, assuming feature independence. In contrast, unless PTIME$=$NP, there is \emph{no bounded approximation} for computing a local cardinally minimal $\delta$-sufficient reason for any $\langle f,\x\rangle$, even when $|\text{D}| = 1$.
\end{restatable}


Overall, these findings strengthen Subsection~\ref{subsection_subset_minimal}, which showed the tractability of computing subset-minimal global explanations in stark contrast to local ones. Here, we further show that approximating cardinally minimal global explanations is tractable, unlike their inapproximable local counterparts.

%Overall, these findings strengthen the results from Section~\ref{alg:subset-minimal}, which established the tractability of computing subset-minimal global explanations in stark contrast to local ones. Specifically, in this section we have demonstrated that approximating cardinally minimal global explanations is tractable, in stark contrast to the inapproximability of their local counterparts.


%\section{Limitations}

\section{Limitations and Future Work}
While many of our most important findings --- particularly those concerning fundamental properties of value functions --- hold generally, we also instantiate them to yield concrete complexity results for specific model classes (e.g., neural networks), distributional assumptions (e.g., empirical distributions), and explanation definitions within our framework. Naturally, other potential settings remain open for analysis. Nonetheless, we believe our findings offer compelling insights into foundational aspects of explanations, along with new tractability and intractability results, which together lay a strong foundation for investigating a broader range of explainability scenarios in future work.

More specifically, our study brings forward several important open theoretical questions:

\begin{enumerate}
\item Firstly, it would be interesting to investigate whether some of the tractable complexity results we obtained for global explainability can be tightened for simpler model classes, such as decision trees, linear models, or other models with inherently tractable structure. Although local explanations for such models are often intractable to compute~\citep{arenas2022computing, subercaseaux2024probabilistic}, the global explanation forms we study in this work exhibit strong structural properties, enabling significantly improved complexity guarantees that may be tightened even further. Furthermore, tightening the approximation guarantees when working with empirical distributions represents a key open direction.
\item Secondly, when \emph{exact} computation of expectations is infeasible, it is natural to ask whether alternative techniques --- such as Fully Polynomial Randomized Approximation Schemes (FPRAS) for model classes like DNF formulas~\citep{meel2019not} (and thus for tree ensembles), or Monte Carlo–based approximations as the ones used in~\citep{subercaseaux2024probabilistic} --- can be employed to approximate these expectations. A key open challenge is determining how to maintain minimality guarantees under approximation error, and what tolerance levels still allow the guarantees to hold.
\item Finally, we believe that studying alternative notions of feature importance --- e.g., other loss measures such as KL divergence~\citep{conmy2023towards} --- may reveal similar patterns of monotonicity, submodularity, or supermodularity when transitioning from local to global explanations. Such properties could enable efficient feature selection under these alternative metrics as well.
\end{enumerate}

%More particularly, there are a few central open theoretical questions that remain open for exploration given our study: \begin{inparaenum}[(i)]
%\item First, it would be interesting to see if some of the tractable complexity results obtained in this work for the global explainability could be even further improved for simple model classes, such as decision trees, linear models, or other models with a tractable structure. \item Second, it remains interesting to see if when expectations cannot be computed exactly, other techniques such as Fully Polynomial Randomized Approximation Schemes (FPRAS) can be obtained for model classes such as DNF formulas (hence implying tractability for tree ensembles), or statistical techniques like the monte carlo approximation method used in~\citep{subercaseaux2025probabilistic} can be used to approximate these expectations. Here, a central question is on how to preserve the minimality guarantees to hold under some tolerance, when the calculation of the value term is not exact. \item Lastly, we believe that studying alternative notions of feature importance (e.g., other loss metrics such as KL divergence, etc.) exhibit a similar phenomenon of montoonicity, submodularity and supermodularity when shifting from the local to the global settings, which can enable efficient feature selection under these metrics.
%\end{inparaenum}

%While many of our most important findings --- particularly those related to fundamental properties of value functions --- hold generally, we instantiate some of these results to provide specific complexity results for particular model classes (e.g., neural networks), distributional assumptions (e.g., empirical distributions), or explanation definitions examined within our framework. Naturally, other potential settings remain open for analysis. Nonetheless, we believe our findings offer compelling insights into foundational aspects of explanations, along with new tractability and intractability results, which together lay a strong foundation for investigating a broader range of explainability scenarios in future work.


\section{Conclusion}
We present a unified framework for evaluating diverse explanations and reveal a stark contrast between local and global sufficient and contrastive reasons. Notably, while the \emph{local} explanation variants lack any structural form common in combinatorial optimization --- such as monotonicity, submodularity, or supermodularity, we prove that their \emph{global} counterparts exhibit precisely these crucial properties: \begin{inparaenum}[(i)] \item \emph{monotonicity}, \item \emph{submodularity} in the case of contrastive reasons, and \item \emph{supermodularity} for sufficient reasons. \end{inparaenum} These proofs form the basis for proving a series of surprising complexity results, showing that global explanations with provable guarantees can be computed efficiently, even for complex model classes such as neural networks. In sharp contrast, we prove that computing the corresponding local explanations remains NP-hard --- even in highly simplified scenarios. Altogether, our results uncover foundational properties of explanations and chart both tractable and intractable frontiers, opening new avenues for future research.

\section*{Acknowledgments}

This work was partially funded by the European Union
(ERC, VeriDeL, 101112713). Views and opinions expressed
are however those of the author(s) only and do not necessarily reflect those of the European Union or the European
Research Council Executive Agency. Neither the European
Union nor the granting authority can be held responsible for
them. This research was additionally supported by a grant
from the Israeli Science Foundation (grant number 558/24).


%We present a unified framework for evaluating diverse explanations and reveal a stark contrast between local and global sufficient and contrastive reasons. Notably, while the \emph{local} explanation variants lack any form of monotonicity, submodularity, or supermodularity, we prove that their \emph{global} counterparts exhibit crucial properties: \begin{inparaenum}[(i)] \item \emph{monotonicity}, \item \emph{submodularity} in the case of contrastive reasons, and \item \emph{supermodularity} for sufficient reasons. \end{inparaenum} These properties are fundamental in combinatorial optimization and enable efficient algorithmic treatment. Hence, these proofs form the basis for proving a series of surprising complexity results, showing that global explanations with provable guarantees can be computed efficiently, even for complex model classes such as neural networks. In sharp contrast, we prove that computing the corresponding local explanations remains NP-hard --- even in highly simplified scenarios. Altogether, our results uncover foundational properties of explanations and chart both tractable and intractable frontiers, opening new avenues for future research.







%In this work, we uncover a surprising contrast between the behavior of local and global sufficient and contrastive reasons. Specifically, we show that while the \emph{local} variants of these explanations lack any form of monotonicity, submodularity, or supermodularity, their \emph{global} counterparts exhibit key structural properties: \begin{inparaenum}[(i)] \item \emph{monotonicity}, \item \emph{submodularity} for contrastive reasons, and \item \emph{supermodularity} for sufficient reasons. \end{inparaenum} These findings lay the groundwork for a series of unexpected complexity results demonstrating that global explanations with provable guarantees can be computed efficiently, even for complex models like neural networks. In stark contrast, we prove that computing the same corresponding local explanations remains NP-hard, even in significantly simplified settings. Overall, our findings reveal fundamental properties of explanations and establish both tractability and intractability results for generating explanations across different settings, thereby paving the way for further research.



%In this work, we introduce a unified framework for analyzing the computational complexity of generating different types of explanations --- ranging from sufficient to contrastive, non-probabilistic to probabilistic, and local to global. We demonstrate that although none of the local value functions exhibit monotonicity, submodularity, or supermodularity at all, the global value functions retain important structural properties: \begin{inparaenum}[(i)] \item monotonicity, \item submodularity in the case of contrastive reasons, and \item supermodularity for sufficient reasons. \end{inparaenum} These insights form the foundation for a series of unexpected results on the efficient computability of global explanations with provable guarantees, even for highly complex models such as neural networks and tree ensembles. In contrast, we show that computing the corresponding local versions of these explanations remains NP-hard, even under substantially more simplified scenarios. Overall, our findings reveal fundamental properties of explanations and establish both tractability and intractability results for generating a wide range of explanations across different settings, thereby paving the way for further research.

\bibliography{iclr2026_conference}
\bibliographystyle{iclr2026_conference}


%In this work we provide a unifying framework for analyzing the complexity of generating different types of explanations including sufficient, contrastive, probabilistic, non-probabilistic and local vs. global explanations. We prove that while both local and global value functions maintain saturation, only the global value function maintains monotonicity and submodularity (for the contrastive case) and supermodularity (for the sufficient case). This insight is the base for proving a set of tractable results for obtaining global explanations with formal gurantees for a diverse set of ML models, including highly complex models such as neural networks and tree ensembles, while we also prove that obtaining the local versions of these explanations is NP-Hard. Overall, we believe that our study highlights fundamental and crucial properties and aspects regarding the feasability of obtaining local and global explanations for ML models, while opening the door for further investiagtion on this topic.




\newpage
\appendix
\onecolumn
\appendix


\setcounter{definition}{0}
\setcounter{proposition}{0}
\setcounter{theorem}{0}
\setcounter{lemma}{0}
\begin{center}\begin{huge} Appendix\end{huge}\end{center}    
\noindent{The appendix contains formalizations and proofs that were mentioned throughout the paper:}




\newlist{MyIndentedList}{itemize}{4}
\setlist[MyIndentedList,1]{%
	label={},
	noitemsep,
	leftmargin=0pt,
}
% \setlist[MyIndentedList]{%
	%     label={},
	%     noitemsep,
	%     }


\begin{MyIndentedList}

    \item 
    \textbf{Appendix~\ref{extended_related_work}} contains extended related work on Formal XAI.

    \item 
    \textbf{Appendix~\ref{models_and_distributions_app}} contains the formalizations of the models and distributions used in this work.
    \item 

    
\textbf{Appendix~\ref{non_uniqueness_appendix}} contains the proof of Proposition~\ref{nonuniqueness_prop}.

    \item 


\textbf{Appendix~\ref{globalmonotonicity_prop_appendix_sec}} contains the proof of Proposition~\ref{globalmonotonicity_prop}.

    \item 

\textbf{Appendix~\ref{supermodularity_proofs_appendix}} contains the proof of Proposition~\ref{globalsupermod_prop}.

    \item 

\textbf{Appendix~\ref{globalsubmodular_prop_sec}} contains the proof of Proposition~\ref{globalsubmodular_prop}.
    \item 


\textbf{Appendix~\ref{subset_minimal_alg_conv_prop_sec}} contains the proof of Proposition~\ref{subset_minimal_alg_conv_prop}.
    
    \item 

\textbf{Appendix~\ref{subset_minimal_dt_theor_sec}} contains the proof of Theorem~\ref{subset_minimal_dt_theor}.


    \item \textbf{Appendix~\ref{empirical_subset_minimal_prop_sec}} contains the proof of Proposition~\ref{empirical_subset_minimal_prop}.


    \item \textbf{Appendix~\ref{seperation_subset_min_emp_theor_sec}} contains the proof of Theorem~\ref{seperation_subset_min_emp_theor}.


        \item \textbf{Appendix~\ref{approx_card_con_prop_sec}} contains the proof of Theorem~\ref{approx_card_con_prop}.


        \item \textbf{Appendix~\ref{approx_card_suff_prop_sec}} contains the proof of Theorem~\ref{approx_card_suff_prop}.

        % \item \textbf{Appendix~\ref{app:approx_val_limit}} discusses the limitations of methods for approximating global and local value functions.
        
        \item \textbf{Appendix~\ref{app:llm-disclosure}} contains an LLM usage disclosure.

\end{MyIndentedList}

%\section{Limitations and Future Work}
%\label{limitations_appendix}


%\section{Extended Related Work}
%\label{related_work_appendix}

\section{Extended Related Work on Formal XAI}
\label{extended_related_work}

Our work falls within the line of research on \emph{formal explainable AI} (formal XAI)~\citep{marques2022delivering}, which studies explanations equipped with provable guarantees~\citep{yu2023eliminating, darwiche2022computation, darwiche2023logic, shih2018symbolic, azzolinbeyond, audemard2021computational, calautti2025complexity}. A central theme in this area concerns the computational complexity of generating such guaranteed explanations~\citep{BaMoPeSu20, bassan2026unifying, marzoukbassanshap, marzouk2025on, marzouk2024tractability, bassan2025additive, amir2024hard, blanc2021provably, blanc2022query, jaakkola2025explainability, amgoud2026axiomatic, geibinger2025and}. Since obtaining formal explanations is often computationally intractable for expressive models such as neural networks~\citep{adolfi2024complexity, BaMoPeSu20} and tree ensembles~\citep{ordyniak2024explaining, bassanmakes}, much of the literature has focused on more tractable model classes~\citep{marques2023no}, including decision trees~\citep{bouniausing, arenas2022computing, bounia2024enhancing, bounia2023approximating}, monotonic models~\citep{marques2021explanations, harzli2022cardinality}, linear models~\citep{marques2020explaining, subercaseaux2024probabilistic}, and additive models~\citep{bassan2026provably}.

Beyond complexity-theoretic investigations, substantial effort has also been devoted to computing formal explanations in practice using automated reasoning tools, such as MaxSAT and MILP solvers for tree ensembles~\citep{audemard2023computing, audemard2022preferred, ignatiev2022using, boumazouza2021asteryx}, and neural network verification frameworks~\citep{wang2021beta, wu2024marabou} to derive certified explanations for neural networks~\citep{wu2024verix, bassan2023towards, bassan2023formally, izza2024distance, hadad2026, ignatiev2019abduction}. Additional approaches seek to alleviate the inherent computational burden through abstractions~\citep{de2025faster, bassan2025explaining, boumazouza2026fame}, relaxations such as smoothing~\citep{xue2023stability, anani2025pixel, jin2025probabilistic}, or training-time interventions~\citep{alvarez2018towards, bassan2025explain, bassan2025self}.

%Our work belongs to the line of research on \emph{formal explainable AI} (formal XAI)~\citep{marques2022delivering}, which studies explanations endowed with provable guarantees~\citep{yu2023eliminating, darwiche2022computation, darwiche2023logic, shih2018symbolic, azzolinbeyond, audemard2021computational, calautti2025complexity}. A central question that revolves around formal XAI is the question of the computational complexity of obtaining explanations with provable guarantees~\citep{BaMoPeSu20, bassan2026unifying, marzoukbassanshap, marzouk2025on, marzouk2024tractability, bassan2025additive, amir2024hard, blanc2021provably, blanc2022query, jaakkola2025explainability, amgoud2026axiomatic}. Since many of these explanations are computationally intractable to obtain for complex models like neural networks~\citep{adolfi2024complexity, BaMoPeSu20} or tree ensembles~\citep{ordyniak2024explaining, bassanmakes}, much of the literature has focused on more tractable model classes~\citep{marques2023no}, including decision trees~\citep{bouniausing, arenas2022computing, bounia2024enhancing, bounia2023approximating}, monotonic models~\citep{marques2021explanations, harzli2022cardinality}, linear models~\citep{marques2020explaining, subercaseaux2025probabilistic}, and additive models~\citep{bassan2026provably}. While much of the focus in formal XAI is on the theoretical complexity of generating such explanations, there also exist a major focus of obtaining such explanations in practice via automated reasoning tools such as MaxSAT or MILP solvers for tree ensembles~\citep{audemard2023computing, audemard2022preferred, ignatiev2022using, boumazouza2021asteryx}, or neural network verification models~\citep{wang2021beta, wu2024marabou} for obtaining formal explanations on neural networks~\citep{wu2024verix, bassan2023towards, izza2024distance,hadad2026, ignatiev2019abduction}. Other frameworks aim to mitigate the high complexity of generating such explanations via various techniques such as abstractions~\citep{de2025faster, bassan2025explaining, boumazouza2026fame}, relaxations such as smoothing~\citep{xue2023stability, anani2025pixel, jin2025probabilistic}, or training-time interventions~\citep{alvarez2018towards, bassan2025explain, bassan2025self}

Finally, within formal XAI, a related line of work studies methods that leverage combinatorial and submodular optimization techniques to improve the computation of \emph{local} probabilistic sufficient reasons~\citep{bouniausing, bounia2023approximating}, particularly for decision tree classifiers. As noted in that literature, as well as in this work, the value-function objective underlying these local explanation variants is not submodular in practice. Nevertheless, greedy algorithms inspired by submodular optimization often serve as effective heuristics and can yield high-quality explanations empirically. In contrast, our work establishes a rigorous theoretical connection between combinatorial and submodular optimization objectives and the \emph{global} variants of sufficient and contrastive reasons. This connection provides formal guarantees and opens the door to a broader range of principled algorithmic implementations.

%Closer to our setting, several frameworks have extended formal guarantees to neural networks~\citep{Laagmirhpanikwma21, bassan2023formally, hadad2026, wu2024verix, labbaf2025space, de2025faster, fel2023don, soria2025formal}. To cope with the inherent computational hardness, these approaches often rely on abstractions~\citep{bassan2025explaining, boumazouza2026fame}, smoothing techniques~\citep{xue2023stability, jin2025probabilistic, anani2025pixel}, or self-explaining architectures~\citep{alvarez2018towards, bassan2025explain, bassan2025self}.

%A complementary line of work considers explanations with formal guarantees for linear models~\citep{marques2020explaining, subercaseaux2025probabilistic}, which constitute a restricted subclass of additive models, as well as the theoretical study of~\citep{bassan2025additive}, which analyzes the computational complexity of generating such explanations.

\section{Model and Distribution Formalizations}
\label{models_and_distributions_app}

In this section, we formalize the models and distributions referenced throughout the paper. Specifically, Subsection~\ref{model_formaliations_subsec} defines the model families, while Subsection~\ref{distribution_formaliations_subsec} formalizes the distributions.


\subsection{Model Formalizations}
\label{model_formaliations_subsec}

In this subsection, we formalize the three base-model types that were analyzed throughout the paper: \begin{inparaenum}[(i)] \item (axis-aligned) decision trees, \item linear classifiers, and \item neural networks with ReLU activations.
\end{inparaenum} %We also formalize \emph{monotonic classifiers}, which we analyze in section~\ref{monotonic_expansion_section}. 

\mysubsection{Decision Trees.} We define a decision tree (DT) as a directed acyclic graph that represents a function $f: \mathbb{F} \to [c]$, where $c \in \mathbb{N}$ denotes the number of classes. The graph encodes the function as follows: \begin{inparaenum}[(i)]
\item Each internal node $v$ is assigned a distinct binary input feature from the set $\{1,\ldots,n\}$;
\item Every internal node $v$ has at most $k$ outgoing edges, each corresponding to a value in $[k]$ assigned to $v$;
\item Along any path $\alpha$ in the decision tree, each variable appears at most once;
\item Each leaf node is labeled with a class from $[c]$.
\end{inparaenum} Thus, assigning values to the inputs $\x \in \mathbb{F}$ uniquely determines a path $\alpha$ from the root to a leaf in the DT, where the function output $f(\x)$ corresponds to a class label $i \in [c]$. The size of the DT, denoted $\vert f \vert$, is defined as the total number of edges in the graph. To allow for flexible modeling, the ordering of input variables $\{1, \ldots, n\}$ may differ across distinct paths $\alpha$ and $\alpha'$, ensuring that no variable is repeated along a single path. 


%Thus, assigning a value to the inputs $\x\in\mathcal{X}$ uniquely determines a specific path $\alpha$ from the root to a leaf in the DT. The function $f(\x)$ is assigned either some $i\in [c]$ or some $i\in \mathbb{R}$ (depending on classification or regression tasks). The size of the DT, denoted as $\vert f \vert$, is measured by the total number of edges in its graph. To allow flexibility of the modeling of the DT, we permit different varying orderings of the input variables $\{1,\ldots,n \}$ across any two distinct paths, $\alpha$ and $\alpha'$. This ensures that no two nodes along any single path $\alpha$ share the same label. We note that in some of the proofs provided in this work, we simplify them by assuming that $k:=2$ or in other words that each feature is defined over binary feature assignments rather than discrete ones. However, this assumption is only for the sake of simplifying the proofs, and our proofs hold also for a general $k$ as well.






%We follow standard conventions in the literature~\citep{BaMoPeSu20, arenas2022computing, bassanlocal} and focus our complexity analysis on \emph{Free Binary Decision Diagrams} (FBDDs)\citep{Le59}. In Section\ref{framework_extensions_section}, we further discuss how these results generalize to other forms of decision trees defined over various input domains. A \emph{Binary Decision Diagram} (BDD) is a directed acyclic graph that represents a Boolean function $f: \mathbb{F} \to \{0,1\}$. The graph encodes the Boolean function through the following structure: \begin{inparaenum}[(i)] 
%\item Each internal node is labeled with a binary input variable from the set $\{1,\ldots,n\}$; 
%\item Each internal node has exactly two outgoing edges, corresponding to the possible values $\{0,1\}$ assigned to its variable; 
%\item Each leaf node is labeled either \emph{true} or \emph{false}.
%\end{inparaenum}

 %We adhere to standard conventions in the field~\citep{BaMoPeSu20, arenas2022computing, bassanlocal} and concentrate our complexity representation on \emph{FBDDs} (Free Binary Decision Diagrams~\citep{Le59}). In section~\ref{framework_extensions_section}, we further explain how these findings can be extended to other types of decision trees with varying input domains. A \emph{binary decision diagram} (\emph{BDD}) is an acyclic-directed graph and serves as a graphical model for a Boolean function $f: \mathbb{F} \to \{0,1\}$. This graph embodies the given Boolean function in the following manner: 
%\begin{inparaenum}[(i)]
%	\item Each internal node $v$ is associated with a unique binary input feature from the set $\{1,\ldots,n\}$;
%	\item Every internal node $v$ has precisely two outgoing edges, corresponding to the values $\{0,1\}$ which are assigned to $v$;
%	\item In the BDD, each variable is encountered no more than once on any given path $\alpha$;
% 	\item Each leaf node is labeled either \emph{true} or \emph{false}.
%\end{inparaenum}

%Hence, assigning values to the input $\x \in \mathbb{F}$ uniquely determines a path $\alpha$ from the root to a leaf in the BDD. The function $f(\x)$ outputs $1$ if the leaf is labeled \emph{true}, and $0$ if it is labeled \emph{false}. The size of a BDD, denoted by $\vert f \vert$, is defined as the total number of edges in the graph. In this work, we focus on the commonly used \emph{Free Binary Decision Diagrams} (\emph{FBDDs}), which allow the input variables $\{1, \ldots, n\}$ to appear in different orders across different paths $\alpha$ and $\beta$. A key constraint of FBDDs is that no variable is tested more than once along any single path. Notably, a \emph{decision tree} is a special case of an FBDD where the graph structure forms a tree.


%Thus, assigning a value to the inputs $\x\in\mathbb{F}$ uniquely determines a specific path $\alpha$ from the root to a leaf in the BDD. The function $f(\x)$ is assigned a $1$ if the terminal node leaf is labeled \true{}, and $0$ if it is labeled \false{}. The size of the BDD, denoted as $\vert f \vert$, is measured by the total number of edges in its graph. Our study focuses on the widely used \emph{Free Binary Decision Diagrams} (\emph{FBDDs}), which permit different varying orderings of the input variables $\{1,\ldots,n \}$ across any two distinct paths, $\alpha$ and $\alpha'$. This ensures that no two nodes along any single path $\alpha$ share the same label. Essentially, a \emph{decision tree} is a type of FBDD with a fundamental graph structure that is a tree.

\textbf{Neural Networks.} We present our \emph{hardness} proofs for neural networks with ReLU activations, though our \emph{tractability results} (i.e., polynomial-time algorithms) apply to any architecture that allows for polynomial-time inference --- a standard assumption. Thus, all separation results between tractable and intractable cases we prove carry over to \emph{any neural architecture} at least as expressive as a standard feed-forward ReLU network, encompassing many widely used models such as ResNets, CNNs, Transformers, Diffusion models, and more. Moreover, note that any ReLU network can be represented as a fully connected network by assigning zero weights and biases to missing connections. Following standard conventions~\citep{BaMoPeSu20, bassanlocal, adolfi2024computational}, we thus assume the network is fully connected. Specifically, our analysis applies to multi-layer perceptrons (MLPs). Formally, an MLP $f$ consists of $t-1$ \emph{hidden layers} ($g^{(j)}$ for $j = 1,\ldots,t-1$) and one output layer $g^{(t)}$, where each layer is defined recursively as:


%We present our proofs for neural networks with ReLU activations, though our \emph{tractability} (i.e., polynomial-time algorithms) extend to any neural network architecture with polynomial-time inference --- a standard requirement. It is easy to verify that any ReLU network can be transformed into a fully connected one by assigning a weight and bias of $0$ to all non-existent connections. In line with standard practice~\citep{BaMoPeSu20, bassanlocal, adolfi2024computational}, we therefore assume the network is fully connected. That is, our analysis applies to multi-layer perceptrons (MLPs). Formally, an MLP $f$ comprises $t-1$ \emph{hidden layers} ($g^{(j)}$ for $j = 1,\ldots,t-1$) and a single output layer $g^{(t)}$. Each layer is defined recursively as:  

\begin{equation}
g^{(j)} := \sigma^{(j)}(g^{(j-1)} W^{(j)} + b^{(j)})
\end{equation}  

where $W^{(j)}$ denotes the weight matrix, $b^{(j)}$ the bias vector, and $\sigma^{(j)}$ the activation function of the $j$-th layer. Accordingly, the model comprises $t$ weight matrices ($W^{(1)}, \ldots, W^{(t)}$), $t$ bias vectors ($b^{(1)}, \ldots, b^{(t)}$), and $t$ activation functions ($\sigma^{(1)}, \ldots, \sigma^{(t)}$).

%where $W^{(j)}$ is the weight matrix, $b^{(j)}$ is the bias vector, and $\sigma^{(j)}$ is the activation function of the $j$-th layer. The model thus consists of $t$ weight matrices ($W^{(1)},\ldots,W^{(t)}$), $t$ bias vectors ($b^{(1)},\ldots,b^{(t)}$), and $t$ activation functions ($\sigma^{(1)},\ldots,\sigma^{(t)}$).



%We focus on neural networks with ReLU activations. It is straightforward to show that any neural network using ReLU activations can be straightforwardly transformed into a fully connected network by assigning ``missing'' connections a weight and bias of $0$ for any non-connected neurons. Following standard conventions~\citep{BaMoPeSu20, bassanlocal, adolfi2024computational}, we assume the network is fully connected. In other words, our analysis pertains to multi-layer perceptrons (MLPs). Formally, an MLP, denoted by $f$, consists of $t-1$ \emph{hidden layers} ($g^{j}$ where $j$ ranges from $1$ to $t-1$) and a single output layer ($g^{t}$). The layers are defined recursively --- each layer $g^{(j)}$ is computed by applying the activation function $\sigma^{(j)}$ to the linear combination of the outputs from the previous layer $g^{(j-1)}$, the corresponding weight matrix $W^{(j)}$, and the bias vector $b^{(j)}$. This is represented as $g^{(j)} := \sigma^{(j)}(g^{(j-1)}W^{(j)} + b^{(j)})$ for each $j$ in $\{1,\ldots,t\}$. The model includes $t$ weight matrices ($W^{(1)},\ldots,W^{(t)}$), $t$ bias vectors ($b^{(1)},\ldots,b^{(t)}$), and $t$ activation functions ($\sigma^{(1)},\ldots,\sigma^{(t)}$).


The input layer is defined as $g^{(0)} = \x \in \{0,1\}^n$, representing the binary input vector. The dimensions of the network are governed by a sequence of positive integers ${d_0, \ldots, d_t}$, with weight matrices and bias vectors given by $W^{(j)} \in \mathbb{Q}^{d_{j-1} \times d_j}$ and $b^{(j)} \in \mathbb{Q}^{d_j}$, respectively. These parameters are learned during training. Since the model functions as a binary classifier over $n$ features, we set $d_0 = n$ and $d_t = 1$. The hidden layers use the ReLU activation function, defined by $\text{ReLU}(x) = \max(0, x)$. Although a sigmoid activation is typically used during training, for interpretability purposes, we assume the output layer applies a threshold-based step function, defined as $\text{step}(\z) = 1$ if $\z \geq 0$ and $\text{step}(\z) = 0$ otherwise.

%In the specified MLP, the function $f$ is defined as the output of the final layer: $f \coloneqq g^{(t)}$. The input layer is denoted by $g^{(0)} = \x \in \{0,1\}^n$, representing the model's input vector. The dimensions of the weight matrices and bias vectors are determined by a sequence of positive integers $\{d_0, \ldots, d_t\}$. We focus on models where all weights and biases are rational numbers, with $W^{(j)} \in \mathbb{Q}^{d_{j-1} \times d_j}$ and $b^{(j)} \in \mathbb{Q}^{d_j}$, and these parameters are learned during training. Since the model serves as a binary classifier over ${1, \ldots, n}$, we set $d_0 = n$ and $d_t = 1$. The main activation function used in the hidden layers is the \emph{ReLU} function, defined as $\text{ReLU}(x) = \max(0, x)$. For the output layer, while a sigmoid function is typically employed in training for classification, we assume a threshold-based \emph{step function} for post-hoc interpretability, defined as $\text{step}(\z) = 1$ if $\z\geq 0$ and $\text{step}(\z) = 0$ otherwise. 


%In the described MLP, the function $f$ is defined to output $f \coloneqq g^{(t)}$. The initial input layer $g^{(0)}$ is denoted by $\x \in {\{0,1\}}^n$, which serves as the model's input. The dimensions of the biases and weight matrices are specified by the sequence of positive integers $\{d_{0}, \ldots, d_{t}\}$. We specifically focus on weights and biases that are rational numbers, represented as $W^{(j)}\in \mathbb{Q}^{d_{j-1}\times d_{j}}$ and $b^{(j)}\in \mathbb{Q}^{d_{j}}$, which are parameters that are optimized during training. Given that the model is a binary classifier for indices $\{1, \ldots, n\}$, it follows that $d_{0} = n$ and $d_{t} = 1$. The primary activation function $\sigma^{(i)}$ that we focus on is the \emph{ReLU} activation function, defined as $reLU(x) = \max(0, x)$, except for the output layer, where a \emph{sigmoid} function is typically used for the classification. Since our focus is only on the \emph{post-hoc} interpretation of the corresponding model, we will equivalently assume the existence of a step function for the final layer activation, where we denote $step(\z) = 1 \iff \z \geq 0$.


\mysubsection{Tree Ensembles.} Many popular ensemble methods exist, but since our goal is to provide post-hoc explanations, we focus on the \emph{inference} phase rather than the training process. Our analysis centers on ensemble families that rely on either \emph{majority voting} (e.g., Random Forests) or \emph{weighted voting} (e.g., XGBoost) during inference. However, as with our treatment of neural networks, our \emph{tractability} results --- namely, polynomial-time algorithms --- extend to any possible ensemble configuration with polynomial-time inference, encompassing an even broader range of ensemble configurations.
    
\mysubsection{Majority Voting Inference.} In \emph{majority voting} inference, the final prediction $f(\x)$ is assigned to the class $j \in [c]$ that receives the majority of votes among the individual predictions $f_i(\x)$ from all $i \in [k]$ (i.e., from each tree in the ensemble).

\begin{equation}
    f(\x):=\operatorname{majority}(\{f_i(\x) \ | \ i \in [k]\}) 
\end{equation}

where $\operatorname{majority}(S)$ denotes the most frequent label in the multiset $S$. If there is a tie, it is resolved by a fixed tie-breaking rule (e.g., lexicographic order or predefined priority).


%are at least $\ceil{\frac{k}{2}}$ base-models within an ensemble $f$ where $f_i(\x)=1$. Put simply, the decision for $f(\x)$ is determined by the majority consensus of the models involved in $f$. Formally, for any $\x\in\mathbb{F}$ we define $f$ as follows:

% \begin{equation}
%f(\x):=\begin{cases}
%    1 \quad if \ \ | \{ \ i \ | \ f_i(\x)=1 \ \}| \ \geq \ \ceil{\frac{k}{2}}\\
%    0 \quad otherwise
%\end{cases}
%\end{equation}


%We use the common notation of this form of voting as \emph{hard-voting} which is usually suitable for \emph{binary classification}, which is our concentraition here. For probablistic classification, another common approach is \emph{soft-voting} for which.



\mysubsection{Weighted Voting Inference.} In \emph{weighted voting} inference, each model in the ensemble is assigned a weight $\phi_i \in \mathbb{Q}$ representing its relative importance. The predicted class is the one with the highest total weight across all models. Formally, for any $\x \in \mathbb{F}$, we define $f$ as:


\begin{equation}
f(\x) := \arg\max_{j \in [c]} \sum_{i=1}^k \phi_i \cdot \mathbb{I}[f_i(\x) = j]
\end{equation}


where $\mathbb{I}$ denotes the identity function.


%\begin{equation}
%\label{weighted_voting_formalization}
%f(\x):=step(\sum_{1\leq i\leq n} \phi_i\cdot f_i(\x))
%\end{equation}

%where $step(\z) = 1 \iff \z \geq 0$.


\subsection{Distribution Formalizations}

\label{distribution_formaliations_subsec}

This subsection formalizes the distribution definitions discussed in the main paper.


\textbf{Empirical Distributions.} The distribution $\mathcal{D}$ over the input features will be defined based on a dataset $\mathbf{D}$ comprising various inputs $\z^1, \z^2,\ldots, \z^{|\mathbf{D}|}$. Here, the distribution of any given input $\x \in \mathbb{F}$ is defined by the frequency of occurrence of $\x$ within $\mathbf{D}$, specifically by:

\begin{equation}
    \textbf{Pr}(\x):=\frac{1}{|\mathbf{D}|}\sum_{i=1}^{|\mathbf{D}|}\mathbb{I}(\z^i=\x)
\end{equation}


\textbf{Independent Distributions.} Formally, given a probability value $p(\x_i) \in [0,1]$ defined for each individual input feature, we say that the distribution $\mathcal{D}$ is independent if the joint probability over inputs is given by $\textbf{Pr}(\x) := \prod_{i \in [n]} p(\x_i)$.

We observe that when $p(\x_i)=p(\x_j)$ holds for all $i,j\in [n]$, the distribution reduces to the \emph{uniform distribution}, a special case of independent distributions.

\textbf{General Distributions.} We note that many of the proofs in this work --- particularly those concerning fundamental properties of value functions --- apply broadly over general distribution assumptions. While empirical distributions (using the training dataset as a proxy) are common in XAI, alternative frameworks for approximating distributions include k-NN resampling from nearby points~\citep{li2023k, almanjahie2018knn, gweon2019k}, copulas for modeling tabular dependencies~\citep{grosser2022copulae}, or more advanced conditional generative models such as CTGAN~\citep{xu2019modeling} and conditional diffusion models~\citep{huang2022fastdiff}. Such choices are common in the counterfactual explanation and algorithmic recourse literature~\citep{karimi2021algorithmic, fokkema2024risks, fokkema2023attribution, verma2024counterfactual}, which are related to contrastive explanations. Furthermore, when the value function is not computed directly from a structural property of the model (e.g., leaf enumeration in decision trees) or from empirical distributions, one may instead approximate it using methods such as Monte Carlo sampling~\citep{hastings1970monte}, following approaches similar to~\citep{subercaseaux2024probabilistic, lopardo2023sea}.

\subsection{Subset vs. Cardinal Minimality}
\label{subset_cardinal_minimality}

In this subsection, we provide a more detailed discussion of the distinction between cardinal and subset minimality. Cardinal minimality offers a substantially stronger guarantee, as it corresponds to a globally minimal explanation size, whereas subset minimality only ensures a local form of minimality.

To see why cardinal minimality is stronger, note that if $S \subseteq [n]$ is a cardinally minimal explanation, then no smaller set $S'$ with $|S'|<|S|$ can qualify as an explanation. Hence, no strict subset $S' \subsetneq S$ is an explanation, which means $S$ is also subset minimal.

However, the reverse does not hold. Consider the function:

\begin{equation}
    f:=\x_1 \vee (\x_2 \wedge \x_3)
\end{equation}


and the assignment $\x := (1,1,1)$, which gives $f(1,1,1)=1$. Fixing feature $\x_1=1$ yields a cardinally minimal (and subset minimal) sufficient reason, since the prediction remains $1$ regardless of $\x_2,\x_3$. But fixing both $\x_2=1$ and $\x_3=1$ also gives a subset minimal explanation—yet not a cardinally minimal one. Thus, while every cardinally minimal explanation is subset minimal, not every subset minimal explanation is cardinally minimal.

%In this subsection, we a more thorough discussion of the difference between cardinal and subset minimality. We note that cardinal minimality is a substantially stronger guarantee as it correlates with a globally minimal explanation size, as opposed to the subset-minimality definition which provides a form of local-minimality.


%First, we explain why cardinal minimality is stronger. By definition, if a subset $S\subseteq [n]$ is a cardinally minimal explanation, then no smaller subset $S'$ with $|S'|<|S|$ can also be an explanation. Consequently, any strict subset $S'\subsetneq S$ cannot be an explanation either, implying that $S'$ is also subset minimal.

%However, \emph{the reverse implication does not hold}. For example, consider the function:

%\begin{equation}
%    f:=\x_1 \vee (\x_2 \wedge \x_3)
%\end{equation}

%and let us assume some assignment $\x:=(1,1,1)$ that implies $f(1,1,1)=1$. Here, fixing feature $1$ (i.e., setting $\x_1=1$) yields a cardinally minimal (and also subset minimal) local sufficient reason since for any assignment over $\x_2$ or $\x_3$ the prediction of $f$ would remain $1$. However, fixing both features $2$ and $3$, i.e., $\x_2=1$ and $\x_3=1$  
%also forms a subset minimal explanation, though it is not cardinally minimal. In other words not any subset minimal explanation is a cardinally minimal one.

%values $[k]$, we can describe a probability function $p:[n]\times [k]\to [0,1]$. For example, $p(1,2)=\frac{1}{2}$, implies that the probability of feature $i=1$, to be set to the value $k=2$ is $\frac{1}{2}$. Then we can define $\mathcal{D}_p$ as an independent distribution over $\mathcal{X}$ iff:

%\begin{equation}
%    \label{eq:explanation}
%    \mathcal{D}_p(\x):=\Big({\displaystyle \prod_{i\in [n], \x_i=j} p(i, j)}\Big)
%\end{equation}

%It is evident that the uniform distribution is a specific instance of $\mathcal{D}_p$, obtained by setting $p(i,j):=\frac{1}{|k|}$ for every $i \in [n]$ and $j \in [k]$.



\section{Proof of Proposition~\ref{nonuniqueness_prop}}

\label{non_uniqueness_appendix}

\nonuniqueness*
%
\begin{proof}
% Assume an arbitrary (non-constant) function $h_{1}(X_{1})$ defined on variables $X_{1} = \{x_1, \dots, x_n\}$, where $X_{1}$ follows a distribution $\fml{D}_{1}$.
% Let $K = \{c_1, \dots, c_m\}$ denote the set of output classes of $h_{1}(X_{1})$.
% Moreover, assume that $S'$ is the unique global subset-minimal $\delta_{1}$-sufficient reason of $h_{1}$.
% We define
% \begin{equation}
% f(X_{1}, X_{2}) :=
% \left \{
% \begin{array}{lcl}
% h_{1}(X_{1}) & \quad & \textnormal{if~$x_{n+1}=0$} \\[3pt]
% c_{m+1} & \quad & \textnormal{if~$x_{n+1}=1$}
% \end{array}
% \right.
% \end{equation}
% %
% where $c_{m+1} \not\in K$ and $x_{n+1} \not\in X_{1}$.
% Assume $x_{n+1}$ follows a uniform distribution $\fml{D}_{2}$,
% and two distributions $\fml{D}_{1}$ and $\fml{D}_{2}$ are independent.
% Let $\fml{D}$ be the new distribution defined on $X_{1}\cup \{x_{n+1}\}$.
% %
% For an arbitrary $S$, if $S = \emptyset$, we have:
% \begin{align*}
% v^{g}_{\text{suff}}(\emptyset)
% &= \sum_{\substack{\x\sim \fml{D} \\ x_{n+1} = 1}}
% \mbf{Pr}(\x) \cdot \mbf{Pr}_{\z\sim \fml{D}}(f(\z) = c_{m+1}) \ + \\
% &\hspace{.7em} \sum_{\substack{\x\sim \fml{D} \\ x_{n+1} = 0}}
% \mbf{Pr}(\x) \cdot \mbf{Pr}_{\z\sim \fml{D}}(f(\z) = h_{1}(\x)).
% \end{align*}
% %
% It is worth noting that $\mbf{Pr}_{\z\sim \fml{D}}(f(\z) = c_{m+1}) = 0$ when $z_{n+1} = 0$,
% and $\mbf{Pr}_{\z\sim \fml{D}}(f(\z) = c_{m+1}) = 1$ when $z_{n+1} = 1$.
% Moreover, $\mbf{Pr}_{\z\sim \fml{D}}(f(\z) = h_{1}(\x)) = \mbf{Pr}_{\z\sim \fml{D}}(h_{1}(\z) = h_{1}(\x))$ when $z_{n+1} = 0$,
% and $\mbf{Pr}_{\z\sim \fml{D}}(f(\z) = h_{1}(\x)) = 0$ when $z_{n+1} = 1$.
% While if $S = \{x_{n+1}\}$, we have:
% \begin{align*}
% v^{g}_{\text{suff}}(\{x_{n+1}\})
% &= \sum_{\substack{\x\sim \fml{D} \\ x_{n+1} = 1}}
% \mbf{Pr}(\x) \cdot \mbf{Pr}_{\z\sim \fml{D}}(c_{m+1} = c_{m+1} \mid z_{n+1} = x_{n+1}) \ + \\
% &\hspace{.7em} \sum_{\substack{\x\sim \fml{D} \\ x_{n+1} = 0}}
% \mbf{Pr}(\x) \cdot \mbf{Pr}_{\z\sim \fml{D}}(h_{1}(\z) = h_{1}(\x) \mid z_{n+1} = x_{n+1}).
% \end{align*}
% %
% Clearly, $v^{g}_{\text{suff}}(\{x_{n+1}\}) > v^{g}_{\text{suff}}(\emptyset)$.
% If $S = S'$, we have:
% \begin{align*}
% v^{g}_{\text{suff}}(S')
% &= \sum_{\substack{\x\sim \fml{D} \\ x_{n+1} = 1}}
% \mbf{Pr}(\x) \cdot \mbf{Pr}_{\z\sim \fml{D}}(f(\z) = c_{m+1} \mid \z_{S'} = \x_{S'})) \ + \\
% &\hspace{.7em} \sum_{\substack{\x\sim \fml{D} \\ x_{n+1} = 0}}
% \mbf{Pr}(\x) \cdot \mbf{Pr}_{\z\sim \fml{D}}(f(\z) = h_{1}(\x) \mid \z_{S'} = \x_{S'})).
% \end{align*}
% %
% The analysis of $\mbf{Pr}_{\z\sim \fml{D}}(f(\z) = c_{m+1} \mid \z_{S'} = \x_{S'}))$ is identical to that of $\mbf{Pr}_{\z\sim \fml{D}}(f(\z) = c_{m+1})$,
% while the analysis of $\mbf{Pr}_{\z\sim \fml{D}}(f(\z) = h_{1}(\x) \mid \z_{S'} = \x_{S'}))$ is identical to that of $\mbf{Pr}_{\z\sim \fml{D}}(f(\z) = h_{1}(\x))$.
% Since $S'$ is a global subset-minimal $\delta_{1}$-sufficient reason of $h_{1}$,
% we can infer that $v^{g}_{\text{suff}}(S') > v^{g}_{\text{suff}}(\emptyset)$.
% Define $\delta = \epsilon + v^{g}_{\text{suff}}(\emptyset)$ such that 
% \[
% \delta \le v^{g}_{\text{suff}}(\{x_{n+1}\}), ~~\text{and}~~ \delta \le v^{g}_{\text{suff}}(S'),
% \]
% %
% this implies that we have at least two global subset-minimal $\delta$-sufficient reasons for the function $f$,
% one is $\{x_{n+1}\}$, and the other is a subset $S'' \subseteq S'$.


It is known that certain Boolean functions admit an exponential number of subset-minimal local (non-probabilistic) sufficient reasons for some input $\x$~\cite{bassanlocal}. In particular, this phenomenon was shown to occur in functions of the following form. We demonstrate that the same function admits an exponential number of \emph{global} and \emph{probabilistic} sufficient reasons. Notably, the mentioned function is a special case of a broader class of \emph{threshold Boolean functions} described in~\cite{wegener2005complexity}. Specifically, for $n = 2k + 1$ with $k \in \mathbb{N}$, this function is defined as:


\begin{equation}
f(\x) :=
\left \{
\begin{array}{lcl}
1 & \quad & \textnormal{if \ \ $\sum_{i=1}^{n}\x_{i} \ge k+1$} \\[3pt]
0 & \quad & \textnormal{otherwise}
\end{array}
\right.
\end{equation}
%


where $\x$ is drawn from a uniform distribution $\mathcal{D}$.
Notably, the function $f$ is \emph{symmetric}, in the sense of symmetric threshold Boolean functions~\cite{wegener2005complexity}: its output depends solely on the number of 1's (or, equivalently, 0's) in the input and is invariant under any permutation of input bits.
Furthermore, the two extreme inputs --- the all-zeros vector $(0,\dots,0)$ and the all-ones vector $(1,\dots,1)$ --- each admit an exponential number of subset-minimal local (non-probabilistic) sufficient reasons.

%where $\x$ is sampled from a uniform distribution $\mathcal{D}_p$.
%It is worth noting that this function $f$ is \emph{symmetric}, as defined by the symmetric threshold Boolean functions in~\cite{wegener2005complexity}, that is,
%its output depends only on the count of 1's (equivalently, on the count of 0's) in the input, 
%and remains unchanged under any permutation of the input bits.
%Moreover, the two extreme points---the all-zeros vector $(0,\dots,0)$ and the all-ones vector $(1,\dots,1)$
%---have the same exponential number of subset-minimal local (non-probabilistic) sufficient reasons.


Since the function $f$ is symmetric, the condition
$v^{g}_{\text{suff}}(S_{1}) = v^{g}_{\text{suff}}(S_{2})$
holds for any two subsets $S_1, S_2 \subseteq [n]$.
Moreover, it can be verified that for any subset $S \subseteq [n]$ of size $k+1$, we have:

%Because the function $f$ is symmetric, we have that the condition
%$v^{g}_{\text{suff}}(S_{1}) = v^{g}_{\text{suff}}(S_{2})$
%for any two $S_{1}\subseteq X$ and $S_{2}\subseteq X$.
%
%Moreover, it can be verified that for any set $S\subseteq X$ with size $k+1$, we have
\begin{equation}
    (\forall i\in S). \quad v^{g}_{\text{suff}}(S) > v^{g}_{\text{suff}}(S\setminus\{i\})
\end{equation}

%$(\forall i\in S). v^{g}_{\text{suff}}(S) > v^{g}_{\text{suff}}(S\setminus\{i\})$.


We also deliberately set $\delta$ to satisfy:

    
\begin{equation}
v^{g}_{\text{suff}}(S\setminus\{i\}) < \delta \le v^{g}_{\text{suff}}(S),
\end{equation}

and also set $k:=\floor{\frac{n}{2}}$. Each of these subsets is \emph{subset-minimal} because any subset of size at most $\floor{\frac{n}{2}} - 1$ fails to be a sufficient reason for $\langle f, \x \rangle$ (with respect to the $\delta$ threshold). Therefore, we can directly apply the same analysis as in~\citep{bassanlocal}. In particular, there are exactly $\binom{n}{\floor{\frac{n}{2}}}$ subset-minimal local sufficient reasons for $\langle f, \x \rangle$. Using Stirling's approximation, we obtain:

\begin{equation}
    \lim_{n\to \infty} \ \ \frac{2\sqrt{2\pi}}{e^2}\cdot \frac{2^n}{\sqrt{n}} \leq \binom{n}{\floor{\frac{n}{2}}}\leq \frac{e}{\pi}\cdot \frac{2^n}{\sqrt{n}}
\end{equation}

This yields the corresponding bound on the number of subset-minimal global $\delta$-sufficient reasons.

\end{proof}



\section{Proof of Proposition~\ref{globalmonotonicity_prop}}
\label{globalmonotonicity_prop_appendix_sec}

\globalmonotonicity*

%\begin{restatable}{proposition}{globalmonotonicity}\label{globalmonotonicity_prop}
%While the local probabilistic setting (for any $\delta$) lacks monotonicity --- i.e., the value functions $v^{\ell}_{\text{con}}$ and $v^{\ell}_{\text{suff}}$ are non-monotonic~\citep{arenas2022computing, izza2023computing, subercaseaux2025probabilistic, izza2024locally} --- in the global probabilistic setting (also for any $\delta$), both value functions $v^g_{\text{con}}$ and $v^g_{\text{suff}}$ are monotonic non-decreasing.
%\end{restatable}

\label{monotonicity_proofs_appendix}

\emph{Proof.} In this section, we prove the monotonicity property for both types of global explanations --- sufficient and contrastive. We begin by establishing the property for global sufficiency. To build intuition, we first focus on the simpler case of Boolean functions, then generalize the result to functions with discrete multi-valued input domains and multiple output classes.
We further extend the result to continuous input domains and show that the monotonicity property holds for any well-defined classification function.
Finally, we demonstrate how the same monotonicity guarantees apply to global contrastive explanations. Importantly, none of our proofs rely on the assumption of feature independence --- the monotonicity property holds strongly for \emph{any} underlying distribution.

Lastly, we present a simple example showing that the monotonicity property does not hold for the local probabilistic sufficient or contrastive value functions.

%In this section we will prove the monotonicity property for both global settings - both sufficient and contrastive. We will start by proving this for the sufficient definition. To build better intuition for the correctness, we start in Subsection~\ref{ssec:bfd} by proving monotonicity for a simple case encompassing boolean functions. We then later in Subsection~\ref{ssec:mfd} extend this to multi-value outputs under \emph{discrete} input domains. In Subsection~\ref{} we extend this to \emph{continous} inputs, and in Subsection~\ref{} show how all of these monotonicity results extend from sufficient explanations to contrastive ones. We note that we \emph{do not} assume feature independence in any of these proofs and that the monotonicity property holds generally in a very strong sense for \emph{any} possible distribution.

\begin{lemma} \label{lem:bfd}
The global sufficient value function $v^g_{\text{suff}}$
is monotonic non-decreasing for Boolean functions,
under any data distribution.
\end{lemma}


\begin{proof}
Given an arbitrary set $S\in [n]$, and a feature $i\not\in S$, we have that:

\begin{equation}
\begin{aligned}
\mbf{Pr}(\x) = \mbf{Pr}(\x_{S}, \x_{\bar{S}}) = \mbf{Pr}(\x_{\bar{S}} \mid \x_{S}) \cdot \mbf{Pr}(\x_{S})
\end{aligned}
\end{equation}


Consequently, this implies that:

\begin{equation}
\begin{aligned}
\mbf{Pr}(\x)
&= \mbf{Pr}(\x_{\overline{S\cup\{i\}}} \mid \x_{S\cup\{i\}}) \cdot \mbf{Pr}(\x_{S\cup\{i\}}) \\
&= \mbf{Pr}(\x_{\overline{S\cup\{i\}}} \mid \x_{S\cup\{i\}}) \cdot \mbf{Pr}(\x_{S},x_i) \\
&= \mbf{Pr}(\x_{\overline{S\cup\{i\}}} \mid \x_{S\cup\{i\}}) \cdot \mbf{Pr}(x_i \mid \x_{S}) \cdot \mbf{Pr}(\x_{S})
\end{aligned}
\end{equation}


Moreover, for any two points $\x, \x' \sim \mathcal{D}$ such that $\x_{S} = \x'_{S}$ and $f(\x) = 1 - f(\x')$, it holds that:

%~\footnote{
%For simplicity, we denote the data distribution by $\fml{D}$ instead of $\fml{D}_p$.
%}

\begin{equation}
\mbf{Pr}_{\z\sim \mathcal{D}}(f(\z) = f(\x) \mid \z_{S} = \x_{S}) = 1 - \mbf{Pr}_{\z\sim \mathcal{D}}(f(\z) = f(\x') \mid \z_{S} = \x'_{S})
\end{equation}

To simplify notation, we occasionally omit the distribution notation $\sim \mathcal{D}$ in the local probability expression $\mbf{Pr}_{\z\sim \mathcal{D}}(f(\z) = f(\x) \mid \z_{S} = \x_{S})$, and use $f^{+}$ and $f^{-}$ to denote the events $f(\z) = 1$ and $f(\z) = 0$, respectively.

We begin with a technical simplification of the probability term to facilitate the proof, and then proceed to establish monotonicity.

%We will begin with a techhnical simplification of the probability term that will fascialitate our proof and will then move to proving monotonicity.

\paragraph{Simplifying the probability term.}


Let $\mathcal{D}_{S}$ denote the distribution restricted to the set $S$.
Given a point $\x \sim \mathcal{D}$ and a set $S$, consider all points $\z \sim \mathcal{D}$ such that $\z_{S} = \x_{S}$ and $f(\z) = b$, where $b \in \{0,1\}$ —
they all share the same local probability $\mbf{Pr}_{\z \sim \mathcal{D}}(f(\z) = b \mid \z_{S} = \x_{S})$.
%
For every $\x \sim \mathcal{D}$ that shares the same $x_{S}$ and the same output $b$, we marginalize over $x_{\bar{S}}$.
This yields $\mbf{Pr}(f(\x) = b, \x_{S}) = \mbf{Pr}(f(\x) = b \mid \x_{S}) \cdot \mbf{Pr}(\x_{S})$,
from which we can infer:



%Let $\fml{D}_{S}$ denote the distribution restricted by the set $S$.
%Given a $\x \sim \mathcal{D}_p$, a set $S$, consider all points $\z \sim \mathcal{D}_p$ such that $\z_{S} = \x_{S}$ and $f(\z) = b$, where $b\in\{0,1\}$,
%they have the same local probability $\mbf{Pr}_{\z\sim \fml{D}}(f(\z) = b \mid \z_{S} = \x_{S})$.
%
%For all $\x \sim \fml{D}$ sharing the same $x_{S}$ and the same output $b$, we marginalize over $x_{\bar{S}}$.
%This gives us $\mbf{Pr}(f(\x) = b, \x_{S}) = \mbf{Pr}(f(\x) = b \mid \x_{S}) \cdot \mbf{Pr}(\x_{S})$,
%and we can infer that: 


\begin{equation}
\mbf{Pr}_{\z\sim \mathcal{D}}(f(\z) = b \mid \z_{S} = \x_{S}) = \mbf{Pr}(f(\x) = b \mid \x_{S}).
\end{equation}


Thus, for $\mbb{E}_{\x\sim \mathcal{D}}[\mbf{Pr}_{\z\sim \mathcal{D}}(f(\z) = f(\x) \mid \z_{S} = \x_{S})]$, we get that:

\begin{equation} \label{eq:gbx01}
\begin{aligned}
&\mbb{E}_{\x\sim \mathcal{D}}[\mbf{Pr}_{\z\sim \mathcal{D}}(f(\z) = f(\x) \mid \z_{S} = \x_{S})] \\
&= \hspace*{-1em} \sum_{\x_{S}\sim \fml{D}_{S}, f^{+}} \mbf{Pr}(f^{+}, \x_{S}) \cdot \mbf{Pr}_{\z\sim \mathcal{D}}(f^{+} \mid \z_{S} = \x_{S}) + \hspace*{-1.3em}
\sum_{\x_{S}\sim \fml{D}_{S}, f^{-}} \mbf{Pr}(f^{-}, \x_{S}) \cdot \mbf{Pr}_{\z\sim \mathcal{D}}(f^{-} \mid \z_{S} = \x_{S}) \\
&= \sum_{\x_{S}\sim \fml{D}_{S}} \mbf{Pr}(\x_{S}) \cdot \bigl[\bigl(\mbf{Pr}_{\z\sim \mathcal{D}}(f^{+} \mid \z_{S} = \x_{S})\bigr)^2 + \bigl(\mbf{Pr}_{\z\sim \mathcal{D}}(f^{-} \mid \z_{S} = \x_{S})\bigr)^2\bigr] \\
&= \sum_{\x_{S}\sim \fml{D}_{S}} \mbf{Pr}(\x_{S}) \cdot \bigl[\bigl(\mbf{Pr}_{\z\sim \mathcal{D}}(f^{+} \mid \z_{S} = \x_{S})\bigr)^2 + \bigl(1 - \mbf{Pr}_{\z\sim \mathcal{D}}(f^{+} \mid \z_{S} = \x_{S})\bigr)^2\bigr],
\end{aligned}
\end{equation}


From which it follows that:

\begin{equation} \label{eq:gbx02}
\begin{aligned}
&\mbb{E}_{\x\sim \mathcal{D}}[\mbf{Pr}_{\z\sim \mathcal{D}}(f(\z) = f(\x) \mid \z_{S\cup\{t\}} = \x_{S\cup\{t\}})] \\
&= \hspace*{-1em} \sum_{\x_{S\cup\{t\}}\sim \fml{D}_{S\cup\{t\}}} \hspace*{-1.3em}
\mbf{Pr}(\x_{S\cup\{t\}}) \cdot \bigl[\bigl(\mbf{Pr}_{\z}(f^{+} \mid \z_{S\cup\{t\}} = \x_{S\cup\{t\}})\bigr)^2 + \bigl(1 - \mbf{Pr}_{\z}(f^{+} \mid \z_{S\cup\{t\}} = \x_{S\cup\{t\}})\bigr)^2\bigr] \\
&= \hspace*{-.5em} \sum_{\substack{\x_{S}\sim \fml{D}_{S} \\ x_t = 1}}
\mbf{Pr}(x_t = 1, \x_{S}) \cdot \bigl[\bigl(\mbf{Pr}_{\z}(f^{+} \mid \z_{S} = \x_{S}, z_t = 1)\bigr)^2 + \bigl(1 - \mbf{Pr}_{\z}(f^{+} \mid \z_{S} = \x_{S}, z_t = 1)\bigr)^2\bigr] \ + \\
&\hspace{.7em} \sum_{\substack{\x_{S}\sim \fml{D}_{S} \\ x_t = 0}}
\mbf{Pr}(x_t = 0, \x_{S}) \cdot \bigl[\bigl(\mbf{Pr}_{\z}(f^{+} \mid \z_{S} = \x_{S}, z_t = 0)\bigr)^2 + \bigl(1 - \mbf{Pr}_{\z}(f^{+} \mid \z_{S} = \x_{S}, z_t = 0)\bigr)^2\bigr]
% &= \sum_{\substack{\x_{S}\sim \fml{D}_{S} \\ x_t = 1}}
% \mbf{Pr}(\x_{S}) \cdot \mbf{Pr}(z_t = 1 \mid \x_{S}) \cdot \bigl[\bigl(\mbf{Pr}_{\z}(f^{+} \mid \z_{S} = \x_{S}, z_t = 1)\bigr)^2 + \bigl(1 - \mbf{Pr}_{\z}(f^{+} \mid \z_{S} = \x_{S}, z_t = 1)\bigr)^2\bigr] \ + \\
% &\hspace{.7em} \sum_{\substack{\x_{S}\sim \fml{D}_{S} \\ x_t = 0}}
% \mbf{Pr}(\x_{S}) \cdot \mbf{Pr}(z_t = 0 \mid \x_{S}) \cdot \bigl[\bigl(\mbf{Pr}_{\z}(f^{+} \mid \z_{S} = \x_{S}, z_t = 0)\bigr)^2 + \bigl(1 - \mbf{Pr}_{\z}(f^{+} \mid \z_{S} = \x_{S}, z_t = 0)\bigr)^2\bigr].
\end{aligned}
\end{equation}

\paragraph{Proving monotonicity.}
%
We now proceed to prove the monotonicity claim, building on the previous simplification. We begin by introducing a few additional notations. Let $g_{1}^{+} := \mbf{Pr}_{\z \sim \mathcal{D}}(f(\z) = 1 \mid \z_{S} = \x_{S}, z_t = 1)$ and $g_{0}^{+} := \mbf{Pr}_{\z \sim \mathcal{D}}(f(\z) = 1 \mid \z_{S} = \x_{S}, z_t = 0)$. Similarly, define $P_{1} := \mbf{Pr}(z_t = 1 \mid \x_{S})$ and $P_{0} := \mbf{Pr}(z_t = 0 \mid \x_{S})$.
Then, the following expectations $\mbb{E}_{\x\sim \mathcal{D}}[\mbf{Pr}_{\z\sim \mathcal{D}}(f(\z) = f(\x) \mid \z_{S} = \x_{S})]$ 
and $\mbb{E}_{\x\sim \mathcal{D}}[\mbf{Pr}_{\z\sim \mathcal{D}}(f(\z) = f(\x) \mid \z_{S\cup\{t\}} = \x_{S\cup\{t\}})]$ can be simplified accordingly~\footnote{
Note that the notations $P_{1}$, $P_{0}$, $g_{1}^{+}$, and $g_{0}^{+}$—as well as all others used throughout the paper—are not fixed constants, but rather depend on the specific input $\x_{S}$. Formally, they should be written as $P_{1}(\x_{S})$, $P_{0}(\x_{S})$, and so on. However, for readability, we omit the explicit dependence on $\x_{S}$.}:
%It is important to note that the notations $P_{1}$, $P_{0}$, $g_{1}^{+}$, and $g_{0}^{+}$,
%as well as all other notations used in the remainder of the paper, are not constants but are associated with a specific input $\x_{S}$.
%More formally, these should be written as $P_{1}(\x_{S})$, $P_{0}(\x_{S})$, etc. 
%For simplicity of presentation, we omit the explicit dependence on $\x_{S}$.



%Let $g_{1}^{+}$ denote $\mbf{Pr}_{\z \sim \fml{D}}(f(\z) = 1 \mid \z_{S} = \x_{S}, z_t = 1)$,
%and $g_{0}^{+}$ denote $\mbf{Pr}_{\z \sim \fml{D}}(f(\z) = 1 \mid \z_{S} = \x_{S}, z_t = 0)$.
%Moreover, let $P_{1}$ denote $\mbf{Pr}(z_t = 1 \mid \x_{S})$, $P_{0}$ denote $\mbf{Pr}(z_t = 0 \mid \x_{S})$.
%$\mbb{E}_{\x\sim \fml{D}}[\mbf{Pr}_{\z\sim \fml{D}}(f(\z) = f(\x) \mid \z_{S} = \x_{S})]$ 
%and $\mbb{E}_{\x\sim \fml{D}}[\mbf{Pr}_{\z\sim \fml{D}}(f(\z) = f(\x) \mid \z_{S\cup\{t\}} = \x_{S\cup\{t\}})]$ can be simplified as follows



\begin{equation}
\begin{aligned}
\mbb{E}_{\x\sim \mathcal{D}}&[\mbf{Pr}_{\z\sim \mathcal{D}}(f(\z) = f(\x) \mid \z_{S} = \x_{S})] \\
&= \sum_{\x_{S} \sim \fml{D}_{S}} \mbf{Pr}(\x_{S}) \cdot [(P_{1} \cdot g_{1}^{+} + P_{0} \cdot g_{0}^{+})^2 + (1 - (P_{1} \cdot g_{1}^{+} + P_{0} \cdot g_{0}^{+}))^2] \\
&= \sum_{\x_{S} \sim \fml{D}_{S}} \mbf{Pr}(\x_{S}) \cdot [2(P_{1} \cdot g_{1}^{+} + P_{0} \cdot g_{0}^{+})^2 - 2(P_{1} \cdot g_{1}^{+} + P_{0} \cdot g_{0}^{+}) + 1].
\end{aligned}
\end{equation}

As a result, we obtain the following:


\begin{equation}
\begin{aligned}
&\mbb{E}_{\x\sim \mathcal{D}}[\mbf{Pr}_{\z\sim \mathcal{D}}(f(\z) = f(\x) \mid \z_{S\cup\{t\}} = \x_{S\cup\{t\}})] \\
&= \sum_{\substack{\x_{S} \sim \fml{D}_{S} \\ x_t = 1}} \mbf{Pr}(\x_{S}) \cdot P_{1} \cdot ((g_{1}^{+})^2 + (1-g_{1}^{+})^2)
+ \sum_{\substack{\x_{S} \sim \fml{D}_{S} \\ x_t = 0}} \mbf{Pr}(\x_{S}) \cdot P_{0} \cdot ((g_{0}^{+})^2 + (1 - g_{0}^{+})^2) \\
&= \sum_{\substack{\x_{S} \sim \fml{D}_{S} \\ x_t = 1}} \mbf{Pr}(\x_{S}) \cdot P_{1} \cdot (2(g_{1}^{+})^2 - 2g_{1}^{+} + 1)
+ \sum_{\substack{\x_{S} \sim \fml{D}_{S} \\ x_t = 0}} \mbf{Pr}(\x_{S}) \cdot P_{0} \cdot (2(g_{0}^{+})^2 - 2g_{0}^{+} + 1) \\
&= \sum_{\x_{S} \sim \fml{D}_{S}} \mbf{Pr}(\x_{S}) \cdot [P_{1} \cdot (2(g_{1}^{+})^2 - 2g_{1}^{+} + 1) + P_{0} \cdot (2(g_{0}^{+})^2 - 2g_{0}^{+} + 1)] \\
&= \sum_{\x_{S} \sim \fml{D}_{S}} \mbf{Pr}(\x_{S}) \cdot [2P_{1} \cdot ((g_{1}^{+})^2 - g_{1}^{+}) + 2P_{0} \cdot ((g_{0}^{+})^2 - g_{0}^{+}) + 1] \\
&= \sum_{\x_{S} \sim \fml{D}_{S}} \mbf{Pr}(\x_{S}) \cdot [2(P_{1} \cdot (g_{1}^{+})^2 + P_{0} \cdot (g_{0}^{+})^2) - 2(P_{1} \cdot g_{1}^{+} + P_{0} \cdot g_{0}^{+}) + 1].
\end{aligned}
\end{equation}


Let $\Delta_{t}(S,f)$ denote
$\mbb{E}_{\x\sim \mathcal{D}}[\mbf{Pr}_{\z\sim \mathcal{D}}(f(\z) = f(\x) \mid \z_{S\cup\{t\}} = \x_{S\cup\{t\}})] - \mbb{E}_{\x\sim \mathcal{D}}[\mbf{Pr}_{\z\sim \mathcal{D}}(f(\z) = f(\x) \mid \z_{S} = \x_{S})]$.
Without loss of generality, we assume that both $P_{1}$ and $P_{0}$ are positive.
Otherwise, it can be verified that $\Delta_{t}(S,f) = 0$.
Then, we get:


\begin{equation} \label{eq:gbx03}
\begin{aligned}
&\Delta_{t}(S,f) \\
&= \sum_{\x_{S} \sim \fml{D}_{S}} \mbf{Pr}(\x_{S} ) \cdot [2(P_{1} \cdot (g_{1}^{+})^2 + P_{0} \cdot (g_{0}^{+})^2) - 2(P_{1} \cdot g_{1}^{+} + P_{0} \cdot g_{0}^{+})^2] \\
% &= \sum_{\x_{S} \sim \fml{D}_{S}} 2 \cdot \mbf{Pr}(\x_{S}) \cdot [(P_{1} \cdot (g_{1}^{+})^2 + P_{0} \cdot (g_{0}^{+})^2) - (P_{1}^2 \cdot (g_{1}^{+})^2 + P_{0}^2 \cdot (g_{0}^{+})^2 + 2 \cdot P_{1} \cdot P_{0} \cdot g_{1}^{+} \cdot g_{0}^{+})] \\
&= \sum_{\x_{S} \sim \fml{D}_{S}} 2 \cdot \mbf{Pr}(\x_{S}) \cdot [(P_{1} - P_{1}^2) (g_{1}^{+})^2 + (P_{0} - P_{0}^2) (g_{0}^{+})^2 - 2 \cdot P_{1} \cdot P_{0} \cdot g_{1}^{+} \cdot g_{0}^{+}] \\
&= \sum_{\x_{S} \sim \fml{D}_{S}} 2 \cdot \mbf{Pr}(\x_{S}) \cdot P_{1} \cdot P_{0} \cdot (g_{1}^{+} - g_{0}^{+})^2 \\
&\ge 0.
\end{aligned}
\end{equation}


Note that $P_{1} - P_{1}^2 = P_{1}(1-P_{1}) = P_{1}P_{0}$ and similarly: $P_{0} - P_{0}^2 = P_{0}(1-P_{0}) = P_{0}P_{1}$).
Using these identities along with the previous result, we conclude that $\mbb{E}_{\x\sim \mathcal{D}}[\mbf{Pr}_{\z\sim \mathcal{D}}(f(\z) = f(\x) \mid \z_{S} = \x_{S})]$ is monotone, thereby concluding our proof.

\end{proof}


\begin{lemma} \label{lem:mfd}
The global sufficient value function $v^g_{\text{suff}}$
is monotonic non-decreasing for functions with discrete multi-valued input domains 
and multiple output classes, under any data distribution.
\end{lemma}

\begin{proof}
Extending the proof from the simplified Boolean case, we now generalize the result to functions whose input domains and output ranges each consist of a finite set of values. Let $\fml{K}$ represent the set of output classes.


We can generalize equation~\ref{eq:gbx01} as follows:

\begin{equation} \label{eq:gbx04}
\begin{aligned}
\mbb{E}_{\x\sim \mathcal{D}}[\mbf{Pr}_{\z\sim \mathcal{D}}(f(\z) = f(\x) \mid \z_{S} = \x_{S})]
= \\ \sum_{\x_{S}\sim \mathcal{D}_{S}} \mbf{Pr}(\x_{S}) \cdot \bigl[\sum_{j \in \fml{K}} \bigl(\mbf{Pr}_{\z\sim \mathcal{D}}(f(\z) = j \mid \z_{S} = \x_{S})\bigr)^2\bigr]
\end{aligned}
\end{equation}


Given the condition $\z_{S} = \x_{S}$, suppose feature $t$ can take $r$ possible values $\{v_1, \dots, v_r\}$.
Then, equation~\ref{eq:gbx02} can be generalized as follows:

\begin{equation} \label{eq:gbx05}
\begin{aligned}
&\mbb{E}_{\x\sim \mathcal{D}}[\mbf{Pr}_{\z\sim \mathcal{D}}(f(\z) = f(\x) \mid \z_{S\cup\{t\}} = \x_{S\cup\{t\}})] \\
&= \sum_{k = 1}^{r} \Bigl[ \sum_{\substack{\x_{S}\sim \fml{D}_{S} \\ x_t = v_k}} 
\mbf{Pr}(\x_{S}) \cdot \mbf{Pr}(z_t = v_k \mid \x_{S}) \cdot \bigl[\sum_{j \in \fml{K}} \bigl(\mbf{Pr}_{\z\sim \mathcal{D}}(f(\z) = j \mid \z_{S} = \x_{S}, z_t = v_k)\bigr)^2\bigr] \Bigr].
\end{aligned}
\end{equation}


Let $P_{k}$ denote $\mbf{Pr}(z_t = v_k \mid \x_{S})$,
Let $g_{k}^{j}$ denote $\mbf{Pr}_{\z\sim \mathcal{D}}(f(\z) = j \mid \z_{S} = \x_{S}, z_t = v_k)$, 
that is, $z_t$ takes the value $v_k$ and the output class is $j$.
Then, we have

\begin{equation}
\mbb{E}_{\x\sim \mathcal{D}}[\mbf{Pr}_{\z\sim \mathcal{D}}(f(\z) = f(\x) \mid \z_{S} = \x_{S})]
= \sum_{\x_{S}\sim \fml{D}_{S}} \mbf{Pr}(\x_{S}) \cdot \Bigl[\sum_{j \in \fml{K}} \bigl(\sum^{r}_{k=1} P_{k} \cdot g_{k}^{j}\bigr)^2 \Bigr]
\end{equation}
%
and
\[
\mbb{E}_{\x\sim \mathcal{D}}[\mbf{Pr}_{\z\sim \mathcal{D}}(f(\z) = f(\x) \mid \z_{S\cup\{t\}} = \x_{S\cup\{t\}})]
= \sum_{\x_{S}\sim \fml{D}_{S}} \mbf{Pr}(\x_{S}) \cdot \Bigl[\sum^{r}_{k=1} P_{k} \cdot \bigl(\sum_{j \in \fml{K}} (g_{k}^{j})^2\bigr) \Bigr]
\]


Combining these two implications, we can conclude that:

\begin{equation} \label{eq:gbx06}
\begin{aligned}
&\Delta_{t}(S,f) \\
&= \sum_{\x_{S}\sim \fml{D}_{S}} \mbf{Pr}(\x_{S}) \cdot
\Bigl[\sum^{r}_{k=1} P_{k} \cdot \bigl(\sum_{j \in \fml{K}} (g_{k}^{j})^2\bigr) - \sum_{j \in \fml{K}} \bigl(\sum^{r}_{k=1} P_{k} \cdot g_{k}^{j}\bigr)^2 \Bigr] \\
&= \sum_{\x_{S}\sim \fml{D}_{S}} \mbf{Pr}(\x_{S}) \cdot
\Bigl[\sum_{j \in \fml{K}} \bigl(\sum^{r}_{k=1} P_{k} \cdot (g_{k}^{j})^2\bigr) - \sum_{j \in \fml{K}} \bigl(\sum^{r}_{k=1} (P_{k} \cdot g_{k}^{j})^2
+ \sum_{k < l} 2 \cdot P_{k} \cdot P_{l} \cdot g_{k}^{j} \cdot g_{l}^{j}\bigr) \Bigr] \\
% &= \sum_{\x_{S}\sim \fml{D}_{S}} \mbf{Pr}(\x_{S}) \cdot
% \Bigl[\sum_{j \in \fml{K}} \bigl(\sum^{r}_{k=1} P_{k} \cdot \bigl(\sum^{r}_{l=1} P_{l}\bigr) \cdot (g_{k}^{j})^2\bigr) - \sum_{j \in \fml{K}} \bigl(\sum^{r}_{k=1} (P_{k} \cdot g_{k}^{j})^2
% + \sum_{k < l} 2 \cdot P_{k} \cdot P_{l} \cdot g_{k}^{j} \cdot g_{l}^{j}\bigr) \Bigr] \\
%&= \sum_{\x_{S}\sim \fml{D}_{S}} \mbf{Pr}(\x_{S}) \cdot
%\Bigl[\sum_{j \in \fml{K}} \bigl(\sum^{r}_{k=1} (P_{k} \cdot g_{k}^{j})^2 + \sum_{k \neq l} P_{k} \cdot P_{l} \cdot (g_{k}^{j})^2 \bigr) - \sum_{j \in \fml{K}} \bigl(\sum^{r}_{k=1} (P_{k} \cdot g_{k}^{j})^2
%+ \sum_{k \neq l} 2 \cdot P_{k} \cdot P_{l} \cdot g_{k}^{j} \cdot g_{l}^{j}\bigr) \Bigr] \\
&= \sum_{\x_{S}\sim \fml{D}_{S}} \mbf{Pr}(\x_{S}) \cdot
\Bigl[\sum_{j \in \fml{K}} \sum_{k \neq l} P_{k} \cdot P_{l} \cdot (g_{k}^{j})^2 - \sum_{j \in \fml{K}} \sum_{k < l} 2 \cdot P_{k} \cdot P_{l} \cdot g_{k}^{j} \cdot g_{l}^{j} \Bigr] \\
&= \sum_{\x_{S}\sim \fml{D}_{S}} \mbf{Pr}(\x_{S}) \cdot \sum_{j \in \fml{K}} 
\Bigl[\sum_{k < l} P_{k} \cdot P_{l} \cdot  (g_{k}^{j} - g_{l}^{j})^2\Bigr] \\
&\ge 0.
\end{aligned}
\end{equation}
%

Hence, $\mbb{E}_{\x\sim \mathcal{D}}[\mbf{Pr}_{\z\sim \mathcal{D}}(f(\z) = f(\x) \mid \z_{S} = \x_{S})]$ is monotone, which completes the proof for this setting.

\end{proof}


\begin{lemma} \label{lem:confd}
The global sufficient value function $v^g_{\text{suff}}$
is monotonic non-decreasing for functions with continuous input domains 
and multiple output classes, under any data distribution.
\end{lemma}
%
\begin{proof}
We extend the monotonicity results to functions where features are defined over continuous domains. One approach is to partition the domain of the continuous features into two disjoint subdomains and then adapt the proofs in~\cref{lem:mfd} by revising the notation accordingly. The core reasoning remains unchanged when transitioning from the discrete to the continuous case—only the notation differs.

Alternatively, we provide a direct proof below.
Assume all features lie in continuous domains, and let $p(\x_{S}, \x_{\bar{S}})$ denote the corresponding joint probability density function. Moreover, let $p(x_i)$ denote the marginal probability density function of $x_i$, and $\fml{X}_i$ denote the domain of feature $i$.
%We further extend the results on monotonicity to functions with features defined over continuous domains.
%Suppose all features have continuous domains, and let $p(\x_{S}, \x_{\bar{S}})$ denote the probability density function.
%For the proofs in~\cref{lem:mfd}, we adapt them by adjusting the notations as follows. In fact, since both summation and integration are linear, the reasoning in the proofs remains unchanged 
%when switching from the discrete to the continuous setting, except for the notation. However, for comletness --- we attach the full proof here:

We define marginal density as
\begin{equation}
p(\x_{S}) := \int_{\x_{\bar{S}} \in \fml{X}_{\bar{S}}} p(\x_{S}, \x_{\bar{S}}) d\x_{\bar{S}},
\end{equation}


where $\fml{X}_{\bar{S}}$ represents the domain of features in $\bar{S}$.

And the conditional density is defined as:

\begin{equation}
p(x_t \mid \x_{S}) := \frac{p(\x_{S}, x_t)}{p(\x_{S})}.
\end{equation}


The local probability of $S$ given $\x$ is defined as
\begin{equation}
\mbf{Pr}_{\z\sim \mathcal{D}}(f(\z) = j \mid \z_{S} = \x_{S})
=\frac{\int_{\z_{\bar{S}} \in \fml{X}_{\bar{S}}} I(f(\x_{S}, \z_{\bar{S}})=j) \cdot p(\x_{S}, \z_{\bar{S}}) d\z_{\bar{S}}}{p(\x_{S})},
\end{equation}

This term can also be decomposed into:

\begin{equation} \label{eq:con_gbx01}
\mbf{Pr}_{\z\sim \mathcal{D}}(f(\z) = j \mid \z_{S} = \x_{S})
= \int_{z_t \in \fml{X}_t} p(z_t \mid \x_{S}) \cdot \mbf{Pr}_{\z\sim \mathcal{D}}(f(\z) = j \mid \z_{S} = \x_{S}, z_t = x_t) dz_t
\end{equation}
via an additional feature $t \not \in S$.

We then modify equation~\ref{eq:gbx04} to
\begin{equation} \label{eq:con_gbx02}
\begin{aligned}
&\mbb{E}_{\x\sim \mathcal{D}}[\mbf{Pr}_{\z\sim \mathcal{D}}(f(\z) = f(\x) \mid \z_{S} = \x_{S})] \\
&= \int_{\x_{S} \in \fml{X}_{S}} p(\x_{S}) \cdot \bigl[\sum_{j \in \fml{K}} \bigl(\mbf{Pr}_{\z\sim \mathcal{D}}(f(\z) = j \mid \z_{S} = \x_{S})\bigr)^2\bigr] d\x_{S},
\end{aligned}
\end{equation}
%
while equation~\ref{eq:gbx05} is changed to
\begin{equation}\label{eq:con_gbx03} 
\begin{aligned}
&\mbb{E}_{\x\sim \mathcal{D}}[\mbf{Pr}_{\z\sim \mathcal{D}}(f(\z) = f(\x) \mid \z_{S\cup\{t\}} = \x_{S\cup\{t\}})] \\
&= \int_{\x_{S} \in \fml{X}_{S}} p(\x_{S}) \cdot \Bigl[ \int_{z_t \in \fml{X}_t} p(z_t \mid \x_{S})
\cdot \Bigl(\sum_{j\in \fml{K}} \bigl(\mbf{Pr}_{\z\sim \mathcal{D}}(f(\z) = j \mid \z_{S\cup\{t\}} = \x_{S\cup\{t\}})\bigr)^2\Bigr) dz_t \Bigr] d\x_{S}.
\end{aligned}
\end{equation}

% In order to mimic the proofs in~\cref{lem:mfd},
% we split the domain $\fml{X}_{t}$ into two disjoint subdomains $\fml{X}_{t}^{(0)}$ and $\fml{X}_{t}^{(1)}$.
% We then transform equation~\eqref{eq:con_gbx01} to
% \begin{align*}
% \mbf{Pr}_{\z\sim \mathcal{D}_p}(f(\z) = j \mid \z_{S} = \x_{S})
% &= \mbf{Pr}(z_t \in \fml{X}_{t}^{(0)} \mid \x_{S}) \cdot \mbf{Pr}_{\z\sim \mathcal{D}_p}(f(\z) = j \mid \z_{S} = \x_{S}, z_t \in \fml{X}_{t}^{(0)}) + \\
% &\hspace{1em} \mbf{Pr}(z_t \in \fml{X}_{t}^{(1)} \mid \x_{S}) \cdot \mbf{Pr}_{\z\sim \mathcal{D}_p}(f(\z) = j \mid \z_{S} = \x_{S}, z_t \in \fml{X}_{t}^{(1)}),
% \end{align*}
% where
% \[
% \mbf{Pr}(z_t \in \fml{X}_{t}^{(k)} \mid \x_{S}) = \int_{z_t \in \fml{X}_{t}^{(k)}} p(z_t \mid \x_{S}) dz_t.
% \]
% %
% This allows us to transform equation~\eqref{eq:con_gbx02} to
% \begin{equation}
% \begin{aligned}
% &\mbb{E}_{\x\sim \mathcal{D}_p}[\mbf{Pr}_{\z\sim \mathcal{D}_p}(f(\z) = f(\x) \mid \z_{S} = \x_{S})] \\
% &= \int_{\x_{S} \in \fml{X}_{S}} p(\x_{S}) \Bigl[\sum_{j \in \fml{K}} \Bigl(
% \sum_{k=0}^{1} \bigl(\mbf{Pr}(z_t \in \fml{X}_{t}^{(k)} \mid \x_{S}) \cdot 
% \mbf{Pr}_{\z\sim \mathcal{D}_p}(f(\z) = j \mid \z_{S} = \x_{S}, z_t \in \fml{X}_{t}^{(k)}) \bigr)
% \Bigr)^2\Bigr] d\x_{S},
% \end{aligned}
% \end{equation}
% %
% and~\eqref{eq:con_gbx03} to
% \begin{equation}
% \begin{aligned}
% &\mbb{E}_{\x\sim \mathcal{D}_p}[\mbf{Pr}_{\z\sim \mathcal{D}_p}(f(\z) = f(\x) \mid \z_{S\cup\{t\}} = \x_{S\cup\{t\}})] \\
% &= \int_{\x_{S} \in \fml{X}_{S}} p(\x_{S}) \Bigl[ \sum_{k=0}^{1} \mbf{Pr}(z_t \in \fml{X}_{t}^{(k)} \mid \x_{S})
% \cdot \Bigl(\sum_{j\in \fml{K}} \bigl(\mbf{Pr}_{\z\sim \mathcal{D}_p}(f(\z) = j \mid \z_{S} = \x_{S}, z_t \in \fml{X}_{t}^{(k)})\bigr)^2\Bigr) \Bigr] d\x_{S}.
% \end{aligned}
% \end{equation}
% %
% We can then leverage the same reasoning in~\cref{lem:mfd} to show that $\Delta_{t}(S,f) \ge 0$.

% Let $Y_{j}(t)$ denote $\mbf{Pr}_{\z\sim \mathcal{D}_p}(f(\z) = j \mid \z_{S} = \x_{S}, z_t = x_t)$, we can prove the monotonicity as follows:
% \begin{equation}
% \begin{aligned}
% &\Delta_{t}(S,f) \\
% &= \mbb{E}_{\x\sim \mathcal{D}_p}[\mbf{Pr}_{\z\sim \mathcal{D}_p}(f(\z) = f(\x) \mid \z_{S\cup\{t\}} = \x_{S\cup\{t\}})] - \mbb{E}_{\x\sim \mathcal{D}_p}[\mbf{Pr}_{\z\sim \mathcal{D}_p}(f(\z) = f(\x) \mid \z_{S} = \x_{S})] \\
% &= \int_{\x_{S} \in \fml{X}_{S}} p(\x_{S}) \cdot \Bigl[ \int_{z_t \in \fml{X}_t} p(z_t \mid \x_{S})
% \cdot \Bigl(\sum_{j\in \fml{K}} \bigl(\mbf{Pr}_{\z\sim \mathcal{D}_p}(f(\z) = j \mid \z_{S\cup\{t\}} = \x_{S\cup\{t\}})\bigr)^2\Bigr) dz_t \Bigr] d\x_{S} \ - \\
% &\hspace{1em} \int_{\x_{S} \in \fml{X}_{S}} p(\x_{S}) \cdot \Bigl[\sum_{j \in \fml{K}} 
% \Bigl(\int_{z_t \in \fml{X}_t} p(z_t \mid \x_{S}) \cdot \mbf{Pr}_{\z\sim \mathcal{D}_p}(f(\z) = j \mid \z_{S} = \x_{S}, z_t = x_t) dz_t\Bigr)^2\Bigr] d\x_{S} \\
% %
% &= \int_{\x_{S} \in \fml{X}_{S}} p(\x_{S}) \cdot \Bigl[ \sum_{j\in \fml{K}}
% \Bigl(\int_{z_t \in \fml{X}_t} p(z_t \mid \x_{S}) \cdot \bigl(\mbf{Pr}_{\z\sim \mathcal{D}_p}(f(\z) = j \mid \z_{S\cup\{t\}} = \x_{S\cup\{t\}})\bigr)^2 dz_t \Bigr) \Bigr] d\x_{S} \ - \\
% &\hspace{1em} \int_{\x_{S} \in \fml{X}_{S}} p(\x_{S}) \cdot \Bigl[\sum_{j \in \fml{K}} 
% \Bigl(\int_{z_t \in \fml{X}_t} p(z_t \mid \x_{S}) \cdot \mbf{Pr}_{\z\sim \mathcal{D}_p}(f(\z) = j \mid \z_{S} = \x_{S}, z_t = x_t) dz_t\Bigr)^2\Bigr] d\x_{S} \\
% %
% &= \int_{\x_{S} \in \fml{X}_{S}} p(\x_{S}) \cdot \Bigl[ \sum_{j\in \fml{K}} \Bigl(\mbb{E}[Y_{j}(t)^2] - (\mbb{E}[Y_{j}(t)])^2\Bigr) \Bigr] d\x_{S} \\
% &\ge 0.
% \end{aligned}
% \end{equation}
% %
% $\Delta_{t}(S,f) \ge 0$ holds as a result of applying the fact that $\mbb{E}[Y_{j}(t)^2] - (\mbb{E}[Y_{j}(t)])^2 \ge 0$.

For simplicity, let $Y_{j}(t)$ denote $\mbf{Pr}_{\z\sim \mathcal{D}}(f(\z) = j \mid \z_{S} = \x_{S}, z_t = v_t)$, and $p_{t\mid S}$ denote $p(z_t \mid \x_{S})$.
Note that both $Y_{j}(t)$ and $p_{t\mid S}$ are associated with a fixed $\x_{S}$
\footnote{Changing $\x_{S}$ alters their values, we omit $\x_S$ in the notation for brevity.}.
We prove monotonicity as follows:
\begin{equation}
\begin{aligned}
&\Delta_{t}(S,f) \\
&= \mbb{E}_{\x\sim \mathcal{D}}[\mbf{Pr}_{\z\sim \mathcal{D}}(f(\z) = f(\x) \mid \z_{S\cup\{t\}} = \x_{S\cup\{t\}})] - \mbb{E}_{\x\sim \mathcal{D}}[\mbf{Pr}_{\z\sim \mathcal{D}}(f(\z) = f(\x) \mid \z_{S} = \x_{S})] \\
&= \int_{\x_{S} \in \fml{X}_{S}} p(\x_{S}) \cdot \Bigl[ \int_{z_t \in \fml{X}_t} p_{t\mid S}
\cdot \Bigl(\sum_{j\in \fml{K}} \bigl(\mbf{Pr}_{\z\sim \mathcal{D}}(f(\z) = j \mid \z_{S\cup\{t\}} = \x_{S\cup\{t\}})\bigr)^2\Bigr) dz_t \Bigr] d\x_{S} \ - \\
&\hspace{1em} \int_{\x_{S} \in \fml{X}_{S}} p(\x_{S}) \cdot \Bigl[\sum_{j \in \fml{K}} 
\Bigl(\int_{z_t \in \fml{X}_t} p_{t\mid S} \cdot \mbf{Pr}_{\z\sim \mathcal{D}}(f(\z) = j \mid \z_{S} = \x_{S}, z_t = x_t) dz_t\Bigr)^2\Bigr] d\x_{S} \\
%
&= \int_{\x_{S} \in \fml{X}_{S}} p(\x_{S}) \cdot \Bigl[ \sum_{j\in \fml{K}} \Bigl(
\int_{\fml{X}_t} p_{t\mid S} \cdot Y_{j}(t)^2 dz_t - \bigl( \int_{\fml{X}_t} p_{t\mid S} \cdot Y_{j}(t) dz_t \bigr)^2 \Bigr) \Bigr] d\x_{S} \\
%
&= \int_{\x_{S} \in \fml{X}_{S}} p(\x_{S}) \cdot \Bigl[ \sum_{j\in \fml{K}} \frac{2}{2} \cdot \Bigl(
\int_{\fml{X}_t} \int_{\fml{X}_{t'}} p_{t\mid S} \cdot p_{t' \mid S} \cdot \bigl(Y_{j}(t)^2 - Y_{j}(t) Y_{j}(t')\bigr) dz_t dz_{t'} \Bigr) \Bigr] d\x_{S} \\
&= \int_{\x_{S} \in \fml{X}_{S}} p(\x_{S}) \cdot \Bigl[ \sum_{j\in \fml{K}} \frac{1}{2} \cdot \Bigl(
\int_{\fml{X}_t} \int_{\fml{X}_{t'}} p_{t\mid S} \cdot p_{t' \mid S} \cdot \bigl(Y_{j}(t) - Y_{j}(t')\bigr)^2 dz_t dz_{t'} \Bigr) \Bigr] d\x_{S} \\
%
&\ge 0.
\end{aligned}
\end{equation}

\end{proof}


\globalmonotonicity*
%
\begin{proof}
As a consequence of~\cref{lem:bfd,lem:mfd,lem:confd}, we conclude that the global sufficient value function $v^g_{\text{suff}}$ is monotonic non-decreasing for any well-defined classification function --- whether its features are Boolean, discrete, continuous, or a combination --- under any data distribution.

%Now consider the global contrastive value function $v^g_{\text{con}}$, where $S$ is the set of features permitted to change. By examining the definitions of both $v^g_{\text{suff}}$ and $v^g_{\text{con}}$, we observe that:
%we infer that
%the global sufficient value function $v^g_{\text{suff}}$ is monotonic non-decreasing 
%for any well-defined classification function (including those with Boolean, discrete, continuous, or hybrid feature domains) under any data distribution.

%For the global contrastive value function $v^g_{\text{con}}$,
%let $S$ denote the set of features that are allowed to change.
%By inspecting the definition of $v^g_{\text{suff}}$ and $v^g_{\text{con}}$, we have:

\begin{equation}
v^g_{\text{con}}(S) = 1 - v^g_{\text{suff}}(\Bar{S})
\end{equation}

This implies that $v^g_{\text{con}}$ is monotonic non-increasing with respect to adding features to $\Bar{S}$. Equivalently, $v^g_{\text{con}}$ is monotonic non-decreasing as features are removed from $\Bar{S}$ and added to $S$, thus completing the proof.

%this means that $v^g_{\text{con}}$ is monotonic non-increasing as we add more features to $\Bar{S}$. In other words, $v^g_{\text{con}}$ is monotonic non-decreasing as we remove features from $\Bar{S}$ and add them to $S$, hence concluding the proof.

\end{proof}





\section{Proof of Proposition~\ref{globalsupermod_prop}}
\label{supermodularity_proofs_appendix}

\globalsupermodular*


\emph{Proof.} This section establishes the supermodularity property for both global sufficient and contrastive explanations. As with monotonicity, we begin with Boolean functions and progressively generalize to multi-valued discrete and continuous input domains, and finally consider the case of well-defined classification functions. Unlike the monotonicity property, however, these proofs require the assumption of feature independence. We also provide a counterexample to show that supermodularity may fail when this assumption is relaxed.


\begin{lemma} \label{lem:bfid}
The global sufficient value function $v^g_{\text{suff}}$
is supermodular for Boolean functions, under the assumption of feature independence.
\end{lemma}


\begin{proof}
Particularly, we will prove that:

\begin{equation}
\Delta_{t}(S,f) \le \Delta_{t}(S',f),
\end{equation}

where $S\subseteq S'$ and $t \not\in S'$. That is, $\Delta_{t}(S,f)$ is monotone.

Clearly, if $S' = S$, then $\Delta_{t}(S,f) \le \Delta_{t}(S',f)$ holds.
Let $S' = S\cup\{i\}$, where $i\not\in S$ and $i \neq t$.
For the local probability $\mbf{Pr}_{\z\sim \mathcal{D}}(f(\z) = f(\x) \mid \z_{S} = \x_{S})$,
under the assumption of feature independence, we have that the following holds:

\begin{equation}
\begin{aligned}
&\mbf{Pr}_{\z}(f(\z) = f(\x) \mid \z_{S} = \x_{S}) \\
&= \mbf{Pr}(z_i = 1) \cdot \mbf{Pr}(z_t = 1) \cdot \mbf{Pr}_{\z}(f(\z) = f(\x) \mid \z_{S} = \x_{S}, z_i = 1, z_t = 1) \ + \\
&\hspace{1em} \mbf{Pr}(z_i = 1) \cdot \mbf{Pr}(z_t = 0) \cdot \mbf{Pr}_{\z}(f(\z) = f(\x) \mid \z_{S} = \x_{S}, z_i = 1, z_t = 0) \ + \\
&\hspace{1em} \mbf{Pr}(z_i = 0) \cdot \mbf{Pr}(z_t = 1) \cdot \mbf{Pr}_{\z}(f(\z) = f(\x) \mid \z_{S} = \x_{S}, z_i = 0, z_t = 1) \ + \\
&\hspace{1em} \mbf{Pr}(z_i = 0) \cdot \mbf{Pr}(z_t = 0) \cdot \mbf{Pr}_{\z}(f(\z) = f(\x) \mid \z_{S} = \x_{S}, z_i = 0, z_t = 0).
\end{aligned}
\end{equation}


Let $P_{1}$, $P_{0}$, $Q_{1}$, and $Q_{0}$ denote $\mbf{Pr}(z_t = 1)$, $\mbf{Pr}(z_t = 0)$, $\mbf{Pr}(z_i = 1)$, and $\mbf{Pr}(z_i = 0)$, respectively.
Moreover, let $g_{11}^{+}$, $g_{10}^{+}$, $g_{01}^{+}$, and $g_{00}^{+}$ denote
$\mbf{Pr}_{\z \sim \mathcal{D}}(f(\z) = 1 \mid \z_{S} = \x_{S}, z_i = 1, z_t = 1)$, 
$\mbf{Pr}_{\z \sim \mathcal{D}}(f(\z) = 1 \mid \z_{S} = \x_{S}, z_i = 1, z_t = 0)$,
$\mbf{Pr}_{\z \sim \mathcal{D}}(f(\z) = 1 \mid \z_{S} = \x_{S}, z_i = 0, z_t = 1)$, and 
$\mbf{Pr}_{\z \sim \mathcal{D}}(f(\z) = 1 \mid \z_{S} = \x_{S}, z_i = 0, z_t = 0)$, respectively.

We further decompose equation~\ref{eq:gbx03} into:

\begin{equation}
\begin{aligned}
\Delta_{t}(S,f)
% &= \sum_{\x_{S}\in \fml{D}_{S}} 2 \cdot \mbf{Pr}(\x_{S}) \cdot \mbf{Pr}(z_t = 1) \cdot \mbf{Pr}(z_t = 0)
% \cdot \bigl[\mbf{Pr}_{\z}(f^{+} \mid \z_{S} = \x_{S}, z_t = 1) - \mbf{Pr}_{\z}(f^{+} \mid \z_{S} = \x_{S}, z_t = 0)\bigr]^2 \\
%
&= \sum_{\x_{S}\sim \fml{D}_{S}} 2 \cdot \mbf{Pr}(\x_{S}) \cdot 
P_{1} P_{0} \cdot \bigl[(Q_{1} \cdot g_{11}^{+} + Q_{0} \cdot g_{01}^{+}) - (Q_{1} \cdot g_{10}^{+} + Q_{0} \cdot g_{00}^{+})\bigr]^2 \\
&= \sum_{\x_{S}\sim \fml{D}_{S}} 2 \cdot \mbf{Pr}(\x_{S}) \cdot 
P_{1} P_{0} \cdot \bigl[Q_{1} \cdot (g_{11}^{+} - g_{10}^{+}) + Q_{0} \cdot (g_{01}^{+} - g_{00}^{+})\bigr]^2.
\end{aligned}
\end{equation}

Similarly, the following also holds:

\begin{equation}
\begin{aligned}
\Delta_{t}(S\cup\{i\},f)
% &= \sum_{\substack{\x_{S}\sim \fml{D}_{S} \\ x_i = 1}} 2 \cdot \mbf{Pr}(\x_{S\cup\{i\}}) \cdot 
% P_{1} P_{0} \cdot \bigl[\mbf{Pr}_{\z}(f^{+} \mid \z_{S\cup\{i\}} = \x_{S\cup\{i\}}, z_t = 1) - \mbf{Pr}_{\z}(f^{+} \mid \z_{S\cup\{i\}} = \x_{S\cup\{i\}}, z_t = 0)\bigr]^2 \ + \\
% &\hspace{1em} \sum_{\substack{\x_{S}\sim \fml{D}_{S} \\ x_i = 0}} 2 \cdot \mbf{Pr}(\x_{S\cup\{i\}}) \cdot
% P_{1} P_{0} \cdot \bigl[\mbf{Pr}_{\z}(f^{+} \mid \z_{S\cup\{i\}} = \x_{S\cup\{i\}}, z_t = 1) - \mbf{Pr}_{\z}(f^{+} \mid \z_{S\cup\{i\}} = \x_{S\cup\{i\}}, z_t = 0)\bigr]^2 \\
&= \sum_{\x_{S}\sim \fml{D}_{S}}
2 \cdot \mbf{Pr}(\x_{S}) \cdot P_{1} P_{0} \cdot \bigl[Q_{1} \cdot (g_{11}^{+} - g_{10}^{+})^2 + Q_{0} \cdot (g_{01}^{+} - g_{00}^{+})^2\bigr].
\end{aligned}
\end{equation}
%
Thus, we get
\begin{equation}
\begin{aligned}
\Delta_{t}(S\cup\{i\},f) - \Delta_{t}(S,f)
&= \sum_{\x_{S}\sim \fml{D}_{S}}
2 \cdot \mbf{Pr}(\x_{S}) \cdot P_{1} P_{0} \cdot Q_{1} Q_{0} \cdot (g_{11}^{+} - g_{10}^{+} - g_{01}^{+} + g_{00}^{+})^2 \\
&\ge 0.
\end{aligned}
\end{equation}
%
This implies that $\Delta_{t}(S,f)$ is monotone, that is, $c(S\cup\{t\},f) - c(S,f) \le c(S'\cup\{t\},f) - c(S',f)$.
\end{proof}


\begin{lemma} \label{lem:mfid}
The global sufficient value function $v^g_{\text{suff}}$
is supermodular for functions with discrete multi-valued input domains 
and multiple output classes, under the assumption of feature independence.
\end{lemma}
%
\begin{proof}
We extend the results of the previous section to multi-valued discrete functions, 
again with the assumption of feature independence.
%
We use $P_{k}$ to denote $\mbf{Pr}(z_t = v_k)$, and $Q_{k'}$ to denote $\mbf{Pr}(z_i = v_k')$.
Moreover, we use $g_{k'k}^{j}$ to denote $\mbf{Pr}_{\z\sim \mathcal{D}}(f(\z) = j \mid \z_{S} = \x_{S}, z_i = v_k', z_t = v_k)$.
We decompose equation~\ref{eq:gbx06} to
%
\begin{equation}
\begin{aligned}
\Delta_{t}(S,f)
% &= \sum_{\x_{S}\sim \fml{D}_{S}} \mbf{Pr}(\x_{S}) \cdot \sum_{j \in \fml{K}} 
% \Bigl[\sum_{k < l} P_{k} P_{l} \cdot \bigl(\mbf{Pr}_{\z}(f(\z) = j \mid \z_{S} = \x_{S}, z_t = v_k) - \mbf{Pr}_{\z}(f(\z) = j \mid \z_{S} = \x_{S}, z_t = v_l)\bigr)^2\Bigr] \\
&= \sum_{\x_{S}\sim \fml{D}_{S}} \mbf{Pr}(\x_{S}) \cdot \sum_{j \in \fml{K}} 
\Bigl[\sum_{k < l} P_{k} P_{l} \cdot \Bigl(\sum_{k' = 1}^{r'} Q_{k'} \cdot g_{k'k}^{j} - \sum_{k' = 1}^{r'} Q_{k'} \cdot g_{k'l}^{j}\Bigr)^2\Bigr] \\
&= \sum_{\x_{S}\sim \fml{D}_{S}} \mbf{Pr}(\x_{S}) \cdot \sum_{j \in \fml{K}} 
\Bigl[\sum_{k < l} P_{k} P_{l} \cdot \Bigl(\sum_{k' = 1}^{r'} Q_{k'} \cdot (g_{k'k}^{j} - g_{k'l}^{j}) \Bigr)^2\Bigr].
\end{aligned}
\end{equation}
%
Likewise, we have
%
\begin{equation}
\begin{aligned}
\Delta_{t}(S\cup\{i\},f)
&= \sum_{k' = 1}^{r'} \Bigl[ \sum_{\substack{\x_{S}\sim \fml{D}_{S} \\ x_i = v_k'}} \mbf{Pr}(\x_{S}) \cdot \mbf{Pr}(z_i = v_k') \cdot \sum_{j \in \fml{K}} 
\Bigl(\sum_{k < l} P_{k} P_{l} \cdot (g_{k'k}^{j} - g_{k'l}^{j})^2\Bigr) \Bigr] \\
&= \sum_{\x_{S}\sim \fml{D}_{S}} \mbf{Pr}(\x_{S}) \cdot \Bigl[ \sum_{k' = 1}^{r'} Q_{k'} \cdot \sum_{j \in \fml{K}} 
\Bigl(\sum_{k < l} P_{k} P_{l} \cdot (g_{k'k}^{j} - g_{k'l}^{j})^2\Bigr) \Bigr] \\
%
&= \sum_{\x_{S}\sim \fml{D}_{S}} \mbf{Pr}(\x_{S}) \cdot \sum_{j \in \fml{K}} 
\Bigl[\sum_{k < l} P_{k} P_{l} \cdot \sum_{k' = 1}^{r'} Q_{k'} \cdot (g_{k'k}^{j} - g_{k'l}^{j})^2\Bigr].
\end{aligned}
\end{equation}
%
Therefore, we get 
\begin{equation}
\begin{aligned}
\Delta_{t}&(S\cup\{i\},f) - \Delta_{t}(S,f) \\
&= \sum_{\x_{S}\sim \fml{D}_{S}} \mbf{Pr}(\x_{S}) \cdot \sum_{j \in \fml{K}} 
\Bigl[\sum_{k < l} \sum_{k' < l'} P_{k} P_{l} Q_{k'} Q_{l'} \cdot (g_{k'k}^{j} - g_{k'l}^{j} - g_{l'k}^{j} + g_{l'l}^{j})^2\Bigr] \\
&\ge 0.
%
\end{aligned}
\end{equation}

\end{proof}


\globalsupermodular*
%
\begin{proof}

We can generalize the supermodularity results established in Lemma~\ref{lem:mfid} to the continuous case, assuming that $\fml{D}_p$ satisfies feature independence.
One approach is to partition the domain of the continuous features into two disjoint subdomains and then adapt the proofs in Lemma~\ref{lem:mfid} by revising the notation accordingly.

Alternatively, we provide a direct proof below.
Under this assumption, the marginal and conditional densities simplify as follows:
%We can extend the supermodularity results proven in~\cref{lem:mfid} to the continuous setting, provided that $\fml{D}_p$ exhibits feature independence.
%
%In such a setting, the marginal density and the conditional density become
\begin{equation}
p(\x_{S}) = \prod_{i\in S} p(x_i),~\textnormal{and}~~p(x_t \mid \x_{S}) = p(x_t).
\end{equation}

Consequently, the local probability of $S$ conditioned on $\x$ becomes:

\begin{equation}
\mbf{Pr}_{\z\sim \mathcal{D}}(f(\z) = j \mid \z_{S} = \x_{S})
=\int_{\z_{\bar{S}} \in \fml{X}_{\bar{S}}} I(f(\x_{S}, \z_{\bar{S}})=j) \cdot \prod_{i\in \bar{S}} p(z_i) d\z_{\bar{S}}.
\end{equation}

% Following a similar approach to the proof in~\cref{lem:confd}, we partition the domain $\fml{X}{t}$ into two disjoint subsets, $\fml{X}{t}^{(0)}$ and $\fml{X}{t}^{(1)}$. Likewise, we divide $\fml{X}{i}$ into $\fml{X}{i}^{(0)}$ and $\fml{X}{i}^{(1)}$. This allows us to directly apply the argument from~\cref{lem:mfid} to establish supermodularity. 

We let $Y_{j}(i,t)$ denote $\mbf{Pr}_{\z\sim \mathcal{D}}(f(\z) = j \mid \z_{S} = \x_{S}, z_i = x_i, z_t = x_t)$, where $i\not\in S$, $t\not\in S$ and $i \neq t$.
Evidently,
\[
Y_{j}(i,t) = Y_{j}(t,i),~~\textnormal{and}~~Y_{j}(t) = \int_{\fml{X}_{i}} p_{i} \cdot Y_{j}(i,t) dz_i.
\]
%
We have:
\begin{equation}
\begin{aligned}
&\Delta_{t}(S,f) \\
&= \int_{\x_{S} \in \fml{X}_{S}} p(\x_{S}) \cdot \Bigl[ \sum_{j\in \fml{K}} \frac{1}{2} \Bigl(
\int_{\fml{X}_t} \int_{\fml{X}_{t'}} p_{t} p_{t'} \bigl(Y_{j}(t) - Y_{j}(t')\bigr)^2 dz_t dz_{t'} \Bigr) \Bigr] d\x_{S} \\
%
&= \int_{\x_{S} \in \fml{X}_{S}} p(\x_{S}) \cdot \Bigl[ \sum_{j\in \fml{K}} \frac{1}{2} \Bigl(
\int_{\fml{X}_t} \int_{\fml{X}_{t'}} p_{t} p_{t'} \bigl(\int_{\fml{X}_{i}} p_{i} \cdot Y_{j}(i,t) dz_i - \int_{\fml{X}_{i}} p_{i} \cdot Y_{j}(i,t') dz_i \bigr)^2 dz_t dz_{t'} \Bigr) \Bigr] d\x_{S} \\
%
&= \int_{\x_{S} \in \fml{X}_{S}} p(\x_{S}) \cdot \Bigl[ \sum_{j\in \fml{K}} \frac{1}{2} \Bigl(
\int_{\fml{X}_t} \int_{\fml{X}_{t'}} p_{t} p_{t'} \bigl(\int_{\fml{X}_{i}} p_{i} \cdot (Y_{j}(i,t) - Y_{j}(i,t')) dz_i \bigr)^2 dz_t dz_{t'} \Bigr) \Bigr] d\x_{S},
\end{aligned}
\end{equation}
where
\begin{align*}
&\int_{\fml{X}_t} \int_{\fml{X}_{t'}} p_{t} p_{t'} \bigl(\int_{\fml{X}_{i}} p_{i} \cdot (Y_{j}(i,t) - Y_{j}(i,t')) dz_i \bigr)^2 dz_t dz_{t'} \\
&= \int_{\fml{X}_t} \int_{\fml{X}_{t'}} \int_{\fml{X}_{i}} \int_{\fml{X}_{i'}} p_{t} p_{t'} p_{i} p_{i'} \cdot
(Y_{j}(i,t) - Y_{j}(i,t')) \cdot (Y_{j}(i',t) - Y_{j}(i',t')) dz_i dz_{i'} dz_t dz_{t'}.
\end{align*}
%
Let $S' = S\cup\{i\}$, then we have:
\begin{equation}
\begin{aligned}
&\Delta_{t}(S',f) \\
&= \int_{\x_{S'} \in \fml{X}_{S'}} p(\x_{S'}) \cdot \Bigl[ \sum_{j\in \fml{K}} \frac{1}{2} \Bigl(
\int_{\fml{X}_t} \int_{\fml{X}_{t'}} p_{t} p_{t'} \cdot \bigl(Y_{j}(i,t) - Y_{j}(i,t')\bigr)^2 dz_t dz_{t'} \Bigr) \Bigr] d\x_{S'} \\
%
&= \int_{\x_{S} \in \fml{X}_{S}} p(\x_{S}) \cdot \Bigl[ \sum_{j\in \fml{K}} \frac{1}{2} \int_{\fml{X}_i} p_{i} \cdot \Bigl(
\int_{\fml{X}_t} \int_{\fml{X}_{t'}} p_{t} p_{t'} \cdot \bigl(Y_{j}(i,t) - Y_{j}(i,t')\bigr)^2 dz_t dz_{t'} \Bigr) dz_i \Bigr] d\x_{S} \\
%
&= \int_{\x_{S} \in \fml{X}_{S}} p(\x_{S}) \cdot \Bigl[ \sum_{j\in \fml{K}} \frac{1}{2} \Bigl(
\int_{\fml{X}_i} \int_{\fml{X}_t} \int_{\fml{X}_{t'}} p_{i} p_{t} p_{t'} \cdot \bigl(Y_{j}(i,t) - Y_{j}(i,t')\bigr)^2 dz_t dz_{t'} dz_i \Bigr) \Bigr] d\x_{S},
\end{aligned}
\end{equation}
%
where
\begin{align*}
&\int_{\fml{X}_i} \int_{\fml{X}_t} \int_{\fml{X}_{t'}} p_{i} p_{t} p_{t'} \cdot \bigl(Y_{j}(i,t) - Y_{j}(i,t')\bigr)^2 dz_t dz_{t'} dz_i \\
&= \int_{\fml{X}_i} \int_{\fml{X}_{i'}} \int_{\fml{X}_t} \int_{\fml{X}_{t'}} p_{i} p_{i'} p_{t} p_{t'} \cdot \bigl(Y_{j}(i,t) - Y_{j}(i,t')\bigr)^2 dz_t dz_{t'} dz_i dz_{i'}.
\end{align*}
%
Let $\alpha(i,t,t')$ denote $Y_{j}(i,t) - Y_{j}(i,t')$ and $\beta(i',t,t')$ denote $Y_{j}(i',t) - Y_{j}(i',t')$, we can infer that
\begin{equation}
\begin{aligned}
&\Delta_{t}(S',f) - \Delta_{t}(S,f) \\
&= \int_{\x_{S} \in \fml{X}_{S}} p(\x_{S}) \Bigl[ \sum_{j\in \fml{K}} \frac{1}{4} \Bigl( 
\int_{\fml{X}_t} \int_{\fml{X}_{t'}} \int_{\fml{X}_{i}} \int_{\fml{X}_{i'}} p_{t} p_{t'} p_{i} p_{i'}
(\alpha(i,t,t') - \beta(i',t,t'))^2 dz_i dz_{i'} dz_t dz_{t'}
 \Bigr) \Bigr] d\x_{S} \\
 &\ge 0.
\end{aligned}
\end{equation}

Consequently, the global sufficient value function $v^g_{\text{suff}}$ retains the supermodularity property for any valid classification function, regardless of whether its features are Boolean, discrete, continuous, or a mix thereof.

%By adopting similar proof ideas from~\cref{lem:confd}, we split the domain $\fml{X}_{t}$ into two disjoint subdomains, $\fml{X}_{t}^{(0)}$ and $\fml{X}_{t}^{(1)}$.
%Similarly, we split the domain $\fml{X}_{i}$ into two disjoint subdomains,
%$\fml{X}_{i}^{(0)}$ and $\fml{X}_{i}^{(1)}$.
%Then we can reuse the proof in~\cref{lem:mfid} to prove supermodularity.
%This further implies that the supermodularity property of the global sufficient value function $v^g_{\text{suff}}$ holds for any well-defined classification function, including those with Boolean, discrete, continuous, or hybrid feature domains.


We now demonstrate that the local probabilistic sufficient value function is neither submodular nor supermodular.
As a simple illustration, fix the uniform distribution on the three variables $\{x_1, x_2, x_3\}$
and define a Boolean function as follows:
\begin{equation}
f(x_1,x_2) = (x_1 \lor x_2) \land x_3.
\end{equation}

Consider the input $\x = (1,1,1)$ which is classified as $1$.
For the local probabilistic sufficient value function $v^{l}_{\text{suff}}$, we have that the following holds:

\begin{align*}
&v^{l}_{\text{suff}}(\emptyset) = \frac{3}{8},~~v^{l}_{\text{suff}}(\{1\}) = \frac{1}{2},~~v^{l}_{\text{suff}}(\{2\}) = \frac{1}{2},~~v^{l}_{\text{suff}}(\{3\}) = \frac{3}{4}, \\
&v^{l}_{\text{suff}}(\{1,2\}) = \frac{1}{2},~~v^{l}_{\text{suff}}(\{1,3\}) = 1,~~v^{l}_{\text{suff}}(\{2,3\}) = 1,~~v^{l}_{\text{suff}}(\{1,2,3\}) = 1.
\end{align*}
%
Since $\bigl(v^{l}_{\text{suff}}(\{1,2\}) - v^{l}_{\text{suff}}(\{2\})\bigr) < \bigl(v^{l}_{\text{suff}}(\{1\}) - v^{l}_{\text{suff}}(\emptyset)\bigr) < \bigl(v^{l}_{\text{suff}}(\{1,3\}) - v^{l}_{\text{suff}}(\{3\})\bigr)$,
the function $v^{l}_{\text{suff}}$ is neither supermodular nor submodular.
 
\end{proof}


\section{Proof of Proposition~\ref{globalsubmodular_prop}}
\label{globalsubmodular_prop_sec}

\globalsubmodular*

\begin{proof}

Once we establish that the global sufficient setting is supermodular, it directly follows that the global contrastive setting is submodular. Let $S \subseteq [n]$ be the set of features allowed to change, and suppose $S \subseteq S'$ with $i \notin S'$. To compare the marginal contributions $v^g_{\text{con}}(S \cup \{i\}) - v^g_{\text{con}}(S)$ and $v^g_{\text{con}}(S' \cup \{i\}) - v^g_{\text{con}}(S')$, we equivalently compare
$\bigl(1 - v^g_{\text{suff}}(\Bar{S} \setminus \{i\})\bigr) - \bigl(1 - v^g_{\text{suff}}(\Bar{S})\bigr) \quad \text{and} \quad \bigl(1 - v^g_{\text{suff}}(\Bar{S'} \setminus \{i\})\bigr) - \bigl(1 - v^g_{\text{suff}}(\Bar{S'})\bigr).$
Since $S \subseteq S'$, it follows that $\Bar{S'} \subseteq \Bar{S}$, and therefore we obtain:


%Let $S\in [n]$ denote the set of features that are allowed to change.
%Let $S\subseteq S'$ and $i\not\in S'$.
%To compare $v^g_{\text{con}}(S\cup\{i\}) - v^g_{\text{con}}(S)$ and $v^g_{\text{con}}(S'\cup\{i\}) - v^g_{\text{con}}(S')$, we essentially compare
%$\bigl(1 - v^g_{\text{suff}}(\Bar{S} \setminus\{i\})\bigr) - \bigl(1 - v^g_{\text{suff}}(\Bar{S})\bigr)$ and $\bigl(1 - v^g_{\text{suff}}(\Bar{S'} \setminus\{i\})\bigr) - \bigl(1 - v^g_{\text{suff}}(\Bar{S'})\bigr)$.
%Since $S\subseteq S'$, we have $\Bar{S'}\subseteq \Bar{S}$, then we get 

\begin{equation}
\bigl(v^g_{\text{suff}}(\Bar{S}) - v^g_{\text{suff}}(\Bar{S} \setminus\{i\})\bigr) \geq
\bigl(v^g_{\text{suff}}(\Bar{S'}) - v^g_{\text{suff}}(\Bar{S'} \setminus\{i\})\bigr)
\end{equation}

Which is equivalent to:

\begin{equation}
\bigl(v^g_{\text{con}}(S\cup\{i\}) - v^g_{\text{con}}(S)\bigr) \geq
\bigl(v^g_{\text{con}}(S'\cup\{i\}) - v^g_{\text{con}}(S')\bigr).
\end{equation}

Thus, the global contrastive value function $v^g_{\text{con}}$ is submodular, hence concluding this proof.


We will now demonstrate that the local probabilistic contrastive value function is neither submodular nor supermodular. As a simple illustration, fix the uniform distribution on the three variables $\{x_1, x_2, x_3\}$
and define a Boolean function as follows:

\begin{equation}
f(x_1,x_2):= (x_1 \lor x_2) \land x_3.
\end{equation}

Consider the input $\x = (1,1,1)$ which is classified as $1$. For the local probabilistic contrastive value function $v^{l}_{\text{con}}$, we have:

\begin{align*}
&v^{l}_{\text{con}}(\emptyset) = 1,~~v^{l}_{\text{con}}(\{1\}) = 1,~~v^{l}_{\text{con}}(\{2\}) = 1,~~v^{l}_{\text{con}}(\{3\}) = \frac{1}{2}, \\
&v^{l}_{\text{con}}(\{1,2\}) = \frac{3}{4},~~v^{l}_{\text{con}}(\{1,3\}) = \frac{1}{2},~~v^{l}_{\text{con}}(\{2,3\}) = \frac{1}{2},~~v^{l}_{\text{con}}(\{1,2,3\}) = \frac{3}{8}.
\end{align*}


Since $\bigl(v^{l}_{\text{con}}(\{1,2\}) - v^{l}_{\text{con}}(\{2\})\bigr) < \bigl(v^{l}_{\text{con}}(\{1,2,3\}) - v^{l}_{\text{con}}(\{2,3\})\bigr) < \bigl(v^{l}_{\text{con}}(\{1\}) - v^{l}_{\text{con}}(\emptyset)\bigr)$,
the function $v^{l}_{\text{con}}$ is neither submodular nor supermodular.

\end{proof}

In the rest of this section, we show that the feature independence assumption is not just sufficient but also necessary for supermodularity. We illustrate this with the following counterexample.

Define a Boolean function:

\begin{equation}
f(x_1, x_2): = x_1
\end{equation}

For this boolean function, it holds that $f(00) = 0$, $f(01) = 0$, $f(10) = 1$, and $f(11) = 1$. Moreover, define a distribution $\mathcal{D}$ on $x_1$ and $x_2$. Recall that the global sufficient value function is:

\[
\mbb{E}_{\x\sim \mathcal{D}}[\mbf{Pr}_{\z\sim \mathcal{D}}(f(\z) = f(\x) \mid \z_{S} = \x_{S})]
= \sum_{\x\sim \mathcal{D}} \mbf{Pr}(\x) \cdot \mbf{Pr}_{\z\sim \mathcal{D}}(f(\z) = f(\x) \mid \z_{S} = \x_{S})
\]
%
In addition, if we can show that
\[
\Delta_{t}(S,f) \le \Delta_{t}(S',f),
\]
where $S\subseteq S'$ and $t \not\in S'$, then the global sufficient value function is supermodular. When $S = \emptyset$, we have that the following holds:

\begin{equation*}
\begin{aligned}
&\mbb{E}_{\x\sim \mathcal{D}}[\mbf{Pr}_{\z\sim \mathcal{D}}(f(\z) = f(\x) \mid \z_{S} = \x_{S})] \\
&= \mbf{Pr}(x_1 = 0, x_2 = 0) \cdot \mbf{Pr}_{\z\sim \mathcal{D}}(f(\z) = 0) \ + \mbf{Pr}(x_1 = 0, x_2 = 1) \cdot \mbf{Pr}_{\z\sim \mathcal{D}}(f(\z) = 0) \ + \\
&\hspace{1em} \mbf{Pr}(x_1 = 1, x_2 = 0) \cdot \mbf{Pr}_{\z\sim \mathcal{D}}(f(\z) = 1) \ + 
\mbf{Pr}(x_1 = 1, x_2 = 1) \cdot \mbf{Pr}_{\z\sim \mathcal{D}}(f(\z) = 1) 
\end{aligned}
\end{equation*}

Which is also equivalent to:

\begin{equation*}
\begin{aligned}
& \mbf{Pr}(x_1 = 0, x_2 = 0) \cdot \bigl(\mbf{Pr}(x_1 = 0, x_2 = 0) + \mbf{Pr}(x_1 = 0, x_2 = 1)\bigr) \ + \\
&\hspace{1em} \mbf{Pr}(x_1 = 0, x_2 = 1) \cdot \bigl(\mbf{Pr}(x_1 = 0, x_2 = 0) + \mbf{Pr}(x_1 = 0, x_2 = 1)\bigr) \ + \\
&\hspace{1em} \mbf{Pr}(x_1 = 1, x_2 = 0) \cdot \bigl(\mbf{Pr}(x_1 = 1, x_2 = 0) + \mbf{Pr}(x_1 = 1, x_2 = 1)\bigr) \ + \\
&\hspace{1em} \mbf{Pr}(x_1 = 1, x_2 = 1) \cdot \bigl(\mbf{Pr}(x_1 = 1, x_2 = 0) + \mbf{Pr}(x_1 = 1, x_2 = 1)\bigr) \\
%
&= \bigl(\mbf{Pr}(x_1 = 0, x_2 = 0) + \mbf{Pr}(x_1 = 0, x_2 = 1)\bigr)^2 \ + \bigl(\mbf{Pr}(x_1 = 1, x_2 = 0) + \mbf{Pr}(x_1 = 1, x_2 = 1)\bigr)^2
\end{aligned}
\end{equation*}


When $S = \{2\}$, we have that the following condition holds:

\begin{equation*}
\begin{aligned}
&\mbb{E}_{\x\sim \mathcal{D}}[\mbf{Pr}_{\z\sim \mathcal{D}}(f(\z) = f(\x) \mid \z_{S} = \x_{S})] \\
&= \mbf{Pr}(x_1 = 0, x_2 = 0) \cdot \mbf{Pr}_{\z\sim \mathcal{D}}(f(\z) = 0 \mid x_2 = 0) \ + \\
&\hspace{1em} \mbf{Pr}(x_1 = 0, x_2 = 1) \cdot \mbf{Pr}_{\z\sim \mathcal{D}}(f(\z) = 0 \mid x_2 = 1) \ + \\
&\hspace{1em} \mbf{Pr}(x_1 = 1, x_2 = 0) \cdot \mbf{Pr}_{\z\sim \mathcal{D}}(f(\z) = 1 \mid x_2 = 0) \ + \\
&\hspace{1em} \mbf{Pr}(x_1 = 1, x_2 = 1) \cdot \mbf{Pr}_{\z\sim \mathcal{D}}(f(\z) = 1 \mid x_2 = 1) \\
\end{aligned}
\end{equation*}

This can also be rearranged as follows:

\begin{equation*}
\begin{aligned}
& \mbf{Pr}(x_1 = 0, x_2 = 0) \cdot \frac{\mbf{Pr}(x_1 = 0, x_2 = 0)}{\mbf{Pr}(x_1 = 0, x_2 = 0) + \mbf{Pr}(x_1 = 1, x_2 = 0)} \ + \\
&\hspace{1em} \mbf{Pr}(x_1 = 0, x_2 = 1) \cdot \frac{\mbf{Pr}(x_1 = 0, x_2 = 1)}{\mbf{Pr}(x_1 = 0, x_2 = 1) + \mbf{Pr}(x_1 = 1, x_2 = 1)} \ + \\
&\hspace{1em} \mbf{Pr}(x_1 = 1, x_2 = 0) \cdot \frac{\mbf{Pr}(x_1 = 1, x_2 = 0)}{\mbf{Pr}(x_1 = 0, x_2 = 0) + \mbf{Pr}(x_1 = 1, x_2 = 0)} \ + \\
&\hspace{1em} \mbf{Pr}(x_1 = 1, x_2 = 1) \cdot \frac{\mbf{Pr}(x_1 = 1, x_2 = 1)}{\mbf{Pr}(x_1 = 0, x_2 = 1) + \mbf{Pr}(x_1 = 1, x_2 = 1)} \\
%
&= \frac{\mbf{Pr}(x_1 = 0, x_2 = 0)^2 + \mbf{Pr}(x_1 = 1, x_2 = 0)^2}{\mbf{Pr}(x_1 = 0, x_2 = 0) + \mbf{Pr}(x_1 = 1, x_2 = 0)} + 
\frac{\mbf{Pr}(x_1 = 0, x_2 = 1)^2 + \mbf{Pr}(x_1 = 1, x_2 = 1)^2}{\mbf{Pr}(x_1 = 0, x_2 = 1) + \mbf{Pr}(x_1 = 1, x_2 = 1)}
%
\end{aligned}
\end{equation*}


In contrast, when $S = \{1\}$, we have the the following holds:

\begin{equation*}
\begin{aligned}
&\mbb{E}_{\x\sim \mathcal{D}}[\mbf{Pr}_{\z\sim \mathcal{D}}(f(\z) = f(\x) \mid \z_{S} = \x_{S})] \\
&= \mbf{Pr}(x_1 = 0, x_2 = 0) \cdot \mbf{Pr}_{\z\sim \mathcal{D}}(f(\z) = 0 \mid x_1 = 0) \ + \\
&\hspace{1em} \mbf{Pr}(x_1 = 0, x_2 = 1) \cdot \mbf{Pr}_{\z\sim \mathcal{D}}(f(\z) = 0 \mid x_1 = 0) \ + \\
&\hspace{1em} \mbf{Pr}(x_1 = 1, x_2 = 0) \cdot \mbf{Pr}_{\z\sim \mathcal{D}}(f(\z) = 1 \mid x_1 = 1) \ + \\
&\hspace{1em} \mbf{Pr}(x_1 = 1, x_2 = 1) \cdot \mbf{Pr}_{\z\sim \mathcal{D}}(f(\z) = 1 \mid x_1 = 1) \\
\end{aligned}
\end{equation*}

This can similarly be rearranged as follows:

\begin{equation*}
\begin{aligned}
& \mbf{Pr}(x_1 = 0, x_2 = 0) \cdot \frac{\mbf{Pr}(x_1 = 0, x_2 = 0) + \mbf{Pr}(x_1 = 0, x_2 = 1)}{\mbf{Pr}(x_1 = 0, x_2 = 0) + \mbf{Pr}(x_1 = 0, x_2 = 1)} \ + \\
&\hspace{1em} \mbf{Pr}(x_1 = 0, x_2 = 1) \cdot \frac{\mbf{Pr}(x_1 = 0, x_2 = 0) + \mbf{Pr}(x_1 = 0, x_2 = 1)}{\mbf{Pr}(x_1 = 0, x_2 = 0) + \mbf{Pr}(x_1 = 0, x_2 = 1)} \ + \\
&\hspace{1em} \mbf{Pr}(x_1 = 1, x_2 = 0) \cdot \frac{\mbf{Pr}(x_1 = 1, x_2 = 0) + \mbf{Pr}(x_1 = 1, x_2 = 1)}{\mbf{Pr}(x_1 = 1, x_2 = 0) + \mbf{Pr}(x_1 = 1, x_2 = 1)} \ + \\
&\hspace{1em} \mbf{Pr}(x_1 = 1, x_2 = 1) \cdot \frac{\mbf{Pr}(x_1 = 1, x_2 = 0) + \mbf{Pr}(x_1 = 1, x_2 = 1)}{\mbf{Pr}(x_1 = 1, x_2 = 0) + \mbf{Pr}(x_1 = 1, x_2 = 1)} = 1
\end{aligned}
\end{equation*}


Finally, when $S = \{1, 2\}$, we have the following condition:

\begin{equation*}
\begin{aligned}
&\mbb{E}_{\x\sim \mathcal{D}}[\mbf{Pr}_{\z\sim \mathcal{D}}(f(\z) = f(\x) \mid \z_{S} = \x_{S})] \\
&= \mbf{Pr}(x_1 = 0, x_2 = 0) \cdot \mbf{Pr}_{\z\sim \mathcal{D}}(f(\z) = 0 \mid x_1 = 0, x_2 = 0) \ + \\
&\hspace{1em} \mbf{Pr}(x_1 = 0, x_2 = 1) \cdot \mbf{Pr}_{\z\sim \mathcal{D}}(f(\z) = 0 \mid x_1 = 0, x_2 = 1) \ + \\
&\hspace{1em} \mbf{Pr}(x_1 = 1, x_2 = 0) \cdot \mbf{Pr}_{\z\sim \mathcal{D}}(f(\z) = 1 \mid x_1 = 1, x_2 = 0) \ + \\
&\hspace{1em} \mbf{Pr}(x_1 = 1, x_2 = 1) \cdot \mbf{Pr}_{\z\sim \mathcal{D}}(f(\z) = 1 \mid x_1 = 1, x_2 = 1) 
\end{aligned}
\end{equation*}

This can likewise be reordered as follows:

\begin{equation*}
\begin{aligned}
&\mbf{Pr}(x_1 = 0, x_2 = 0) \cdot \frac{\mbf{Pr}(x_1 = 0, x_2 = 0)}{\mbf{Pr}(x_1 = 0, x_2 = 0)} + 
\mbf{Pr}(x_1 = 0, x_2 = 1) \cdot \frac{\mbf{Pr}(x_1 = 0, x_2 = 1)}{\mbf{Pr}(x_1 = 0, x_2 = 1)} \ + \\
&\hspace{1em} \mbf{Pr}(x_1 = 1, x_2 = 0) \cdot \frac{\mbf{Pr}(x_1 = 1, x_2 = 0)}{\mbf{Pr}(x_1 = 1, x_2 = 0)} + 
\mbf{Pr}(x_1 = 1, x_2 = 1) \cdot \frac{\mbf{Pr}(x_1 = 1, x_2 = 1)}{\mbf{Pr}(x_1 = 1, x_2 = 1)} \\
%
&= 1
\end{aligned}
\end{equation*}

Let $t = 2$. Evidently, we have $\Delta_{t}(\{1\},f) = 0$.
Next, let us analyse $\Delta_{t}(\emptyset,f)$.
We note that if we assume feature independence, we observe that:

\begin{equation*}
\begin{aligned}
&\bigl(\mbf{Pr}(x_1 = 0, x_2 = 0) + \mbf{Pr}(x_1 = 0, x_2 = 1)\bigr)^2 \ + \bigl(\mbf{Pr}(x_1 = 1, x_2 = 0) + \mbf{Pr}(x_1 = 1, x_2 = 1)\bigr)^2 \\
&= \mbf{Pr}(x_1 = 0)^2 + \mbf{Pr}(x_1 = 1)^2,
\end{aligned}
\end{equation*}

Additionally, another condition that holds is:

\begin{equation*}
\begin{aligned}
&\frac{\mbf{Pr}(x_1 = 0, x_2 = 0)^2 + \mbf{Pr}(x_1 = 1, x_2 = 0)^2}{\mbf{Pr}(x_1 = 0, x_2 = 0) + \mbf{Pr}(x_1 = 1, x_2 = 0)} + 
\frac{\mbf{Pr}(x_1 = 0, x_2 = 1)^2 + \mbf{Pr}(x_1 = 1, x_2 = 1)^2}{\mbf{Pr}(x_1 = 0, x_2 = 1) + \mbf{Pr}(x_1 = 1, x_2 = 1)} \\
&= \mbf{Pr}(x_1 = 0)^2 + \mbf{Pr}(x_1 = 1)^2.
\end{aligned}
\end{equation*}

This implies $\Delta_{t}(\emptyset,f) = 0$, that is, $\Delta_{t}(\emptyset,f) \le \Delta_{t}(\{1\},f)$. However, without the assumption of feature independence, we cannot guarantee $\Delta_{t}(\emptyset,f) = 0$. For example, let:

\[
\mbf{Pr}(x_1 = 0, x_2 = 0) = P_{00}, \
\mbf{Pr}(x_1 = 0, x_2 = 1) = P_{01}
\]
and 
\[
\mbf{Pr}(x_1 = 1, x_2 = 0) = P_{10}, \
\mbf{Pr}(x_1 = 1, x_2 = 1) = P_{11}
\]


Then, by simplifying $\Delta_{t}(\emptyset,f)$ we get that the following holds:
\begin{equation}
\Delta_{t}(\emptyset,f) = 
\Big[\frac{P_{00}^2 + P_{10}^2}{P_{00} + P_{10}} + \frac{P_{01}^2 + P_{11}^2}{P_{01} + P_{11}}\Bigr]
- \Big[\bigl(P_{00} + P_{01}\bigr)^2 \ + \bigl(P_{10} + P_{11}\bigr)^2\Bigr]
\end{equation}


We then let $P_{00} = 0.4$, $P_{01} = 0.2$, $P_{10} = 0.1$, and $P_{11} = 0.3$,
and get that the following holds:

\begin{equation}
\Delta_{t}(\emptyset,f) = 
\Big[\frac{(0.4)^2 + (0.1)^2}{0.5} + \frac{(0.2)^2 + (0.3)^2}{0.5}\Bigr]
- \Big[\bigl(0.4 + 0.2\bigr)^2 \ + \bigl(0.1 + 0.3\bigr)^2\Bigr]
= 0.08.
\end{equation}

This implies that $\Delta_{t}(\emptyset,f) > \Delta_{t}({1},f)$, indicating that --- unlike in the monotonicity setting where it was not required --- feature independence is a necessary condition for supermodularity.
%\end{remark}



%\section{Subset Minimal Computational Complexity Proofs}


\section{Proof of Proposition~\ref{subset_minimal_alg_conv_prop}}
\label{subset_minimal_alg_conv_prop_sec}

\subsetminimalalgconv*

\emph{Proof.} We note that this proposition follows directly from our proof that both $v^g_{\text{suff}}$ and $v^g_{\text{con}}$ are monotone non-decreasing (as established in Proposition~\ref{globalmonotonicity_prop}), along with other relevant conclusions drawn in prior work. In the remainder of this section, we elaborate on how this monotonicity property leads to the stated corollary.

The convergence of Algorithm~\ref{alg:subset-minimal} to a subset-minimal $\delta$ global sufficient or contrastive reason with respect to the value functions $v^{g}_{\text{con}}$ and $v^{g}{\text{suff}}$ follows directly from the monotonicity property, for which prior work~\citep{ignatiev2019abduction, wu2024verix, arenas2022computing} showed that this type of greedy algorithm yields a subset-minimal explanation when the underlying value function is monotone non-decreasing. The algorithm halts when, for all $i \in S$, the value $v(S \setminus \{i\})$ drops below $\delta$, ensuring that although $v(S) \geq \delta$ (a maintained invariant), removing any element causes the condition to fail. Due to monotonicity, this implies that no proper subset of $S$ satisfies the condition, guaranteeing that $S$ is subset-minimal.

We now turn to explain why Algorithm~\ref{alg:subset-minimal} does not converge to a subset-minimal solution in the local setting. We note that it is well known that the local probabilistic sufficient value function 
(and likewise the local probabilistic contrastive value function) is not monotone~\cite{arenas2022computing, izza2023computing, subercaseaux2024probabilistic, izza2024locally}. As a simple illustration, fix the uniform distribution on the two binary variables $\{x_1, x_2\}$ and define a Boolean function as follows:

\begin{equation}
f(x_1,x_2) := x_1 \lor x_2.
\end{equation}


Consider the input $\x = (0,1)$ which is classified as $1$, then:

\begin{equation}
v^{\ell}_{\text{suff}}(\emptyset) = \frac{3}{4},~~v^{\ell}_{\text{suff}}(\{1\}) = \frac{1}{2},~~v^{\ell}_{\text{suff}}(\{2\}) = 1,~~v^{\ell}_{\text{suff}}(\{1,2\}) = 1.
\end{equation}

Because $v^{\ell}_{\text{suff}}(\emptyset) > v^{\ell}_{\text{suff}}(\{1\}) < v^{\ell}_{\text{suff}}(\{1,2\})$, the function $v^{\ell}_{\text{suff}}$ is not monotone. Likewise, it can be computed that:

\begin{equation}
v^{\ell}_{\text{con}}(\emptyset) = 1,~~v^{\ell}_{\text{con}}(\{1\}) = 1,~~v^{\ell}_{\text{con}}(\{2\}) = \frac{1}{2},~~v^{\ell}_{\text{con}}(\{1,2\}) = \frac{3}{4}.
\end{equation}

(For $v^{\ell}_{\text{con}}$, the $S$ denotes the set of features that are allowed to vary.)
Because $v^{\ell}_{\text{con}}(\emptyset) > v^{\ell}_{\text{con}}(\{2\}) < v^{\ell}_{\text{con}}(\{1,2\})$, the function $v^{\ell}_{\text{con}}$ is also non-monotone.


$\qedsymbol{}$



\section{Proof of Theorem~\ref{subset_minimal_dt_theor}}
\label{subset_minimal_dt_theor_sec}

\subsetminimaldt*

\emph{Proof.} We begin by noting that for decision trees, assuming that $\mathcal{D}$ represents independent distributions, the computation of the value function:

\begin{equation}
v^g_{\text{suff}}(S)=\mathbb{E}_{\x\sim\mathcal{D}}[
\textbf{Pr}_{\z\sim\mathcal{D}}(f(\x)=f(\z) \ | \ \x_S=\z_S)]
\end{equation}


Can be carried out in polynomial time using the following procedure. Thanks to the tractability of decision trees, we can iterate over all pairs of leaf nodes, each corresponding to partial assignments $\x_S$ and $\z_{S'}$. For each such pair, if both leaf nodes yield the same prediction under $f$, we compute the corresponding term in the expected value $\mathbb{E}_{\x\sim\mathcal{D}}[
\textbf{Pr}_{\z\sim\mathcal{D}}(f(\x)=f(\z) \mid \x_S=\z_S)]$ by multiplying the respective feature-wise probabilities over the shared features of the two vectors. Under the feature independence assumption, these probabilities decompose, allowing the full expectation to be computed by summing the contributions of all such matching leaf pairs.

Since all probabilities involved are provided as part of the input and each step involves only polynomial-time operations, the entire procedure runs in polynomial time. This establishes that every step in Algorithm~\ref{alg:subset-minimal} is efficient, and thus the algorithm as a whole runs in polynomial time. Finally, combining this with Lemma~\ref{globalmonotonicity_prop}, which proves that $v^g_{\text{suff}}$ is monotone non-decreasing, we conclude that Algorithm~\ref{alg:subset-minimal} converges to a subset-minimal explanation. This is because the algorithm halts when, for all $i \in S$, the value $v(S \setminus \{i\})$ falls below $\delta$, while $v(S) \geq \delta$ is preserved throughout. By monotonicity, this guarantees that no strict subset of $S$ satisfies the condition, ensuring $S$ is indeed subset-minimal.

%can be computed in polynomial-time using the following procedure. Due to the tractable nature of decision trees we can iterate over pairs of leaf nodes which correspond to two partial assignments of vectors $\x_S$ and $\z_{S'}$. Given these partial assignments we can calculate and in case both of these leafs lead to the same prediction of $f$, the corresponding partial probability term of the expected value 5$\mathbb{E}_{\x\sim\mathcal{D}_p}[
%\textbf{Pr}_{\z\sim\mathcal{D}_p}(f(\x)=f(\z) \ | \ \x_S=\z_S)]$ can be calculated by multiplying them by the corresponding probability values for each feature in the mutual features of the two vectors. Due to the feature independence assumption, we get that all of these sums are sepreable by sum --- and hence summing all values together will give us the total expected value $\mathbb{E}_{\x\sim\mathcal{D}_p}[
%\textbf{Pr}_{\z\sim\mathcal{D}_p}(f(\x)=f(\z) \ | \ \x_S=\z_S)]$. We also note that all of these calclations are calculated over terms given as part of the input, and hence this entire process is polynomial which shows that every step in Algorithm~\ref{alg:subset-minimal} is polynomial and hence the entire algorithm is polynomial. Now, we are only left to connect this with our results in Lemma~\ref{globalmonotonicity_prop} which showed the non-decreasing monotonicity of $v^g_{\text{suff}}$, and this directly shows that Algorithm~\ref{alg:subset-minimal} converges to a subset-minimal explanation (since for all $i \in S$, the value $v(S \setminus \{i\})$ drops below $\delta$, ensuring that although $v(S) \geq \delta$ (a maintained invariant), removing any element causes the condition to fail. Due to monotonicity, this implies that no proper subset of $S$ satisfies the condition, guaranteeing that $S$ is subset-minimal).


For the remaining part of the claim, the result follows from~\citep{arenas2022computing}, which showed that, assuming P$\neq$NP, there is no polynomial-time algorithm for computing a subset-minimal $\delta$-sufficient reason for decision trees. While this was proven specifically under the uniform distribution, the hardness clearly extends to independent distributions, which include this case. Combined, these results establish the complexity separation between the local and global variants, thus concluding our proof.

%For the other part of the claim, this result was established by~\citep{arenas2022computing} which showed that assuming P$\neq$NP, there is no polynomial-time algorithm for computing a subset-minimal $\delta$ sufficient reason for decision trees (this was particualrly proven under the uniform distributions, but clearly hardness for independent distributions which encompasses this class holds too). This together, gives us the complexity seperation between the local and the global variant, thus concluding our proof.

$\qedsymbol{}$



\section{Proof of Proposition~\ref{empirical_subset_minimal_prop}}
\label{empirical_subset_minimal_prop_sec}


\empiricalsubsetminimal*


\emph{Proof.} The computation of the global sufficient probability function:

\begin{equation}
v^g_{\text{suff}}(S)=\mathbb{E}_{\x\sim\mathcal{D}}[
\textbf{Pr}_{\z\sim\mathcal{D}}(f(\x)=f(\z) \ | \ \x_S=\z_S)]
\end{equation}

or the computation of the global contrastive probability function


\begin{equation}
v^g_{\text{con}}(S)=\mathbb{E}_{\x\sim\mathcal{D}}[
\textbf{Pr}_{\z\sim\mathcal{D}}(f(\x)\neq f(\z) \ | \ \x_{\Bar{S}}=\z_{\Bar{S}})]
\end{equation}


Can be performed in polynomial time when $\mathcal{D}$ is selected from the class of empirical distributions. This is achieved by iterating over pairs of instances $\x,\z$ within the dataset and running an inference through $f(\x)$ and $f(\z)$ to compute the expected values at each step. Consequently, determining whether the probability function exceeds or falls below a given threshold $\delta$ can also be accomplished in polynomial time. Furthermore, leveraging the proof of non-decreasing monotonicity (Proposition~\ref{monotonicity_proofs_appendix}) of the global probability function, we can utilize algorithm~\ref{alg:subset-minimal} for contrastive reason search or for sufficient reason search. By executing a linear number of polynomial queries, a subset-minimal explanation (whether a subset minimal contrastive reason or a subset minimal sufficient reason) can thus be obtained in polynomial time.

$\qedsymbol{}$


\section{Proof of Theorem~\ref{seperation_subset_min_emp_theor}}
\label{seperation_subset_min_emp_theor_sec}

\seperationsubsetminemp*


\emph{Proof.} The first part of the proof follows directly from Proposition~\ref{empirical_subset_minimal_prop}, which established that this holds for any model, since both $v^g_{\text{suff}}(S)=\mathbb{E}_{\x\sim\mathcal{D}}[
\textbf{Pr}_{\z\sim\mathcal{D}}(f(\x)=f(\z) \ | \ \x_S=\z_S)]$ and $v^g_{\text{con}}(S)=\mathbb{E}_{\x\sim\mathcal{D}}[
\textbf{Pr}_{\z\sim\mathcal{D}}(f(\x)\neq f(\z) \ | \ \x_{\Bar{S}}=\z_{\Bar{S}})]$ can be computed in polynomial time from the empirical distribution. For the second part, we proceed by proving two lemmas: one establishing the intractability of computing the sufficient case, and the other addressing the contrastive case.





%The first part of the proof is a direct corollary of Proposition~\ref{empirical_subset_minimal_prop} that showed that this holds for any model due to the possible polynomial-time computation of both $v^g_{\text{suff}}(S)=\mathbb{E}_{\x\sim\mathcal{D}_p}[
%\textbf{Pr}_{\z\sim\mathcal{D}_p}(f(\x)=f(\z) \ | \ \x_S=\z_S)]$ and $v^g_{\text{con}}(S)=\mathbb{E}_{\x\sim\mathcal{D}_p}[
%\textbf{Pr}_{\z\sim\mathcal{D}_p}(f(\x)\neq f(\z) \ | \ \x_{\Bar{S}}=\z_{\Bar{S}})]$ based on empirical distributions. For the second part of the proof, we will prove the following lemma:

%\textbf{Complexity separation for decision trees.} We start our analysis with decision trees --- which is considered 

%\begin{proposition}
%Assuming that $f$ is either a neural network or a tree ensemble, and given that $\mathcal{D}_p$ is an empirical distribution, then computing a local probabilistic subset minimal sufficient reason is coNP hard, while computing a local probabilistic subset minimal contrastive reason is NP-Hard.
%\end{proposition}

%\begin{proposition}
%For any model $f$ whose inference can be obtained in polynomial time, and $\mathcal{D}_p$ is an empirical distribution, then computing a subset minimal global probabilistic sufficient or contrastive reason is in PTIME.
%\end{proposition}



\begin{lemma}
\label{subset_minimal_emp_intract_lemma}
Unless PTIME$=$ NP, then there is no polynomial time algorithm for computing a subset minimal local $\delta$-sufficient reason for either a tree ensemble or a neural network, under empirical distributions.
\end{lemma}

\emph{Proof.} We will first prove this claim for neural networks and then extend the result to tree ensembles. We will establish this claim by demonstrating that if a polynomial-time algorithm exists for computing a subset-minimal local $\delta$ sufficient reason for neural networks, then it would enable us to solve the classic NP-complete \emph{CNF-SAT} problem in polynomial time. The CNF-SAT problem is defined as follows:

\vspace{0.5em} % Add vertical space 

\noindent\fbox{%
    \parbox{\columnwidth}{%
\mysubsection{CNF-SAT}:

\textbf{Input}: A formula in conjunctive normal form (CNF): $\phi$.

\textbf{Output}: \yes{}, if there exists an assignment to the $n$ literals of $\phi$ such that $\phi$ is evaluated to True,  and \no{} otherwise
    }%
}


Our proof will also utilize the following lemma (with its proof provided in~\cite{BaMoPeSu20}):

\begin{lemma}
    \label{boolean_circuit_mlp_lemma}
    Any boolean circuit $\phi$ can be encoded into an equivalent MLP over the binary domain $\{0,1\}^n\to \{0,1\}$ in polynomial time.
\end{lemma}

\emph{Proof.} We will actually establish hardness for a simpler, more specific case, which will consequently imply hardness for the more general setting. In this case, we assume that the empirical distribution $\mathcal{D}$ consists of a single element, which, for simplicity, we take to be the zero vector $\mathbf{0}_n$. Since the sufficiency of $S$ in this scenario depends only on the local instance $\mathbf{x}$ and the single instance $\mathbf{0}_n$ in our empirical distribution, the nature of the input format (whether discrete, continuous, etc.) does not affect the result. Therefore, the hardness results hold universally across all these settings.

Similarly to $\mathbf{0}_n$, we define $\mathbf{1}_n$ as an $n$-dimensional vector consisting entirely of ones. Given an input $\phi$, we initially assign $\mathbf{1}_n$ to $\phi$, effectively setting all variables to True. If $\phi$ evaluates to True under this assignment, then a satisfying truth assignment exists. Therefore, we assume that $\phi(\mathbf{1}_n) = 0$. We now introduce the following definitions:

 %We denote $\mathbf{1}_n$ as a vector of size $n$ containing only $1$'s, and $\mathbf{0}_n$ as a vector of size $n$ containing only $0$'s.  Given some input $\phi$, we can first assign $\mathbf{1}_n$ to $\phi$ (setting all variables to True). If $\phi$ evaluates to True, then there exists a Truth assignment to $\phi$, so the reduction can construct some trivially correct encoding. We can hence assume that $\phi(\mathbf{1}_n)=0$. We now denote the following:

 \begin{equation}
 \begin{aligned}
 \phi_2: = (x_1 \wedge x_2 \wedge \ldots x_n), \\
          \phi':=\phi \vee \phi_2
 \end{aligned}
 \end{equation}

 $\phi'$, while no longer a CNF, can still be transformed into an MLP $f$ using Lemma~\ref{boolean_circuit_mlp_lemma}, ensuring that $f$ behaves equivalently to $\phi'$ over the domain $\{0,1\}^n$. Given our assumption that computing a subset-minimal local $\delta$-sufficient reason is feasible in polynomial time, we can determine one for the instance $\langle f,\x:=\mathbf{1}_n\rangle$, $\delta:=1$, noting that the empirical distribution $\mathcal{D}$ is simply defined over the single data point $\mathbf{0}_n$.

We now assert that the subset-minimal $\delta$-sufficient reason generated for $\langle f, \x \rangle$ encompasses the entire input space, i.e., $S = \{1, \ldots, n\}$, if and only if $\langle \phi \rangle \not\in$\emph{CNF-SAT}.


Let us assume that $S = \{1, \ldots, n\}$. Since $S$ is a $\delta$-sufficient reason for $\langle f, \x \rangle$, this simply means that setting the complementary set to any value maintains the prediction. Since the complementary set is $\emptyset $ in this case, this trivially holds. The fact that $S$ is subset-minimal means that any other subset $S' \subseteq S$ satisfies $v(S') < \delta = 1$. Since the probability function $\textbf{Pr}_{\z\sim \mathcal{D}}(f(\x) = f(\z) \mid \x_S = \z_S)$ is determined by a single point (the distribution contains only the point $\mathbf{0}_n$), the probability function can only take values of $1$ or $0$. Hence, we also know that $v(S') = 0$. This tells us that, aside from the subset $S = \{1, \ldots, n\}$, for any subset $S' \subseteq S$, fixing the features in $ S'$ to  $1$ and the rest to $0$ does not result in a classification outcome of $1$. Since the $\phi_2$ component within $\phi'$ is True only if all features are assigned $ 1$, this directly implies that $ \phi $ is assigned False for any of these inputs. Since we already know that $ \phi$ does not return a True answer for the vector assignment $ \mathbf{1}_n $ (as verified at the beginning), and now we have established that the same holds for all other input vectors, we conclude that $ \phi \not\in $ \emph{CNF-SAT}.

Now, suppose that $ S$ is a subset that is strictly contained within $ \{1,\ldots,n\} $. Given that $S $ is sufficient under our definitions of the distribution $ \mathcal{D}$ with $ \delta=1 $, we can apply the same reasoning as before to conclude that $ v(S)=1 $. This implies that setting the features in $ S $ to $ 1 $ while setting the remaining features to $0 $ ensures that the function $ f $ evaluates to $ 1$. Since this assignment is necessarily not a vector consisting entirely of ones, it follows that the $ \phi_2 $ component within $ \phi' $ must be False. Consequently, the $\phi$ component must be True, which implies that  $\langle \phi\rangle\in$\emph{CNF-SAT}. This completes the proof.

We have established that the following claim holds for neural networks. However, extending the proof to tree ensembles requires a minor and straightforward adaptation. To do so, we will utilize the following lemma, which has been noted in several previous works~\citep{ordyniak2024explaining, audemard2022trading}:

\begin{lemma}
    \label{boolean_circuit_mlp_lemma}
    Any CNF or DNF $\phi$ can be encoded into an equivalent random forest classifier over the binary domain $\{0,1\}^n\to \{0,1\}$ in polynomial time.
\end{lemma}

\emph{Proof.} We observe that we can apply the same process used in our proof for neural networks, where we encoded $\phi'$ into an equivalent neural network. However, $\phi'$ is no longer a valid CNF due to our construction (though encoding it into an MLP was not an issue, as any Boolean circuit can be transformed into an MLP). Nevertheless, since $\phi'$ consists of a conjunction of only two terms, we can easily represent it as an equivalent CNF:


\begin{equation}
    \phi':= (c_1 \vee \phi_2) \wedge (c_2 \vee \phi_2) \wedge \ldots (c_m \vee \phi_2)
\end{equation}

Where each $c_i$ is a disjunction of a few terms. Consequently, $\phi'$ is a valid CNF, allowing us to transform it into an equivalent random forest classifier. The reduction we outlined for MLPs applies directly to these models as well, thereby completing the proof for both model families.

%where each $c_i$ is a disjunction of a few terms. Hence, we have that $\phi'$ is a valid CNF and hence we can trasnform it to an equivalent random forest classifier and the reduction we described for MLPs will directly work for these models as well. This concludes the proof for both model families.

%Now, let us assume that $S$ is some subset that is strictly contained within $\{1,\ldots,n\}$. Since $S$ is sufficient by our definitions of the distribution $\mathcal{D}_p$, and with $\delta=1$, we can use the same reasoning as above to imply that $v(S)=1$ which means that setting the features in $S$ to $1$, and the rest to $0$ sets the function $f$ to $1$. Since we know that this assignment is necessairily not a vector of only ones, then we know that the $\phi_2$ component within $\phi'$ is necessairily False. This means that the $\phi$ component must be True and hence $\langle \phi\rangle\in$\emph{CNF}. This concludes the proof.

%Let us assume that $S = \{1, \ldots, n\}$. Since $S$ is a $\delta$-sufficient reason for $\langle f,\x\rangle$ this simply means that setting the complementry set to any value maintains the prediction. Since the complementry set is $\emptyset$ in this case, then this trivially holds. The fact that $S$ is subset-minimal, means that any other subset $S'\subseteq S$ maintains that $v(S')<\delta=1$. Since the probability function $\textbf{Pr}_{\z\sim \mathcal{D}}(f(\x)=f(\z) | \x_S=\z_S$ is determined by one single point (the distribution contains only the point $\mathbf{0}_n$, then the probability function can be only $1$ or $0$, and hence we also know that $v(S')=0$. This tells us that other than the subset $S=\{1,\ldots, n\}$, for any subset $S'\subseteq S$, then fixing the features in $S'$ to $1$, and the rest to $0$ does not set the classification outcome to $1$. Since the $\phi_2$ component within $\phi'$ is True only if all features are assigned $1$, this directly implies that $\phi$ is assigned False for any ot these inputs. Since we already know that $\phi$ does not return a True answer for the vector assignment $\mathbf{1}_n$ (we have checked that in the beginning), and now we know that the same holds for all other input vectors, then we get that $\phi\not\in$\emph{CNF-SAT}. 
 
% $\phi'$, is no longer a CNF, however, we can still use Lemma~\ref{boolean_circuit_mlp_lemma} to transform $\phi'$ into an MLP $f$ which behaves equivalently to $\phi'$ under the domain $\{0,1\}^n$. We then can use the fact that we assumed that computing a local $\delta$-sufficient reason can be computed in polynomial time, and compute one for the instance $\langle f,\x:=\mathbf{1}_n\rangle$, $\delta:=1$ where we recall that the empirical distribution $\mathcal{
%D}_p$ is simply defined over the single data point $\mathbf{0}_n$.

% We first note that $\phi_2(\x) = 1$, and hence both $\phi'(\x) = 1$ and $f(\x) = 1$. If $\langle \phi \rangle \in$ \emph{CNF-SAT}, then there exists an assignment of variables that evaluates $\phi$ to True. This means there is some assignment in $\{0,1\}^n$ that evaluates $f$ to True. Since we have assumed that this is not the assignment $\x = \mathbf{1}_n$, it must be some assignment $\x' \in \{0,1\}^n \neq \x$. This implies that there exists an assignment of size $k < n$ that evaluates $\mathbf{1}_n$ to True. Consequently, there exists a subset $S$ where $|S| \leq k$ and $f(\mathbf{1}_S; \mathbf{0}_{\Bar{S}}) = f(\x_S; \x'_{\Bar{S}}) = 1 = f(\mathbf{1}_n)$.

%We first note that $\phi_2(\x)=1$ and hence both $\phi'(\x)=1$ and $f(\x)=1$. If $\langle \phi\rangle\in$ \emph{CNF-SAT} then there exists some assignment to variables which evaluates $\phi$ to True. This implies that there exists some assignment over $\{0,1\}^n$ which evaluates $f$ to true. We have assumed that this is not the assignment $\x=\mathbf{1}_n$, and hence it is some assignment $\x'\in\{0,1\}^n\neq \x$. This implies that there exists some assignment to $\mathbf{1}_n$ of size $k<n$ or which evaluates to true, and hence, there exists a subset $S$ for which $|S|\leq k$ and $f(\mathbf{1}_S;\mathbf{0}_{\Bar{S}})=f(\x_S;\x'_{\Bar{S}})=1=f(\mathbf{1}_n)$.

%For the second direction, we assume that $\langle \phi \rangle \not\in$ \emph{CNF-SAT}. This implies that no assignment evaluates $\phi$ to True. Consequently, there is no assignment with fewer than $n$ variables that evaluates $\phi_2 = (x_1 \wedge x_2 \wedge \ldots \wedge x_n)$ to True. Therefore, there is no assignment of size $k = n-1$ or less that evaluates either $\phi$ or $\phi_2$ to True. From this, we conclude that there is no subset $S \subseteq [n]$ of size $k$ or less for which $f(\x_S; \x'_{\Bar{S}}) \neq f(\x) = f(\mathbf{1}_n) = 1$, thus completing the reduction.



%For the second direction, we assume that $\langle \phi\rangle\not\in$ \emph{CNF-SAT}. This implies that no assignment evaluates to True in $\phi$. There clearly does not exist any assignment with strictly less than $n$ assignments that evaluate $\phi_2=(x_1 \wedge x_2 \wedge \ldots x_n)$ to True. Hence, we get that there is no assignment of size $k=n-1$ or less that evaluates either $\phi$ or $\phi_2$ to true. From here, we get that there is no subset $S\in[n]$ of size $k$ or less for which $f(\x_S;\x'_{\Bar{S}}) \neq f(\x)=f(\mathbf{1}_n)=1$, which concludes the reduction.

$\qedsymbol{}$


\begin{lemma}
\label{intract_subset_con_lemma}
Unless PTIME$=$ NP, then there is no polynomial time algorithm for computing a subset minimal local $\delta$-contrastive reason for either a tree ensemble or a neural network, under empirical distributions.
\end{lemma}

\emph{Proof.} We will present a proof analogous to the one in Lemma~\ref{subset_minimal_emp_intract_lemma}. Specifically, we will once again utilize the classical NP-hard CNF-SAT problem defined in  Lemma~\ref{subset_minimal_emp_intract_lemma}. In particular, given a Boolean formula $\phi$, we will demonstrate that determining a subset-minimal contrastive reason --- whether for a neural network or a tree ensemble --- allows us to decide the satisfiability of $\phi$. 


%We will perform a similar proof to the one given in Proposition~\ref{}. Particularly, we will again levarge both the classic NP-hard CNF-SAT problem defined in proposittion~\ref{} as well as both Lemma~\ref{} and Lemma~\ref{}. Particularly, given some boolean formula $\phi$ we will show that by obtaining a subset minimal contrastive reason either for a neural network or a tree ensemble, we can decide whether $\phi$ is satisfiable.

First, we check whether assigning all variables in $\phi$ to $1$ evaluates $\phi$ to True. If so, the formula is satisfiable, and we have determined its satisfiability. Otherwise, we use Lemma~\ref{boolean_circuit_mlp_lemma} to encode the CNF formula as a neural network $f$. Next, we compute a subset-minimal $\delta$-sufficient reason by setting $\mathcal{D}$ as the empirical distribution containing only a single data point $\mathbf{0}_n$, following a similar procedure to Lemma~\ref{subset_minimal_emp_intract_lemma}. Additionally, we set $\delta=1$ and compute a subset-minimal $\delta$-contrastive reason concerning $\langle f,\x\rangle$.

%First, we will check whether assigning all variables in $\phi$ to $1$ evaluates $\phi$ to True. If this is the case, the formula is satisfiable and we have determined its satisfiability. If not, we will use Lemma~\ref{} to encode the CNF formula as a neural network $f$. We will then compute a subset minimal $\delta$ sufficient reason by taking $\mathcal{D}_p$ to be the empirical distribution containing only a single data point $\mathbf{0}_n$ (similarly to the procedure done in Proposition~\ref{}). We additinoally will set $\delta=1$ and will compute a subset minimal $\delta$ contrastive reason concerning $\langle f,\x\rangle$. 

We will now demonstrate that if any subset-minimal $\delta$ contrastive reason obtained for $\phi$ is satisfiable, then $\phi$ itself is satisfiable. Conversely, if no subset-minimal $\delta$ contrastive reason is obtained, then $\phi$ is unsatisfiable. The validity of this claim follows a reasoning similar to that provided in Lemma~\ref{subset_minimal_emp_intract_lemma}. Specifically, the term $\mathbf{Pr}_{\z\sim\mathcal{D}}(f(\x)\neq f(\z) \mid \x_{\Bar{S}}=\z_{\Bar{S}})$, where the distribution $\mathcal{D}$ considers sampling from a single datapoint, can set the probability to either 0 or 1. Furthermore, since we are searching for a $\delta=1$ contrastive reason, this is equivalent to asking whether there exists an assignment that changes the classification of $f$, which corresponds to modifying the assignment of $\phi$ from False to True. If such an assignment exists, then $\phi$ is satisfiable. However, if no subset-minimal contrastive reason exists, then no subset of features fixed to zero --- when the complementary set is set to ones --- evaluates to true. This is equivalent to stating that no assignment evaluates $f$ (and consequently $\phi$) to True, implying that $\phi$ is unsatisfiable.

%We will now show that if any subset minimal $\delta$ contrastive reason that is obtained $\phi$ is satisfiable, whereas if no subset minimal $\delta$ contrastive reason is obtained, $\phi$ is unsatisfiable. The vlaiditiy of this claim holds for a similar reason to the reasonining provided in Proposition~\ref{}. Particularly, the term $\mathbf{Pr}_{\z\in\mathcal{D}}(f(\x)\neq f(\z) | \x_{\Bar{S}}=\z_{\Bar{S}})$, where the distribution $\mathcal{D}$ considers sampling from a single datapoint can set the probability to be either 0 or 1. Moreover, since we are searching for a $\delta=1$ contrastive reason, this equivalent to asking whether there exists an assingment that changes the classification of $f$ which is equivalent to changing the assignment of $\phi$ from False to True. If there exists any assignment of this form then $\phi$ is satisfiable. However, if no subset minimal contrastive reason exists then we also know that no subset of features that we fix to zeros, when we set the complementry set to ones, is evaluated to true which is equivalent to stating that no assignment evaluates $f$ (and hence also $\phi$) to True, which means $\phi$ is unsatisfiable.

To extend the proof from neural networks to tree ensembles, we can follow the same procedure outlined in Lemma~\ref{subset_minimal_emp_intract_lemma}, encoding the CNF formula into an equivalent random forest classifier. Consequently, the proof remains valid for tree ensembles, thereby concluding the proof.


$\qedsymbol{}$


%\begin{proposition}
%Assuming that $f$ is either a neural network or a tree ensemble, and given that $\mathcal{D}_p$ is an empirical distribution, then computing a local probabilistic subset minimal sufficient reason is coNP hard, while computing a local probabilistic subset minimal contrastive reason is NP-Hard.
%\end{proposition}


%\section{Background on Empirical Distributions}




%\section{Simplifying the Probability Term}

%Let us consider a subset $S \subseteq \{1, \ldots, n\}$. For any assignment to the features in $S$ that appears in $\mathbf{D}$, we suppose there are $k$ distinct assignments. We denote these distinct partial assignments to the features in $S$ as $\y^1_S, \y^2_S, \ldots, \y^k_S$, where $\y^i_S$ represents a specific partial assignment to the features in $S$, ensuring that each partial assignment is unique.


%Now, let us assume we have some subset $S\subseteq \{1,\ldots, n\}$. We can look at any possible assignment to the features in $S$ that is contained within $\mathbf{D}$. We will assume that there is generally $k$ different assignments. Let us denote these partial assignments to the features in $S$ by $\y^1_S, \y^2_S,\ldots \y^k_S$ where $\y^i_S$ denotes a partial assignment to the features in $S$, and it holds that all partial assignments are different (distinct partial assigments).

%We can also represent the subsets of inputs in $\textbf{D}$ that correspond to the partial assignment $\y^i_S$ as $\textbf{D}_i$. Consequently, we can express $\mathbf{D}$ as the union of these subsets: $\mathbf{D} = \mathbf{D}_1 \cup \mathbf{D}_2 \cup \ldots \cup \mathbf{D}_k$, with the total number of elements in $\mathbf{D}$ given by the sum of the elements in each subset: $|\mathbf{D}| = |\mathbf{D}_1| + |\mathbf{D}_2| + \ldots + |\mathbf{D}_k|$. For clarity in our discussion, let us define $d_i := |\mathbf{D}_i|$; thus, the equation simplifies to $|\mathbf{D}| = d_1 + d_2 + \ldots + d_k$.


%Let us also denote the corresponding inputs within $\textbf{D}$ which have the partial assignment of $\y^i_S$ as $\textbf{D}_i$. Overall we have that: $\mathbf{D}=\mathbf{D}_1\cup \mathbf{D}_2\cup \ldots \cup \mathbf{D}_k$ and that $|\mathbf{D}|=|\mathbf{D}_1|+|\mathbf{D}_2|+\ldots +|\mathbf{D}_k|$. To simplify the presentation, we will also denote $d_i:=|\mathbf{D}_i|$, and hence $|\mathbf{D}|=d_1+d_2+\ldots d_k$.

%Now, we have that the following relations hold:

%\begin{equation}
%\begin{aligned}
    %\textbf{Pr}_{\x.\z\sim\mathcal{D}_p}(f(\x)=f(\z) | \x_S=\z_S)=\\
%\frac{\textbf{Pr}_{\x.\z\sim\mathcal{D}_p}((f(\x)=f(\z))\cap (\x_S=\z_S))}{\textbf{Pr}_{\x.\z\sim\mathcal{D}_p}(\x_S=\z_S)}=\\
%\frac{\bigcup_{i=1}^{k}\textbf{Pr}_{\x.\z\sim\mathcal{D}_p}((f(\x)=f(\z))\cap (\x_S=\z_S=\y_S^i))}{\bigcup_{i=1}^k\textbf{Pr}_{\x.\z\sim\mathcal{D}_p}(\x_S=\z_S=\y_S^i)}=\\
%\frac{\sum_{i=1}^{k}\textbf{Pr}_{\x.\z\sim\mathcal{D}_p}((f(\x)=f(\z))\cap (\x_S=\z_S=\y_S^i))}{\sum_{i=1}^k\textbf{Pr}_{\x.\z\sim\mathcal{D}_p}(\x_S=\z_S=\y_S^i)}
 %   \end{aligned}
%\end{equation}

%The final equality results from the definition of empirical distributions and the mutual exclusivity of the events in question.


%Where the last equality is derived due to the definition of empirical distributions, and the mutual exclusivity of the corresponding events.

%We will now break down and simplify each term one by one, beginning with the first term:


%We will now simplify each one of the terms independently. We start with the first term:


%\begin{equation}
%\begin{aligned}
%\textbf{Pr}_{\x.\z\sim\mathcal{D}_p}(\x_S=\z_S=\y_S^i)= \frac{|\mathbf{D}_i|^2}{|\mathbf{D}|^2}=\frac{({d_i})^2}{|\mathbf{D}|^2}
%    \end{aligned}
%\end{equation}

%This is because sampling $\x,\z$ from $\mathcal{D}_p$ equates to sampling pairs of vectors $\x,\z$, each from the dataset $\mathbf{D}$, resulting in a total of $|\mathbf{D}|^2$ such pairs. Additionally, according to our definition, the subset of $\mathbf{D}$ where the features in $S$ are assigned $\y_S^i$ corresponds to the subset $\mathbf{D}_i \subseteq \mathbf{D}$. In iterating over pairs, there are precisely $|\mathbf{D}_i|^2 = d_i^2$ such pairs.



%This is due to the fact that sampling $\x,\z$ from $\mathcal{D}_p$ is equivalent to sampling pairs of vectors $\x,\z$ each from the dataset $\mathbf{D}$ (there are in total $|\mathbf{D}|^2$ such pairs. Moreover from our definition the subset of $\mathbf{D}$ for which the features in $S$ are assigned $\y_S^i$ are those within the subset $\mathbf{D}_i\subseteq \mathbf{D}$. When iterating over pairs there are exactly $|\mathbf{D}_i|^2=d_i^2$ such pairs.

%Now, moving on to the second term. we have that:


%\begin{equation}
%\begin{aligned}
    %\textbf{Pr}_{\x.\z\sim\mathcal{D}_p}((f(\x)=f(\z))\cap (\x_S=\z_S=\y_S^i))= \\ \textbf{Pr}_{\x.\z\sim\mathcal{D}_p}(((f(\x)=f(\z)=1)\cap (\x_S=\z_S=\y_S^i))\cup \\ ((f(\x)=f(\z)=0)\cap (\x_S=\z_S=\y_S^i)))=\\
    %\textbf{Pr}_{\x.\z\sim\mathcal{D}_p}((f(\x)=f(\z)=1)\cap (\x_S=\z_S=\y_S^i))+\\
    %\textbf{Pr}_{\x.\z\sim\mathcal{D}_p}((f(\x)=f(\z)=0)\cap (\x_S=\z_S=\y_S^i))
%\end{aligned}
%\end{equation}

%where the last equation is a result of mutual exclusivity.


%We will introduce the following notation: Let $\mathbf{D}_i$ be the subset of vectors $\z$ in $\mathbf{D}$ that correspond to the partial assignment $\y_S^i$. We divide $\mathbf{D}_i$ into two distinct subsets: $\mathbf{D}_i^1$ and $\mathbf{D}_i^0$. The subset $\mathbf{D}_i^1$ includes all vectors $\z$ in $\mathbf{D}$ that, under the partial assignment $\y_S^i$, produce an output value of $1$, that is, $\forall \z \in \mathbf{D}_i^1, f(\z) = 1$. Conversely, $\mathbf{D}_i^0$ comprises all vectors $\z$ in $\mathbf{D}$ that, with the same partial assignment, yield an output value of $0$, that is, $\forall \z \in \mathbf{D}_i^0, f(\z) = 0$. For ease of reference, we denote $d_i^1 := |\mathbf{D}_i^1|$ and $d_i^0 := |\mathbf{D}_i^0|$.

%We will now use the following notation: we split the set $\mathbf{D}_i$ which contains all vectors $\z$ in $\mathbf{D}$ with the partial assignment $\y_S^i$ into two disjoint sets: $\mathbf{D}_1^1$, which denotes all vectors $\z$ in $\mathbf{D}$ with the partial assignment $\y_S^i$ and which get an output value of $1$ (or in other words it holds that $\forall \z\in\mathbb{D}_i^1$ it holds that $f(\z)=1$. The second set is $\mathbf{D}_1^0$, which denotes all vectors $\z$ in $\mathbf{D}$ with the partial assignment $\y_S^i$ and which get an output value of $0$ (or in other words it holds that $\forall \z\in\mathbb{D}_i^1$ it holds that $f(\z)=0$. Similarly to before, for simplicity we will denote $d_i^1:=|\mathbf{D}_i^1|$ and $d_i^0:=|\mathbf{D}_i^0|$.

%Now, we have that: 

%\begin{equation}
%    \begin{aligned}
    %\textbf{Pr}_{\x.\z\sim\mathcal{D}_p}((f(\x)=f(\z)=1)\cap (\x_S=\z_S=\y_S^i))= \frac{|\mathbf{D}_i^1|^2}{|\mathbf{D}|^2}=\frac{{d_i^1}^2}{|\mathbf{D}|^2}
    %\end{aligned}
%\end{equation}

%Following the logic previously established, sampling $\x, \z$ from the distribution $\mathcal{D}_p$ corresponds to selecting pairs of vectors $\x, \z$ from the dataset $\mathbf{D}$. Consequently, the total number of pairs available is $\mathbf{D}^2$. For those pairs where the partial assignment $\y_S^i$ is assigned $1$, the exact count is ${d_i^1}^2$.


%Since, following a similar reasoning as above sampling $\x,\z$ from the distribution $\mathcal{D}_p$ is equivalent to sampling pairs of vectors $\x,\z$ from the dataset $\mathbf{D}$. Overall, we have $\mathbf{D}^2$ total pairs, and the number of pairs that have a partial assignment $\y_S^i$ and are assigned $1$ is exactly ${d_i^1}^2$.

%Moreover, for the exact same reasoning we have that:

%\begin{equation}
%    \begin{aligned}
    %\textbf{Pr}_{\x.\z\sim\mathcal{D}_p}((f(\x)=f(\z)=0)\cap (\x_S=\z_S=\y_S^i))= \frac{|\mathbf{D}_i^0|^2}{|\mathbf{D}|^2}=\frac{{d_i^0}^2}{|\mathbf{D}|^2}
  %  \end{aligned}
%\end{equation}

%And in total, we have that:

%\begin{equation}
%\begin{aligned}
    %\textbf{Pr}_{\x.\z\sim\mathcal{D}_p}((f(\x)=f(\z))\cap (\x_S=\z_S=\y_S^i))= \\ 
    %\textbf{Pr}_{\x.\z\sim\mathcal{D}_p}((f(\x)=f(\z)=1)\cap (\x_S=\z_S=\y_S^i))+\\
    %\textbf{Pr}_{\x.\z\sim\mathcal{D}_p}((f(\x)=f(\z)=0)\cap (\x_S=\z_S=\y_S^i))=\\
    %\frac{{d_i^1}^2+{d_i^0}^2}{|\mathbf{D}|^2}=\frac{{d_i^1}^2+{(d_i-d_i^1)}^2}{|\mathbf{D}|^2}
%\end{aligned}
%\end{equation}

%Now, we get back to our original term and we have that:


%\begin{equation}
%\begin{aligned}
    %\textbf{Pr}_{\x.\z\sim\mathcal{D}_p}(f(\x)=f(\z) | \x_S=\z_S)=\\
%\frac{\sum_{i=1}^{k}\textbf{Pr}_{\x.\z\sim\mathcal{D}_p}((f(\x)=f(\z))\cap (\x_S=\z_S=\y_S^i))}{\sum_{i=1}^k\textbf{Pr}_{\x.\z\sim\mathcal{D}_p}(\x_S=\z_S=\y_S^i)}=\\
%\frac{\sum_{i=1}^k \big(\frac{{d_i^1}^2+{d_i^0}^2}{|\mathbf{D}|^2}\big)}{\sum_{i=1}^k \big(\frac{{d_i}^2}{|\mathbf{D}|^2}\big)}=\frac{\sum_{i=1}^k  {d_i^1}^2+{d_i^0}^2}{\sum_{i=1}^k {d_i}^2}=\frac{\sum_{i=1}^k  \big({d_i^1}^2+{(d_i-d_i^1)^2}\big)}{\sum_{i=1}^k {d_i}^2}
%    \end{aligned}
%\end{equation}

%We first note the simple observation that since $d_i$ is always greater than $d_i^1$, and they are both positive then it holds that for all $i$:

%\begin{equation}
%    (d_i^1)^2+(d_i-d_i^1)^2\leq d_i^2
%\end{equation}

%This is true since the following is equivalent to stating that $2d_i^1(d_i^1-d_i)\leq 0$, which trivially holds. From this we also get that:

%\begin{equation}
%    \sum_{i=1}^k  \big({d_i^1}^2+{(d_i-d_i^1)^2}\big)\leq \sum_{i=1}^k {d_i}^2
%\end{equation}

%which is aligned with the fact that the value of $\textbf{Pr}_{\x.\z\sim\mathcal{D}_p}(f(\x)=f(\z) | \x_S=\z_S)$ is always between $0$ and $1$.


%We will now prove that $\textbf{Pr}_{\x.\z\sim\mathcal{D}_p}(f(\x)=f(\z) | \x_S=\z_S)$ is monotonic with respect to the subset $S$. 

%As previously specified, the subset $S$ divides the dataset $\mathbf{D}$ into $k$ unique subsets $\mathbf{D}_1,\mathbf{D}_2,\ldots,\mathbf{D}_k$, with each subset corresponding to vectors that share a particular partial assignment of the features in $S$ (this partial assignment is represented by $\y_S^i$).

%As defined earlier, the subset $S$ splits the dataset $\mathbf{D}$ into $k$ distinct subsets $\mathbf{D}_1,\mathbf{D}_2,\ldots,\mathbf{D}_k$ each representing vectors with some different partial assignment to the features in $S$ (and the partial assignment is denoted by $\y_S^i$).

%Similarly, for the subset $S \cup {i}$, we can identify $k'$ unique subsets: $\mathbf{D}'_1, \mathbf{D}'_2, \ldots, \mathbf{D}'_{k'}$, each corresponding to different assignments of the features in $S \cup {i}$ that partition the dataset $\mathbf{D}$.


%In the same way, we can define for a subset $S\cup\{i\}$ a total of $k'$ distinct subsets: $\mathbf{D}'_1,\mathbf{D}'_2,\ldots,\mathbf{D}'_{k'}$ which split the dataset $\mathbf{D}$ into assignments to the features of $S\cup\{i\}$.

%Naturally, we observe that $k' \geq k$ because there are more potential partial assignments available to $S \cup {i}$ in $\mathbf{D}$ than to $S$ alone. Additionally, we consider the relationships among $\mathbf{D}_1, \mathbf{D}_2, \ldots, \mathbf{D}_k$ and $\mathbf{D}'_1, \mathbf{D}'_2, \ldots, \mathbf{D}'_{k'}$. For any given subset $\mathbf{D}_j$, it may either appear identically in $\mathbf{D}'_1, \mathbf{D}'_2, \ldots, \mathbf{D}'_{k'}$. This occurs when every possible vector within $\mathbf{D}_j$ assigns consistently the same value (either $0$ or $1$) to feature $i$ across all partial assignments $\y_S^j$. Consequently, in these scenarios, the values of $d_i$, as well as $d_i^1$ and $d_i^0$, remain unchanged within the subsets that divide $\mathbf{D}$ concerning the subset $S$ and those that divide it concerning $S \cup {i}$.



%It of course holds that $k'\geq k$ since there are more possible partial assignments to $S\cup \{i\}$ within $\mathbf{D}$ than to $S$. We also note that the following conditions can hold between the relationship between $\mathbf{D}_1,\mathbf{D}_2,\ldots,\mathbf{D}_k$ and: $\mathbf{D}'_1,\mathbf{D}'_2,\ldots,\mathbf{D}'_{k'}$. Given some subset $\mathbf{D}_j$ it is either the case that $\mathbf{D}_j$ appears exactly as is within $\mathbf{D}'_1,\mathbf{D}'_2,\ldots,\mathbf{D}'_{k'}$. This happens when for any possible vector within $\mathbf{D}_j$ it holds that all partial assignment $\y_S^j$ have the exact same assignment to the feature $i$ (either always $0$ or always $1$). Hence, in these cases we will have that the values of $d_i$, and hence also of $d_i^1$ and $d_i^0$ remain the same within the subsets that split the dataset $\mathbf{D}$ with respect to the subset $S$ and those that split $\mathbf{D}$ with respect to the subset $S\cup \{i\}$.

%Now, the second scenario arises when, for the partial assignment $\y_S^j$, there are vectors in $\mathbf{D}$ that correspond to $\y_S^j$ with an assignment of $0$ to feature $i$, and also vectors in $\mathbf{D}$ that correspond to $\y_S^j$ with an assignment of $1$ to feature $i$. In this case, the set $\mathbf{D}_j$ is divided into two distinct subsets contained within $\mathbf{D}'_1, \mathbf{D}'_2, \ldots, \mathbf{D}'_{k'}$, which correspond to two different assignments for the features in $S \cup {i}$: one where feature $i$ is assigned a $0$ and another where feature $i$ is assigned a $1$.


%Now, the second possible scenario happens when for the partial assignment $\y_S^^j$ there are \emph{both} vectors within $\mathbf{D}$ with partial assignments to $\y_S^j$ and with an assignment to $0$ for feature $i$, and vectors within $\mathbf{D}$ with partial assignments to $\y_S^j$ and with an assignment to $1$ for feature $i$. In this particular case the set $\mathbf{D}_j$ will be split into two distinct sets that are contained within $\mathbf{D}'_1,\mathbf{D}'_2,\ldots,\mathbf{D}'_{k'}$ corresponding to two different assignments to the features in $S\cup\{i\}$: either to the assignment to features $\y_S^j$ with a $0$ assignemtn to feature $i$, and an assignemnt to features $\y_S^j$ with a $1$ assignemnt to feature $i$.

%For clarity, we'll use $d:=d_j$ to refer to the size of subset $\mathbf{D}_j$ and denote the sizes of the two subsets into which $\mathbf{D}_j$ is divided within $\mathbf{D}$ as $d'$ and $r:=d-d'$. This ensures the sizes of both subsets sum to $d$. Retaining our prior symbols, let $d_1$ represent the count of vectors within $d_1$ classified as $1$, and $d_0$ as the count of vectors classified as $0$. Naturally, it follows that $d_0=d-d_1$. Similarly, for $d'$ and $r$, we'll use the notations $d'_1$, $r_1$ for vectors classified as $1$ and $d'_0$, $r_0$ for those classified as $0$, where $d'_0=d'-d'_1$ and $r_0=r-r_1$ hold respectively.



%For simplicity, let us denote $d:=d_j$ (corresponding to our subset $\mathbf{D}_j$), and let us denote the sizes of the two sets to which $\mathbf{D}_j$ is split to within $\mathbf{D}$ by $d'$ and $r:=d-d'$ (this of course holds since the sum of the two sets needs to add up to $d$). Moreover, let us keep our previous notation and denote $d_1$ as the number of vectors within $d_1$ that were classified as $1$, and $d_0$ as the number of assignments classified as $0$. It autonmatically holds that: $d_0=d-d_1$. Moreover we will use similar notations for $d'$ and $r$, by using the notation of $d'_1$, $r_1$ and $d'_0$, $r_0$ for $1$ and $0$ classification. It again holds that $d'_0=d'-d'_1$ and that $r_0=r-r_1$.

%In summary, we will examine how splitting the subset $\mathbf{D}_i$ into two distinct subsets within $\mathbf{D}'_i$ leads to an improvement in the calculated probability function.


%In the denominator, the portion of the sum $\sum_{i=1}^k (d_i)^2$ attributed to our subset is $(d_j)^2$, which we refer to as $d^2$. After the division, this portion becomes $(d')^2 + (d-d')^2$. Thus, the contribution of subset $\mathbf{D}_j$ to the probability function before the split is $d^2$, and following the split, it is $(d')^2 + (d-d')^2$.

%Now, let's shift our focus to the nominator. Prior to the division, the relative contribution of $\mathbf{D}_j$ was:

%\begin{equation}
%    \begin{aligned}
%        (d_1)^2+(d_0)^2=(d_1)^2+(d-d_1)^2
%    \end{aligned}
%\end{equation}

%After the split, the contribution is:

%\begin{equation}
%    \begin{aligned}
%        (d'_1)^2+(d'_0)^2+(r_1)^2+(r_0)^2= (d'_1)^2+(d'-d'_1)^2+(r_1)^2+(r-r_1)^2=\\
%        (d'_1)^2+(d'-d'_1)^2+(r_1)^2+(d-d'-r_1)^2
%    \end{aligned}
%\end{equation}


%Overall, we see that partitioning $\mathbf{D}j$ into two subsets when transitioning from the probability calculation $\textbf{Pr}_{\x.\z\sim\mathcal{D}p}(f(\x)=f(\z) | \x_S=\z_S)$ to $\textbf{Pr}_{\x.\z\sim\mathcal{D}p}(f(\x)=f(\z) | \x_{S\cup{i}}=\z_{S\cup{i}})$ yields the following ``gain'' in the numerator:


%In total, we have that splitting $\mathbf{D}_j$ into two subsets when moving from the proability calculation of $\textbf{Pr}_{\x.\z\sim\mathcal{D}_p}(f(\x)=f(\z) | \x_S=\z_S)$ to the probability calculation of $\textbf{Pr}_{\x.\z\sim\mathcal{D}_p}(f(\x)=f(\z) | \x_{S\cup\{i\}}=\z_{S\cup\{i\}})$ provides the following ``gain'' go the nominator:

%\begin{equation}
%   (d')^2+(d-d')^2-d^2
%\end{equation}

%and the following ``gain'' to the denominator:

%\begin{equation}
%    (d'_1)^2+(d'-d'_1)^2+(r_1)^2+(d-d'-r_1)^2-(d_1)^2-(d-d_1)^2
%\end{equation}


%We will now subtract the gain of the numerator from the gain of the denominator:

%\begin{equation}
%\begin{aligned}
 %   ((d'_1)^2+(d'-d'_1)^2+(r_1)^2+(d-d'-r_1)^2-(d_1)^2-(d-d_1)^2)-\\
 %   ((d')^2+(d-d')^2-d^2)
%\end{aligned}
%\end{equation}

%Let us simplify this term:

%\begin{equation}
%\begin{aligned}
%    ((d'_1)^2+(d'-d'_1)^2+(r_1)^2+(d-d'-r_1)^2-(d_1)^2-(d-d_1)^2)-\\
%    ((d')^2+(d-d')^2-d^2)=\\
%    ((d'_1)^2+(d'-d'_1)^2-(d')^2)+((r_1)^2+(d-d'-r_1)^2-(d-d')^2)-\\ ((d_1)^2+(d-d_1)^2-d^2)=\\
 %   (2(d'_1)^2-2d'_1d') + (2(r_1)^2-2r_1(d-d'))-(2(d_1)^2-2dd_1)=\\
%    2(d'_1(d'_1-d')+r_1(r_1-(d-d'))-d_1(d_1-d))
%\end{aligned}
%\end{equation}

%We also understand that $d_1=d'_1+r_1$, indicating that the total number of `1's before the split equals the number after the split. Similarly, $d_0=d'_0+r_0$ implies that the total number of zeros remains unchanged before and after the split. This implies: $d-d_1=(d'-d'_1)+(d-d'-r_1)=d-d'_1-r_1$.

%We additionally know that $d_1=d'_1+r_1$ (the total number of ``1''s before the split is equal to the number after the split). In the same manner we know that $d_0=d'_0+r_0$ (total number of zeros before and after the split is the same). This is equivalent to saying that: $d-d_1=(d'-d'_1)+(d-d'-r_1)=d-d'_1-r_1$.

%Let us denote $A:=d'_1$, $B=d'-d'_1$, $C=r_1$ and $D=(d-d')-r_1$. 

%We first notice that $A+C=d'_1+r_1=d_1$ and that also: $B+D=d'-d'_1+(d-d')-r_1=d-d'_1-r_1=d-d_1$. 


%Therefore, it is clear that the final sum we have arrived at is equivalent to:


%\begin{equation}
%\begin{aligned}
 %   2(d'_1(d'_1-d')+r_1(r_1-(d-d'))-d_1(d_1-d))=\\
%    2(A(-B)+C(-D)-(A+C)(-B-D))=2(AD+BC)
%\end{aligned}
%\end{equation}

%Alternatively, we can express the final sum as follows:


%\begin{equation}
%\begin{aligned}
%d'_1 \cdot ((d-d')-r_1) + (d'-d'_1)\cdot r_1
%\end{aligned}
%\end{equation}

%It is evident that both $r_1$ and $d'_1$, which represent quantities of vectors in the general dataset $\mathbf{D}$, are non-negative. Additionally, it is observed that $d'-d'_1=d'_0$, which also remains non-negative, represents the total count of zeros among the values in the subset corresponding to $d'$ post-split. Furthermore, the expression $(d-d')-r_1=r-r_1=r_0$ indicates the number of zeros in the $f$ part of the sum following the split. Consequently, the overall sum is non-negative, expressed as follows:


%Clearly, both $r_1$ and $d'_1$ which denote quantities of vectors within the general dataset $\mathbf{D}$ are non-negative. Moreover, we note that $d'-d'_1=d'_0$ is also non-negative since it denotes the total number of zero's within the values within the subset corresponding to the $d'$ after the split (and hence is a non-negative number). Moreover, $(d-d')-r_1=r-r_1=r_0$ denotes the nubmer of zeros within the $f$ paert of the sum after the split. Hence we have that the total sum is non-negative and is equal to:


%\begin{equation}
%\begin{aligned}
%2(d'_1 \cdot r_0 + d'_0\cdot r_1)
%\end{aligned}
%\end{equation}


%To summarize, when transitioning from the probability $\textbf{Pr}_{\x,\z\sim\mathcal{D}p}(f(\x)=f(\z) | \x_S=\z_S)$ to the probability $\textbf{Pr}_{\x,\z\sim\mathcal{D}p}(f(\x)=f(\z) | \x_{S\cup{i}}=\z_{S\cup{i}})$, the dataset $\mathbf{D}$, originally divided into $k$ subsets $\mathbf{D}_1,\ldots \mathbf{D}_k$ based on the partial assignments of $S$, is now divided into $k'$ subsets $\mathbf{D}'_1,\ldots \mathbf{D}'_k$ based on the partial assignments of $S\cup{i}$.


%To summarize, we have that when shifting from the probability $\textbf{Pr}_{\x.\z\sim\mathcal{D}_p}(f(\x)=f(\z) | \x_S=\z_S)$ to the probability calculation of $\textbf{Pr}_{\x.\z\sim\mathcal{D}_p}(f(\x)=f(\z) | \x_{S\cup\{i\}}=\z_{S\cup\{i\}})$ the dataset $\mathbf{D}$ which is split into $k$ sets $\mathbf{D}_1,\ldots \mathbf{D}_k$ corresponding to the partial assignments of $S$, is now split into $k'$ sets $\mathbf{D}'_1,\ldots \mathbf{D}'_k$ corresponding to the partial assignments of $S\cup\{i\}$. 

%Every set originally in $\mathbf{D}_1, \ldots, \mathbf{D}_k$ either remains unchanged in $\mathbf{D}'_1, \ldots, \mathbf{D}'_k$, thus not affecting the probability function, or is divided into two distinct subsets within $\mathbf{D}'_1, \ldots, \mathbf{D}'_k$. When a subset $\mathbf{D}_j$ is split into two parts, the change in the probability function due to the separation increases the numerator more than the denominator. The precise difference between the increases in the numerator and the denominator is:


%\begin{equation}
%\begin{aligned}
%2(d'_1 \cdot r_0 + d'_0\cdot r_1)
%\end{aligned}
%\end{equation}

%Thus, if we consider the total variation for any $\mathbf{D}j$ split into two subsets, we find that the overall probability function is greater in the calculation of $\textbf{Pr}_{\x,\z\sim\mathcal{D}p}(f(\x)=f(\z) | \x_{S\cup{i}}=\z_{S\cup{i}})$ compared to $\textbf{Pr}_{\x,\z\sim\mathcal{D}_p}(f(\x)=f(\z) | \x_S=\z_S)$, thereby establishing monotonicity.


%\section{Proof of Supermodularity under Empirical Distributions}

%TBD

%\section{Expanding the Proof to General Distributions}

%We have not made any assumptions about the size of the dataset $\mathbf{D}$ or the distribution of vectors within it. Therefore, we can adapt various distributions to fit this scenario. Specifically, we can set the number of partial assignments for a subset $S$ in dataset $\mathbf{D}$ to $k:=2^{n}$, ensuring it includes every possible input from $\mathbb{F}$ exactly once. This approach aligns perfectly with the uniform distribution.

%Clearly, we did not make any assumption regarding the size of the dataset $\mathbf{D}$ or the spread of vectors within it. We can hence adjust many different distributions to account at this description. Particularly, we can take the number of partial assignments for some subset $S$ in the dataset $\mathbf{D}$ to be $k:=2^{n}$ and simply contains any possible input within $\mathbb{F}$ exactly once. This will correspond exactly to the \emph{uniform} distribution.

%However, we are able to utilize significantly more expressive distributions by increasing the size of $\mathbf{D}$ as needed, thereby controlling the precise distribution of any input $\x \in \mathbb{F}$. This approach allows for the representation of any discrete distribution, provided each probability $\mathbf{Pr}(\x)$ is a rational number and can thus be represented by dividing its occurrence over the dataset $\mathbf{D}$. It is important to note that this framework does \emph{not} presume independence among features and can accommodate distributions with correlations between features.

%\todo{Here Shahaf is working on changing the proof to be with general distribution notations (not just empirical)}

%\section{Simplifying the Probability Term}

%Let us consider a subset $S \subseteq \{1, \ldots, n\}$. For any assignment to the features in $S$ that appears in $\mathbf{D}$, we suppose there are $k$ distinct assignments. We denote these distinct partial assignments to the features in $S$ as $\y^1_S, \y^2_S, \ldots, \y^k_S$, where $\y^i_S$ represents a specific partial assignment to the features in $S$, ensuring that each partial assignment is unique.


%Now, let us assume we have some subset $S\subseteq \{1,\ldots, n\}$. We can look at any possible assignment to the features in $S$ that is contained within $\mathbf{D}$. We will assume that there is generally $k$ different assignments. Let us denote these partial assignments to the features in $S$ by $\y^1_S, \y^2_S,\ldots \y^k_S$ where $\y^i_S$ denotes a partial assignment to the features in $S$, and it holds that all partial assignments are different (distinct partial assigments).

%We can also represent the subsets of inputs in $\textbf{D}$ that correspond to the partial assignment $\y^i_S$ as $\textbf{D}_i$. Consequently, we can express $\mathbf{D}$ as the union of these subsets: $\mathbf{D} = \mathbf{D}_1 \cup \mathbf{D}_2 \cup \ldots \cup \mathbf{D}_k$, with the total number of elements in $\mathbf{D}$ given by the sum of the elements in each subset: $|\mathbf{D}| = |\mathbf{D}_1| + |\mathbf{D}_2| + \ldots + |\mathbf{D}_k|$. For clarity in our discussion, let us define $d_i := |\mathbf{D}_i|$; thus, the equation simplifies to $|\mathbf{D}| = d_1 + d_2 + \ldots + d_k$.


%Let us also denote the corresponding inputs within $\textbf{D}$ which have the partial assignment of $\y^i_S$ as $\textbf{D}_i$. Overall we have that: $\mathbf{D}=\mathbf{D}_1\cup \mathbf{D}_2\cup \ldots \cup \mathbf{D}_k$ and that $|\mathbf{D}|=|\mathbf{D}_1|+|\mathbf{D}_2|+\ldots +|\mathbf{D}_k|$. To simplify the presentation, we will also denote $d_i:=|\mathbf{D}_i|$, and hence $|\mathbf{D}|=d_1+d_2+\ldots d_k$.

%Now, we have that the following relations hold:

%\begin{equation}
%\begin{aligned}
    %\textbf{Pr}_{\x.\z\sim\mathcal{D}_p}(f(\x)=f(\z) | \x_S=\z_S)=\\
%\frac{\textbf{Pr}_{\x.\z\sim\mathcal{D}_p}((f(\x)=f(\z))\cap (\x_S=\z_S))}{\textbf{Pr}_{\x.\z\sim\mathcal{D}_p}(\x_S=\z_S)}=\\
%\frac{\bigcup_{i=1}^{k}\textbf{Pr}_{\x.\z\sim\mathcal{D}_p}((f(\x)=f(\z))\cap (\x_S=\z_S=\y_S^i))}{\bigcup_{i=1}^k\textbf{Pr}_{\x.\z\sim\mathcal{D}_p}(\x_S=\z_S=\y_S^i)}=\\
%\frac{\sum_{i=1}^{k}\textbf{Pr}_{\x.\z\sim\mathcal{D}_p}((f(\x)=f(\z))\cap (\x_S=\z_S=\y_S^i))}{\sum_{i=1}^k\textbf{Pr}_{\x.\z\sim\mathcal{D}_p}(\x_S=\z_S=\y_S^i)}
%    \end{aligned}
%\end{equation}

%The final equality results from the definition of empirical distributions and the mutual exclusivity of the events in question.


%Where the last equality is derived due to the definition of empirical distributions, and the mutual exclusivity of the corresponding events.

%We will now break down and simplify each term one by one, beginning with the first term:


%We will now simplify each one of the terms independently. We start with the first term:


%\begin{equation}
%\begin{aligned}
%\textbf{Pr}_{\x.\z\sim\mathcal{D}_p}(\x_S=\z_S=\y_S^i)= (\textbf{Pr}_{\z\sim\mathcal{D}_p}(\z_S=\y_S^i))^2
%    \end{aligned}
%\end{equation}

%This is because sampling $\x,\z$ from $\mathcal{D}_p$ equates to sampling pairs of vectors $\x,\z$, each from the dataset $\mathbf{D}$, resulting in a total of $|\mathbf{D}|^2$ such pairs. Additionally, according to our definition, the subset of $\mathbf{D}$ where the features in $S$ are assigned $\y_S^i$ corresponds to the subset $\mathbf{D}_i \subseteq \mathbf{D}$. In iterating over pairs, there are precisely $|\mathbf{D}_i|^2 = d_i^2$ such pairs.



%This is due to the fact that sampling $\x,\z$ from $\mathcal{D}_p$ is equivalent to sampling pairs of vectors $\x,\z$ each from the dataset $\mathbf{D}$ (there are in total $|\mathbf{D}|^2$ such pairs. Moreover from our definition the subset of $\mathbf{D}$ for which the features in $S$ are assigned $\y_S^i$ are those within the subset $\mathbf{D}_i\subseteq \mathbf{D}$. When iterating over pairs there are exactly $|\mathbf{D}_i|^2=d_i^2$ such pairs.

%Now, moving on to the second term. we have that:


%\begin{equation}
%\begin{aligned}
    %\textbf{Pr}_{\x.\z\sim\mathcal{D}_p}((f(\x)=f(\z))\cap (\x_S=\z_S=\y_S^i))= \\ \textbf{Pr}_{\x.\z\sim\mathcal{D}_p}(((f(\x)=f(\z)=1)\cap (\x_S=\z_S=\y_S^i))\cup \\ ((f(\x)=f(\z)=0)\cap (\x_S=\z_S=\y_S^i)))=\\
    %\textbf{Pr}_{\x.\z\sim\mathcal{D}_p}((f(\x)=f(\z)=1)\cap (\x_S=\z_S=\y_S^i))+\\
    %\textbf{Pr}_{\x.\z\sim\mathcal{D}_p}((f(\x)=f(\z)=0)\cap (\x_S=\z_S=\y_S^i))
%\end{aligned}
%\end{equation}

%where the last equation is a result of mutual exclusivity.


%We will introduce the following notation: Let $\mathbf{D}_i$ be the subset of vectors $\z$ in $\mathbf{D}$ that correspond to the partial assignment $\y_S^i$. We divide $\mathbf{D}_i$ into two distinct subsets: $\mathbf{D}_i^1$ and $\mathbf{D}_i^0$. The subset $\mathbf{D}_i^1$ includes all vectors $\z$ in $\mathbf{D}$ that, under the partial assignment $\y_S^i$, produce an output value of $1$, that is, $\forall \z \in \mathbf{D}_i^1, f(\z) = 1$. Conversely, $\mathbf{D}_i^0$ comprises all vectors $\z$ in $\mathbf{D}$ that, with the same partial assignment, yield an output value of $0$, that is, $\forall \z \in \mathbf{D}_i^0, f(\z) = 0$. For ease of reference, we denote $d_i^1 := |\mathbf{D}_i^1|$ and $d_i^0 := |\mathbf{D}_i^0|$.

%We will now use the following notation: we split the set $\mathbf{D}_i$ which contains all vectors $\z$ in $\mathbf{D}$ with the partial assignment $\y_S^i$ into two disjoint sets: $\mathbf{D}_1^1$, which denotes all vectors $\z$ in $\mathbf{D}$ with the partial assignment $\y_S^i$ and which get an output value of $1$ (or in other words it holds that $\forall \z\in\mathbb{D}_i^1$ it holds that $f(\z)=1$. The second set is $\mathbf{D}_1^0$, which denotes all vectors $\z$ in $\mathbf{D}$ with the partial assignment $\y_S^i$ and which get an output value of $0$ (or in other words it holds that $\forall \z\in\mathbb{D}_i^1$ it holds that $f(\z)=0$. Similarly to before, for simplicity we will denote $d_i^1:=|\mathbf{D}_i^1|$ and $d_i^0:=|\mathbf{D}_i^0|$.

%Now, we have that: 




%\begin{equation}
    %\begin{aligned}\textbf{Pr}_{\x.\z\sim\mathcal{D}_p}((f(\x)=f(\z)=1)\cap (\x_S=\z_S=\y_S^i))= \\ (\textbf{Pr}_{\z\sim\mathcal{D}_p}((f(\z)=1)\cap (\z_S=\y_S^i)))^2
    %\end{aligned}
%\end{equation}

%Following the logic previously established, sampling $\x, \z$ from the distribution $\mathcal{D}_p$ corresponds to selecting pairs of vectors $\x, \z$ from the dataset $\mathbf{D}$. Consequently, the total number of pairs available is $\mathbf{D}^2$. For those pairs where the partial assignment $\y_S^i$ is assigned $1$, the exact count is ${d_i^1}^2$.


%Since, following a similar reasoning as above sampling $\x,\z$ from the distribution $\mathcal{D}_p$ is equivalent to sampling pairs of vectors $\x,\z$ from the dataset $\mathbf{D}$. Overall, we have $\mathbf{D}^2$ total pairs, and the number of pairs that have a partial assignment $\y_S^i$ and are assigned $1$ is exactly ${d_i^1}^2$.

%Moreover, for the exact same reasoning we have that:

%\begin{equation}
%    \begin{aligned}
    %\textbf{Pr}_{\z\sim\mathcal{D}_p}((f(\x)=f(\z)=0)\cap (\x_S=\z_S=\y_S^i))= \\
    %(\textbf{Pr}_{\z\sim\mathcal{D}_p}((f(\z)=0)\cap (\z_S=\y_S^i)))^2
    %\end{aligned}
%\end{equation}

%And in total, we have that:

%\begin{equation}
%\begin{aligned}
    %\textbf{Pr}_{\x.\z\sim\mathcal{D}_p}((f(\x)=f(\z))\cap (\x_S=\z_S=\y_S^i))= \\ 
    %\textbf{Pr}_{\x.\z\sim\mathcal{D}_p}((f(\x)=f(\z)=1)\cap (\x_S=\z_S=\y_S^i))+\\
    %\textbf{Pr}_{\x.\z\sim\mathcal{D}_p}((f(\x)=f(\z)=0)\cap (\x_S=\z_S=\y_S^i))=\\
    %(\textbf{Pr}_{\z\sim\mathcal{D}_p}((f(\z)=1)\cap (\z_S=\y_S^i)))^2 + (\textbf{Pr}_{\z\sim\mathcal{D}_p}((f(\z)=0)\cap (\z_S=\y_S^i)))^2= \\
    %(\textbf{Pr}_{\z\sim\mathcal{D}_p}((f(\z)=1)\cap (\z_S=\y_S^i)))^2 + (\textbf{Pr}_{\z\sim\mathcal{D}_p}(\z_S=\y_S^i)-\textbf{Pr}_{\z\sim\mathcal{D}_p}((f(\z)=1)\cap (\z_S=\y_S^i)))^2
%\end{aligned}
%\end{equation}

%Now, we get back to our original term and we have that:


%\begin{equation}
%\begin{aligned}
    %\textbf{Pr}_{\x.\z\sim\mathcal{D}_p}(f(\x)=f(\z) | \x_S=\z_S)=\\
%\frac{\sum_{i=1}^{k}\textbf{Pr}_{\x.\z\sim\mathcal{D}_p}((f(\x)=f(\z))\cap (\x_S=\z_S=\y_S^i))}{\sum_{i=1}^k\textbf{Pr}_{\x.\z\sim\mathcal{D}_p}(\x_S=\z_S=\y_S^i)}=\\
%\frac{\sum_{i=1}^k \big(\frac{{d_i^1}^2+{d_i^0}^2}{|\mathbf{D}|^2}\big)}{\sum_{i=1}^k \big(\frac{{d_i}^2}{|\mathbf{D}|^2}\big)}=\frac{\sum_{i=1}^k  {d_i^1}^2+{d_i^0}^2}{\sum_{i=1}^k {d_i}^2}=\frac{\sum_{i=1}^k  \big({d_i^1}^2+{(d_i-d_i^1)^2}\big)}{\sum_{i=1}^k {d_i}^2}
%    \end{aligned}
%\end{equation}

%We first note the simple observation that since $d_i$ is always greater than $d_i^1$, and they are both positive then it holds that for all $i$:

%\begin{equation}
%    (d_i^1)^2+(d_i-d_i^1)^2\leq d_i^2
%\end{equation}

%This is true since the following is equivalent to stating that $2d_i^1(d_i^1-d_i)\leq 0$, which trivially holds. From this we also get that:

%\begin{equation}
%    \sum_{i=1}^k  \big({d_i^1}^2+{(d_i-d_i^1)^2}\big)\leq \sum_{i=1}^k {d_i}^2
%\end{equation}

%which is aligned with the fact that the value of $\textbf{Pr}_{\x.\z\sim\mathcal{D}_p}(f(\x)=f(\z) | \x_S=\z_S)$ is always between $0$ and $1$.


%We will now prove that $\textbf{Pr}_{\x.\z\sim\mathcal{D}_p}(f(\x)=f(\z) | \x_S=\z_S)$ is monotonic with respect to the subset $S$. 

%As previously specified, the subset $S$ divides the dataset $\mathbf{D}$ into $k$ unique subsets $\mathbf{D}_1,\mathbf{D}_2,\ldots,\mathbf{D}_k$, with each subset corresponding to vectors that share a particular partial assignment of the features in $S$ (this partial assignment is represented by $\y_S^i$).

%As defined earlier, the subset $S$ splits the dataset $\mathbf{D}$ into $k$ distinct subsets $\mathbf{D}_1,\mathbf{D}_2,\ldots,\mathbf{D}_k$ each representing vectors with some different partial assignment to the features in $S$ (and the partial assignment is denoted by $\y_S^i$).

%Similarly, for the subset $S \cup {i}$, we can identify $k'$ unique subsets: $\mathbf{D}'_1, \mathbf{D}'_2, \ldots, \mathbf{D}'_{k'}$, each corresponding to different assignments of the features in $S \cup {i}$ that partition the dataset $\mathbf{D}$.


%In the same way, we can define for a subset $S\cup\{i\}$ a total of $k'$ distinct subsets: $\mathbf{D}'_1,\mathbf{D}'_2,\ldots,\mathbf{D}'_{k'}$ which split the dataset $\mathbf{D}$ into assignments to the features of $S\cup\{i\}$.

%Naturally, we observe that $k' \geq k$ because there are more potential partial assignments available to $S \cup {i}$ in $\mathbf{D}$ than to $S$ alone. Additionally, we consider the relationships among $\mathbf{D}_1, \mathbf{D}_2, \ldots, \mathbf{D}_k$ and $\mathbf{D}'_1, \mathbf{D}'_2, \ldots, \mathbf{D}'_{k'}$. For any given subset $\mathbf{D}_j$, it may either appear identically in $\mathbf{D}'_1, \mathbf{D}'_2, \ldots, \mathbf{D}'_{k'}$. This occurs when every possible vector within $\mathbf{D}_j$ assigns consistently the same value (either $0$ or $1$) to feature $i$ across all partial assignments $\y_S^j$. Consequently, in these scenarios, the values of $d_i$, as well as $d_i^1$ and $d_i^0$, remain unchanged within the subsets that divide $\mathbf{D}$ concerning the subset $S$ and those that divide it concerning $S \cup {i}$.



%It of course holds that $k'\geq k$ since there are more possible partial assignments to $S\cup \{i\}$ within $\mathbf{D}$ than to $S$. We also note that the following conditions can hold between the relationship between $\mathbf{D}_1,\mathbf{D}_2,\ldots,\mathbf{D}_k$ and: $\mathbf{D}'_1,\mathbf{D}'_2,\ldots,\mathbf{D}'_{k'}$. Given some subset $\mathbf{D}_j$ it is either the case that $\mathbf{D}_j$ appears exactly as is within $\mathbf{D}'_1,\mathbf{D}'_2,\ldots,\mathbf{D}'_{k'}$. This happens when for any possible vector within $\mathbf{D}_j$ it holds that all partial assignment $\y_S^j$ have the exact same assignment to the feature $i$ (either always $0$ or always $1$). Hence, in these cases we will have that the values of $d_i$, and hence also of $d_i^1$ and $d_i^0$ remain the same within the subsets that split the dataset $\mathbf{D}$ with respect to the subset $S$ and those that split $\mathbf{D}$ with respect to the subset $S\cup \{i\}$.

%Now, the second scenario arises when, for the partial assignment $\y_S^j$, there are vectors in $\mathbf{D}$ that correspond to $\y_S^j$ with an assignment of $0$ to feature $i$, and also vectors in $\mathbf{D}$ that correspond to $\y_S^j$ with an assignment of $1$ to feature $i$. In this case, the set $\mathbf{D}_j$ is divided into two distinct subsets contained within $\mathbf{D}'_1, \mathbf{D}'_2, \ldots, \mathbf{D}'_{k'}$, which correspond to two different assignments for the features in $S \cup {i}$: one where feature $i$ is assigned a $0$ and another where feature $i$ is assigned a $1$.


%Now, the second possible scenario happens when for the partial assignment $\y_S^^j$ there are \emph{both} vectors within $\mathbf{D}$ with partial assignments to $\y_S^j$ and with an assignment to $0$ for feature $i$, and vectors within $\mathbf{D}$ with partial assignments to $\y_S^j$ and with an assignment to $1$ for feature $i$. In this particular case the set $\mathbf{D}_j$ will be split into two distinct sets that are contained within $\mathbf{D}'_1,\mathbf{D}'_2,\ldots,\mathbf{D}'_{k'}$ corresponding to two different assignments to the features in $S\cup\{i\}$: either to the assignment to features $\y_S^j$ with a $0$ assignemtn to feature $i$, and an assignemnt to features $\y_S^j$ with a $1$ assignemnt to feature $i$.

%For clarity, we'll use $d:=d_j$ to refer to the size of subset $\mathbf{D}_j$ and denote the sizes of the two subsets into which $\mathbf{D}_j$ is divided within $\mathbf{D}$ as $d'$ and $r:=d-d'$. This ensures the sizes of both subsets sum to $d$. Retaining our prior symbols, let $d_1$ represent the count of vectors within $d_1$ classified as $1$, and $d_0$ as the count of vectors classified as $0$. Naturally, it follows that $d_0=d-d_1$. Similarly, for $d'$ and $r$, we'll use the notations $d'_1$, $r_1$ for vectors classified as $1$ and $d'_0$, $r_0$ for those classified as $0$, where $d'_0=d'-d'_1$ and $r_0=r-r_1$ hold respectively.



%For simplicity, let us denote $d:=d_j$ (corresponding to our subset $\mathbf{D}_j$), and let us denote the sizes of the two sets to which $\mathbf{D}_j$ is split to within $\mathbf{D}$ by $d'$ and $r:=d-d'$ (this of course holds since the sum of the two sets needs to add up to $d$). Moreover, let us keep our previous notation and denote $d_1$ as the number of vectors within $d_1$ that were classified as $1$, and $d_0$ as the number of assignments classified as $0$. It autonmatically holds that: $d_0=d-d_1$. Moreover we will use similar notations for $d'$ and $r$, by using the notation of $d'_1$, $r_1$ and $d'_0$, $r_0$ for $1$ and $0$ classification. It again holds that $d'_0=d'-d'_1$ and that $r_0=r-r_1$.

%In summary, we will examine how splitting the subset $\mathbf{D}_i$ into two distinct subsets within $\mathbf{D}'_i$ leads to an improvement in the calculated probability function.


%In the denominator, the portion of the sum $\sum_{i=1}^k (d_i)^2$ attributed to our subset is $(d_j)^2$, which we refer to as $d^2$. After the division, this portion becomes $(d')^2 + (d-d')^2$. Thus, the contribution of subset $\mathbf{D}_j$ to the probability function before the split is $d^2$, and following the split, it is $(d')^2 + (d-d')^2$.

%Now, let's shift our focus to the nominator. Prior to the division, the relative contribution of $\mathbf{D}_j$ was:

%\begin{equation}
%    \begin{aligned}
%        (d_1)^2+(d_0)^2=(d_1)^2+(d-d_1)^2
%    \end{aligned}
%\end{equation}

%After the split, the contribution is:

%\begin{equation}
%    \begin{aligned}
%        (d'_1)^2+(d'_0)^2+(r_1)^2+(r_0)^2= (d'_1)^2+(d'-d'_1)^2+(r_1)^2+(r-r_1)^2=\\
%        (d'_1)^2+(d'-d'_1)^2+(r_1)^2+(d-d'-r_1)^2
%    \end{aligned}
%\end{equation}


%Overall, we see that partitioning $\mathbf{D}j$ into two subsets when transitioning from the probability calculation $\textbf{Pr}_{\x.\z\sim\mathcal{D}p}(f(\x)=f(\z) | \x_S=\z_S)$ to $\textbf{Pr}_{\x.\z\sim\mathcal{D}p}(f(\x)=f(\z) | \x_{S\cup{i}}=\z_{S\cup{i}})$ yields the following ``gain'' in the numerator:


%In total, we have that splitting $\mathbf{D}_j$ into two subsets when moving from the proability calculation of $\textbf{Pr}_{\x.\z\sim\mathcal{D}_p}(f(\x)=f(\z) | \x_S=\z_S)$ to the probability calculation of $\textbf{Pr}_{\x.\z\sim\mathcal{D}_p}(f(\x)=f(\z) | \x_{S\cup\{i\}}=\z_{S\cup\{i\}})$ provides the following ``gain'' go the nominator:

%\begin{equation}
%   (d')^2+(d-d')^2-d^2
%\end{equation}

%and the following ``gain'' to the denominator:

%\begin{equation}
%    (d'_1)^2+(d'-d'_1)^2+(r_1)^2+(d-d'-r_1)^2-(d_1)^2-(d-d_1)^2
%\end{equation}


%We will now subtract the gain of the numerator from the gain of the denominator:

%\begin{equation}
%\begin{aligned}
%    ((d'_1)^2+(d'-d'_1)^2+(r_1)^2+(d-d'-r_1)^2-(d_1)^2-(d-d_1)^2)-\\
%    ((d')^2+(d-d')^2-d^2)
%\end{aligned}
%\end{equation}

%Let us simplify this term:

%\begin{equation}
%\begin{aligned}
%    ((d'_1)^2+(d'-d'_1)^2+(r_1)^2+(d-d'-r_1)^2-(d_1)^2-(d-d_1)^2)-\\
%    ((d')^2+(d-d')^2-d^2)=\\
%    ((d'_1)^2+(d'-d'_1)^2-(d')^2)+((r_1)^2+(d-d'-r_1)^2-(d-d')^2)-\\ ((d_1)^2+(d-d_1)^2-d^2)=\\
%    (2(d'_1)^2-2d'_1d') + (2(r_1)^2-2r_1(d-d'))-(2(d_1)^2-2dd_1)=\\
%    2(d'_1(d'_1-d')+r_1(r_1-(d-d'))-d_1(d_1-d))
%\end{aligned}
%\end{equation}

%We also understand that $d_1=d'_1+r_1$, indicating that the total number of `1's before the split equals the number after the split. Similarly, $d_0=d'_0+r_0$ implies that the total number of zeros remains unchanged before and after the split. This implies: $d-d_1=(d'-d'_1)+(d-d'-r_1)=d-d'_1-r_1$.

%We additionally know that $d_1=d'_1+r_1$ (the total number of ``1''s before the split is equal to the number after the split). In the same manner we know that $d_0=d'_0+r_0$ (total number of zeros before and after the split is the same). This is equivalent to saying that: $d-d_1=(d'-d'_1)+(d-d'-r_1)=d-d'_1-r_1$.

%Let us denote $A:=d'_1$, $B=d'-d'_1$, $C=r_1$ and $D=(d-d')-r_1$. 

%We first notice that $A+C=d'_1+r_1=d_1$ and that also: $B+D=d'-d'_1+(d-d')-r_1=d-d'_1-r_1=d-d_1$. 


%Therefore, it is clear that the final sum we have arrived at is equivalent to:


%\begin{equation}
%\begin{aligned}
%    2(d'_1(d'_1-d')+r_1(r_1-(d-d'))-d_1(d_1-d))=\\
%    2(A(-B)+C(-D)-(A+C)(-B-D))=2(AD+BC)
%\end{aligned}
%\end{equation}

%Alternatively, we can express the final sum as follows:


%\begin{equation}
%\begin{aligned}
%d'_1 \cdot ((d-d')-r_1) + (d'-d'_1)\cdot r_1
%\end{aligned}
%\end{equation}

%It is evident that both $r_1$ and $d'_1$, which represent quantities of vectors in the general dataset $\mathbf{D}$, are non-negative. Additionally, it is observed that $d'-d'_1=d'_0$, which also remains non-negative, represents the total count of zeros among the values in the subset corresponding to $d'$ post-split. Furthermore, the expression $(d-d')-r_1=r-r_1=r_0$ indicates the number of zeros in the $f$ part of the sum following the split. Consequently, the overall sum is non-negative, expressed as follows:


%Clearly, both $r_1$ and $d'_1$ which denote quantities of vectors within the general dataset $\mathbf{D}$ are non-negative. Moreover, we note that $d'-d'_1=d'_0$ is also non-negative since it denotes the total number of zero's within the values within the subset corresponding to the $d'$ after the split (and hence is a non-negative number). Moreover, $(d-d')-r_1=r-r_1=r_0$ denotes the nubmer of zeros within the $f$ paert of the sum after the split. Hence we have that the total sum is non-negative and is equal to:


%\begin{equation}
%\begin{aligned}
%2(d'_1 \cdot r_0 + d'_0\cdot r_1)
%\end{aligned}
%\end{equation}


%To summarize, when transitioning from the probability $\textbf{Pr}_{\x,\z\sim\mathcal{D}p}(f(\x)=f(\z) | \x_S=\z_S)$ to the probability $\textbf{Pr}_{\x,\z\sim\mathcal{D}p}(f(\x)=f(\z) | \x_{S\cup{i}}=\z_{S\cup{i}})$, the dataset $\mathbf{D}$, originally divided into $k$ subsets $\mathbf{D}_1,\ldots \mathbf{D}_k$ based on the partial assignments of $S$, is now divided into $k'$ subsets $\mathbf{D}'_1,\ldots \mathbf{D}'_k$ based on the partial assignments of $S\cup{i}$.


%To summarize, we have that when shifting from the probability $\textbf{Pr}_{\x.\z\sim\mathcal{D}_p}(f(\x)=f(\z) | \x_S=\z_S)$ to the probability calculation of $\textbf{Pr}_{\x.\z\sim\mathcal{D}_p}(f(\x)=f(\z) | \x_{S\cup\{i\}}=\z_{S\cup\{i\}})$ the dataset $\mathbf{D}$ which is split into $k$ sets $\mathbf{D}_1,\ldots \mathbf{D}_k$ corresponding to the partial assignments of $S$, is now split into $k'$ sets $\mathbf{D}'_1,\ldots \mathbf{D}'_k$ corresponding to the partial assignments of $S\cup\{i\}$. 

%Every set originally in $\mathbf{D}_1, \ldots, \mathbf{D}_k$ either remains unchanged in $\mathbf{D}'_1, \ldots, \mathbf{D}'_k$, thus not affecting the probability function, or is divided into two distinct subsets within $\mathbf{D}'_1, \ldots, \mathbf{D}'_k$. When a subset $\mathbf{D}_j$ is split into two parts, the change in the probability function due to the separation increases the numerator more than the denominator. The precise difference between the increases in the numerator and the denominator is:


%\begin{equation}
%\begin{aligned}
%2(d'_1 \cdot r_0 + d'_0\cdot r_1)
%\end{aligned}
%\end{equation}

%Thus, if we consider the total variation for any $\mathbf{D}j$ split into two subsets, we find that the overall probability function is greater in the calculation of $\textbf{Pr}_{\x,\z\sim\mathcal{D}p}(f(\x)=f(\z) | \x_{S\cup{i}}=\z_{S\cup{i}})$ compared to $\textbf{Pr}_{\x,\z\sim\mathcal{D}_p}(f(\x)=f(\z) | \x_S=\z_S)$, thereby establishing monotonicity.


%\section{Proof of Supermodularity under Empirical Distributions}

%TBD

%\section{Expanding the Proof to General Distributions}

%We have not made any assumptions about the size of the dataset $\mathbf{D}$ or the distribution of vectors within it. Therefore, we can adapt various distributions to fit this scenario. Specifically, we can set the number of partial assignments for a subset $S$ in dataset $\mathbf{D}$ to $k:=2^{n}$, ensuring it includes every possible input from $\mathbb{F}$ exactly once. This approach aligns perfectly with the uniform distribution.

%Clearly, we did not make any assumption regarding the size of the dataset $\mathbf{D}$ or the spread of vectors within it. We can hence adjust many different distributions to account at this description. Particularly, we can take the number of partial assignments for some subset $S$ in the dataset $\mathbf{D}$ to be $k:=2^{n}$ and simply contains any possible input within $\mathbb{F}$ exactly once. This will correspond exactly to the \emph{uniform} distribution.

%However, we are able to utilize significantly more expressive distributions by increasing the size of $\mathbf{D}$ as needed, thereby controlling the precise distribution of any input $\x \in \mathbb{F}$. This approach allows for the representation of any discrete distribution, provided each probability $\mathbf{Pr}(\x)$ is a rational number and can thus be represented by dividing its occurrence over the dataset $\mathbf{D}$. It is important to note that this framework does \emph{not} presume independence among features and can accommodate distributions with correlations between features.


%However, we can take much more expressive distributions - we can make $\mathbf{D}$ as large as we want and control the exact spread of any input $\x\in\mathbb{F}$, which practically can incorporate any discrete distribution (at least assuming each probability $\mathbf{Pr}(\x)$ is a rational number --- and hence can be described by splitting its spread over the dataset $\mathbf{D}$). We note that this description does \emph{not} assume feature independence and can take distributions where there is correlations between features.


%\section{Intractability of Global Absolute SHAP}

%The efficiency property is:

%\begin{equation}
%\begin{aligned}
%    \sum_{i=1}^{n}\phi_{f,\x}(i)=v([n])-v(\emptyset)=f(\x)-\mathbb{E}_{\z\sim\mathcal{D}_p}[f(\z)]
%\end{aligned}
%\end{equation}

%For the particular case where we have that $\mathcal{D}_p$ is the uniform distribution we have that:

%\begin{equation}
%\begin{aligned}
%    \sum_{i=1}^{n}\phi_{f,\x}(i)=v([n])-v(\emptyset)=f(\x)-\mathbb{E}_{\z\sim\mathcal{D}_p}[f(\z)]=f(\x)-\frac{\#\text{SAT}}{2^n}
%\end{aligned}
%\end{equation}

%We first prove that under uniform distributions (this in practice will hold under any independent distribution) the sum of all Global SHAP values when used under the definition of Frye et al. of calculating $\mathbb{E}_{\x\sim\mathcal{D}_p}[\phi_{f.\x}(i)]$ is always zero.

%\begin{equation}
%\begin{aligned}
 %   \mathbb{E}_{\x\sim\mathcal{D}_p}[\sum_{i=1}^{n}\phi_{f,\x}(i)]=\mathbb{E}_{\x\sim\mathcal{D}_p}[f(\x)-\frac{\#\text{SAT}}{2^n}]=\\ \mathbb{E}_{\x\sim\mathcal{D}_p}[f(\x)]-\mathbb{E}_{\x\sim\mathcal{D}_p}[\frac{\#\text{SAT}}{2^n}]=\frac{\#\text{SAT}}{2^n}-\frac{\#\text{SAT}}{2^n}=0
%\end{aligned}
%\end{equation}

%This is equivalent to saying that:

%\begin{equation}
%\begin{aligned}
%    \mathbb{E}_{\x\sim\mathcal{D}_p}[\sum_{i=1}^{n}\phi_{f,\x}(i)]=    \sum_{\x\in\{0.1\}^n}\sum_{i=1}^{n}\phi_{f,\x}(i)=0
%\end{aligned}
%\end{equation}

%We will use the following notation: if we have a sum of numbers $\sum ()$, we can look at both the sum of all of the positive parts of this sum, as well as the sum of all negative parts of this sum. We will denote $\sum ()_+$ by the positive part and $\sum ()_-$ by the negative part. From the previous implication that $\sum_{\x\in\{0.1\}^n}\sum_{i=1}^{n}\phi_{f,\x}(i)=0$ we have that:

%\begin{equation}
%\begin{aligned}
%    \sum_{\x\in\{0.1\}^n} (\sum_{i\in \{1,\ldots,n\}}(\phi_{f,\x}(i)))_+=      \sum_{\x\in\{0.1\}^n} (\sum_{i\in \{1,\ldots,n\}}(\phi_{f,\x}(i)))_-
%\end{aligned}
%\end{equation}

%which also implies that:







%\begin{equation}
%\begin{aligned}
%    \sum_{\x\in\{0,1\}^n}|\sum_{i=1}^{n}\phi_{f,\x}(i)|=\sum_{\x\in\{0,1\}^n}(S_{\x,+}(i)-S_{\x,-}(i))=\sum_{\x\in\{0,1\}^n}(2S_{\x,+}(i))
%\end{aligned}
%\end{equation}

%\begin{equation}
%\begin{aligned}
%    \sum_{\x\in\{0,1\}^n}|\sum_{i=1}^{n}\phi_{f,\x}(i)|=\sum_{\x\in\{0,1\}^n}(S_{\x,+}(i)-S_{\x,-}(i))=\sum_{\x\in\{0,1\}^n}(2S_{\x,+}(i))
%\end{aligned}
%\end{equation}

%\begin{equation}
%\begin{aligned}
%    \sum_{\x\in\{0,1\}^n}|\sum_{i=1}^{n}\phi_{f,\x}(i)|=\sum_{\x\in\{0,1\}^n}|f(\x)-\frac{\#\text{SAT}}{2^n}|= \\
 %   \#\text{SAT}(1-\frac{\#\text{SAT}}{2^n})+(2^n-\#\text{SAT})(-\frac{\#\text{SAT}}{2^n})\cdot(-1)=\frac{2\#\text{SAT}}{2^n}
%\end{aligned}
%\end{equation}


%\begin{equation}
%\begin{aligned}
 %   \sum_{\x\in\{0,1\}^n}\sum_{i=1}^{n}|\phi_{f,\x}(i)|
%\end{aligned}
%\end{equation}



\section{Proof of Theorem~\ref{approx_card_con_prop}}
\label{approx_card_con_prop_sec}

\approxcardcon*

\emph{Proof.} We divide the proof of the theorem into two lemmas, covering both the approximation guarantee and the result on the absence of a bounded approximation.


\begin{lemma}
\label{approx_card_con_lemma}
    Given a neural network or tree ensemble $f$ and an empirical distribution $\mathcal{D}$ over a fixed dataset $\text{D}$, Algorithm~\ref{alg:cardinally-minimal} yields a \emph{constant} $\mathcal{O}\left(\ln\left(\frac{v^g_{con}([n])}{\min_{i \in [n]} v^g_{con}(\{i\})}\right)\right)$-approximation, bounded by $\mathcal{O}(\ln(|\text{D}|))$, for computing a global cardinally minimal $\delta$-contrastive reason for $f$, assuming feature independence.
\end{lemma}


\emph{Proof.} We will in fact prove a stronger claim, showing that this holds for \emph{any} model, provided the trivial condition that its inference time is computable in polynomial time, along with one additional mild condition that we will detail later --- both of which apply to both our neural network and tree ensemble formalizations.

We begin by noting that, since we are working with empirical distributions, the computation of the global contrastive probability value function:

\begin{equation}
v^g_{\text{con}}(S)=\mathbb{E}_{\x\sim\mathcal{D}}[
\textbf{Pr}_{\z\sim\mathcal{D}}(f(\x)\neq f(\z) \ | \ \x_{\Bar{S}}=\z_{\Bar{S}})]
\end{equation}

Can be computed in polynomial time by iterating over all pairs $\x, \z$ in the dataset D, as previously established in Proposition~\ref{globalmonotonicity_prop}. Since Algorithm~\ref{alg:cardinally-minimal} performs only a linear number of these polynomial-time queries, its total runtime is therefore polynomial.


Regarding the approximation, the classical work by Wolsey et al.~\citep{wolsey1982analysis} established a harmonic-series-based approximation guarantee for monotone non-decreasing submodular functions $v$ with integer values. More generally, their result yields an approximation factor of $\mathcal{O}\left(\ln\left(\frac{v([n])}{\min_{i \in [n]} v(\{i\})}\right)\right)$. We showed in Proposition~\ref{globalmonotonicity_prop} that the value function $v^g_{\text{con}}$ is monotone non-decreasing, and under feature independence, Proposition~\ref{globalsubmodular_prop} establishes its submodularity. Combined with the fact that Algorithm~\ref{alg:cardinally-minimal} runs in polynomial time, this directly yields an approximation guarantee of $\mathcal{O}\left(\ln\left(\frac{v^g_{\text{con}}([n])}{\min_{i \in [n]} v^g_{\text{con}}(\{i\})}\right)\right)$.


%We have proven in proposition~\ref{monotonicity_proofs_appendix} that the value function $v^g_{\text{con}}$ is monotone non decreasing, and under feature independence in Proposition~\ref{globalsubmodular_prop} that it is submodular. Hence, the combination of both these results automatically give us that Algorithm~\ref{alg:cardinally-minimal} which we have already proven that operates in polynomial time, gives an approximation factor of $\mathcal{O}\left(\ln\left(\frac{v^g_{\text{con}}([n])}{\min_{i \in [n]} v^g_{\text{con}}(\{i\})}\right)\right)$.

%Regarding the approximation, the classical work of Wolsey et al.~\citep{wolsey1982analysis} has provided a harmonic-series based approximation for integer valued monotone non-decreasing submodular functions $v$ which more generally gives an approximation factor of $\mathcal{O}\left(\ln\left(\frac{v([n])}{\min_{i \in [n]} v(\{i\})}\right)\right)$. 


However, we must also show that the expression $\mathcal{O}\left(\ln\left(\frac{v^g_{\text{con}}([n])}{\min_{i \in [n]} v^g_{\text{con}}(\{i\})}\right)\right)$ is both finite and bounded by $\mathcal{O}(\ln(|\text{D}|))$, implying that it is effectively constant, due to the assumption of a fixed dataset $|\text{D}|$. To ensure this, we begin by confirming that the expression is well-defined and finite. We then proceed to establish the desired bounds. To do so, we introduce a preprocessing step in which we eliminate ``redundant'' elements—those that could theoretically cause the denominator $\min_{i \in [n]} v^g_{\text{con}}(\{i\})$ to be zero. We begin by formally defining what we mean by redundancy:


%However, we note that we additionally need to prove that $\mathcal{O}\left(\ln\left(\frac{v^g_{\text{con}}([n])}{\min_{i \in [n]} v^g_{\text{con}}(\{i\})}\right)\right)$ is bounded by $\mathcal{O}\left(ln(|\text{D}|\right)$ and is constant. We first need to make sure that it is even finite, and then will move on to explain the other bounds. We do this using a preprocessing phase, where we detach all ``redundant'' elements which, theoretically can cause the $\min_{i \in [n]} v^g_{\text{con}}(\{i\})$ term to be equal to $0$. In order to do that, we first define redundancy in the following way:

\begin{definition}
    Let $\mathcal{D}$ denote some empirical distribution over a dataset $\text{D}$. Then we say that some feature $i\in[n]$ is \emph{redundant} with respect to $\mathcal{D}$ if for any pair $\x,\z\in\text{D}$ it holds that $f(\x_{[n]\setminus \{i\}};\z_{\{i\}})=f(\x)$.
\end{definition}


Here, the notation $f(\x_{[n]\setminus \{i\}};\z_{\{i\}})$ indicates that all features in $[n]\setminus\{i\}$ are fixed to their values in $\x$, while feature $i$ is set to its value in $\z$. Notably, this is equivalent to defining:

%Where the notation $f(\x_{[n]\setminus \{i\}};\z_{\{i\}})$ means that we fix the features $[n]\setminus\{i\}$ to their corresponding features in $\x$ and the feature $i$ to its corresponding feature in $\z$. Particularly, we note that this definition is equivalent to defining that:

\begin{equation}
    \forall S\subseteq [n],\z\in\text{D} \quad f(\x_S;\z_{\Bar{S}})=f(\x_{S\setminus \{i\}};\z_{\Bar{S}\cup\{i\}})
\end{equation}


As before, the notation $(\x_S; \z_{\Bar{S}})$ indicates that the features in $S$ are fixed to their values in $\x$, while the features in $\overline{S}$ are fixed to their values in $\z$.
%Where, similarly to the previous notation $(\x_S;\z_{\Bar{S}})$ denotes that we fix the features in $S$ to their corresponding values in $\x$ and those in $\overline{S}$ to their corresponding values in $\z$. 
Thus, in this sense, if we ``remove'' feature $i$ from the input space and define a new function \$f'\$ over the reduced space $[n']:=[n]\setminus \{i\}$, then for any input of size $n'$, the output of $f'$ will exactly match the output of $f$ when applied to the same features $[n] \setminus \{i\}$. This removal can be done in polynomial time for both tree ensembles and neural networks: for neural networks, it involves detaching the corresponding input neuron from the network; for tree ensembles, it involves removing any splits on that feature from all decision trees. Given the empirical nature of the contrastive value function $v^g_{\text{con}}(S)$ --- as previously discussed in this proof and in Proposition~\ref{globalmonotonicity_prop} --- we can compute each $v^g_{\text{con}}(\{i\})$ for all $i \in [n]$ by iterating over all pairs $\x, \z \in \mathcal{D}$ and checking whether:


%Hence, in this sense if we ``remove'' feature $i$ from our input space and define an alternative function $f'$ over an input space $[n']:=[n]\setminus \{i\}$ then the output of any instance of size $n'$ that is passed through $f'$ will correlate exactly to any output of $f$ over the features $[n]\setminus \{i\}$. We note that this process of removing a feature from our associated model $f$ can be performed in polynomial time both for the family of tree ensembles and neural networks. For a neural network this will correspond to detaching the corresponding input neuron from all other layers. In a tree ensemble, this will correspond to removin any corresponding split corresponding to that input in each one of its associated decision trees. Since given the empircial distribution nature of $v^g_{\text{con}}(S)$ due to the exact same reasoning explained earlier in this proof as well as in Proposition~\ref{globalmonotonicity_prop}, we can compute each value $v^g_{\text{con}}(\{i\})$ for every $i\in [n]$ and check whether it is strictly redundant by iterating over all pairs $\x,\z\in$D, and checking whether:

\begin{equation}    
f(\x_{[n]\setminus \{i\}};\z_{\{i\}})\neq f(\x)
\end{equation}

If this condition holds, it indicates that feature $i$ is not redundant. Conversely, if the condition fails for all pairs considered during iteration, then $i$ is deemed redundant. Once all redundant features have been identified with respect to the empirical distribution $\mathcal{D}$, we can remove them and construct an equivalent model $f'$. As discussed above, this transformation can be performed in polynomial time for both neural networks and decision trees.

%if this condition holds, this implies that $i$ is not redundant, and if this does not hold for all pairs that we iterate over, then $i$ is identified as redundant. Finally, after identifying all redundant features with respect to the empirical distribution $\mathcal{D}_p$ --- we can remove them and construct an equivalent model $f'$ --- where, for the reasons explained above, this can be done in polynomial time both for neural networks and decision trees.


To conclude our proof, we observe that $v^g_{\text{con}}(\{i\})$ equals $0$ for some empirical distribution $\mathcal{D}$ if and only if feature $i$ is redundant with respect to $\mathcal{D}$. Therefore, once the preprocessing step removes all redundant features from $f$ and we construct the resulting function $f'$, we ensure that $\min_{i \in [n]} v^g_{\text{con}}(\{i\})>0$.


%Finally, to conclude our proof, we note that $v^g_{\text{con}}(\{i\})$ is equal to $0$ for some empirical distribution $\mathcal{D}_p$ if and only if feature $i$ is redundant with respect to $\mathcal{D}_p$. Hence, after performing the preprocessing phase of removing all redundant features from $f$ and constructing $f'$ we know that $\min_{i \in [n]} v^g_{\text{con}}(\{i\})>0$.


We now present the more precise bounds referenced in our proof. Specifically, we demonstrate that the approximation factor is $\mathcal{O}(\ln(|\text{D}|))$. Given that the empirical distribution $\text{D}$ is fixed, this yields a constant approximation. While tighter bounds may be achievable, our goal here is solely to establish that the bound is \emph{constant} --- a property that will later sharply contrast with the local explainability setting, where no bounded approximation exists.


%We will now explain the more accurate bounds we mention in our proof. We will particularly show the proof of our $\mathcal{O}(\ln(|\text{D}|))$ and since we have assumed that the empirical distribution $\text{D}$ is fixed, this gives us a constant approximation. We do note that tighter bounds can be obtained, but we only have a purpose here to show that this bound is \emph{constant} (which will stand in stark contrast to the local explainability setting we will show later that does not have any bounded approximation). 


More specifically, we know that the probability function $\textbf{Pr}_{\z\sim\mathcal{D}}(f(\x)\neq f(\z) \ | \ \x_{\Bar{S}}=\z_{\Bar{S}})$ is, by the definition of empirical distributions, at least $\frac{1}{|\text{D}|}$. Therefore, the definition of $\mathbb{E}_{\x\sim\mathcal{D}}[
\textbf{Pr}_{\z\sim\mathcal{D}}(f(\x)\neq f(\z) \ | \ \x_{\Bar{S}}=\z_{\Bar{S}})]$ (i.e., the value function $v^g_{\text{con}}$) is at least $\frac{1}{|\text{D}|^2}$. In particular, this also implies that:


%More specifically, we know that the probability function $\textbf{Pr}_{\z\sim\mathcal{D}_p}(f(\x)\neq f(\z) \ | \ \x_{\Bar{S}}=\z_{\Bar{S}})$ is, due to the definition of empirical distributions greater or equal to $\frac{1}{|\text{D}|}$. Hence the definition of $\mathbb{E}_{\x\sim\mathcal{D}_p}[
%\textbf{Pr}_{\z\sim\mathcal{D}_p}(f(\x)\neq f(\z) \ | \ \x_{\Bar{S}}=\z_{\Bar{S}})]$ (i.e., the value function $v^g_{\text{con}}$) is greater or equal than $\frac{1}{|\text{D}|^2}$. More specifically, this also shows us that:

\begin{equation}
    \min_{i \in [n]} v^g_{\text{con}}(\{i\})\geq \frac{1}{|\text{D}|^2}
\end{equation}

Consequently, we obtain that:


\begin{equation}
\ln\left(\frac{v^g_{con}([n])}{\min_{i \in [n]} v^g_{con}(\{i\})}\right)\leq 
\ln(|\text{D}|^2)=\mathcal{O}(\ln(|\text{D}|)
\end{equation}

Which concludes our proof.

$\qedsymbol{}$








\begin{lemma}
\label{no_bounded_approx_con_lemma}
    Unless PTIME=NP, there is no bounded approximation for computing a cardinally minimal local $\delta$-contrastive reason for any $\langle f,\x\rangle$ where $f$ is either a neural network or a tree ensemble, even when $|\text{D}|=1$. 
\end{lemma}

\emph{Proof.} Assume, for contradiction, that there exists a bounded approximation algorithm for computing a cardinally minimal local $\delta$-contrastive reason for some $\langle f,\x\rangle$, where $f$ is a neural network or a tree ensemble—even when $|\text{D}| = 1$. However, Lemma~\ref{intract_subset_con_lemma} establishes that even with a single baseline $\z = 0_n$ (i.e., when the entire dataset is just one instance), deciding whether a contrastive reason exists is NP-hard unless PTIME = NP. Therefore, if a polynomial-time algorithm could yield a bounded approximation for a cardinally minimal contrastive reason, it would contradict this result, as such an approximation would implicitly decide the existence of a contrastive explanation.

$\qedsymbol{}$


%Assume towards contradiction that there exists a bounded approximation for computing a cardinally minimal local $\delta$-contrastive reason for some $\langle f,\x\rangle$ where $f$ is either a neural network or a tree ensemble, even when $|\text{D}|=1$. We have proven in Lemma~\ref{intract_subset_con_lemma} that even if we take a single baseline $\z=0_n$ to be our entire dataset $|\text{D}|$, that unless PTIME=NP, there is no polynomial time algorithm for deciding whether there exists one contrastive reason with respect to $\langle f,\x\rangle$ (which in turn gave us the intractabity for obtaining a subset-minimal explanation). If we have a polynomial time algorithm that gives us a bounded approximation for a cardinally minimal sufficient reason, as long as we derive a bounded approximation we can know that there does exist a cardinally minimal contrastive explanation which stands in contrast to what we have shown in Lemma~\ref{intract_subset_con_lemma}.



\section{Proof of Theorem~\ref{approx_card_suff_prop}}
\label{approx_card_suff_prop_sec}

\approxcardsuff*

\emph{Proof.} We divide the proof into two lemmas: one establishing the approximation guarantee for the global case, and the other demonstrating the intractability of the local case.


\begin{lemma}
        Given a neural network or tree ensemble $f$ and an empirical distribution $\mathcal{D}$ over a fixed dataset $\text{D}$, Algorithm~\ref{alg:cardinally-minimal} yields a \emph{constant} $\mathcal{O}\left(\frac{1}{1-k^f}+\ln\left(\frac{v^g_{\text{suff}}([n])}{\min_{i \in [n]} v^g_{\text{suff}}(\{i\})}\right)\right)$-approximation for computing a global cardinally minimal $\delta$-sufficient reason for $f$, assuming feature independence.
\end{lemma}

\emph{Proof.} The proof will follow a similar approach to that of Lemma~\ref{approx_card_con_lemma}, where we showed that for both neural networks and tree ensembles, Algorithm~\ref{alg:cardinally-minimal} achieves an approximation factor of $\ln\left(\frac{v^g_{\text{con}}([n])}{\min_{i \in [n]} v^g_{\text{con}}(\{i\})}\right)$. After applying the preprocessing step, this ratio is guaranteed to be finite and bounded by $\mathcal{O}(\ln(|\text{D}|))$.


%The proof will follow a similar line to our proof of Lemma~\ref{approx_card_con_lemma}, where we established that for both neural networks and tree ensembles the Algorithm~\ref{alg:cardinally-minimal} yields a $\ln\left(\frac{v^g_{\text{con}}([n])}{\min_{i \in [n]} v^g_{\text{con}}(\{i\})}\right)$ bound which after the described preprocessing phase, is indeed both finite and also bounded by $\mathcal{O}(\ln(|\text{D}|))$. 

Similar to Lemma~\ref{approx_card_con_lemma}, we will again prove a stronger claim --- namely, that the result holds for \emph{any} model, assuming the trivial condition that its inference time is polynomially computable, and the additional condition concerning the removal of redundant features, as described in Lemma~\ref{approx_card_con_lemma}. As discussed there, both conditions are satisfied by our neural network and tree ensemble formalizations.


Here as well, since we are working with empirical distributions, the computation of the global sufficient probability value function:

\begin{equation}
v^g_{\text{suff}}(S)=\mathbb{E}_{\x\sim\mathcal{D}}[
\textbf{Pr}_{\z\sim\mathcal{D}}(f(\x)= f(\z) \ | \ \x_{S}=\z_{S})]
\end{equation}

can be computed in polynomial time by iterating over all pairs $\x, \z$ in the dataset $\text{D}$, as previously shown in both Lemma~\ref{approx_card_con_lemma} and Proposition~\ref{globalmonotonicity_prop}. Since Algorithm~\ref{alg:cardinally-minimal} makes only a linear number of such polynomial-time queries, its overall runtime is polynomial.


Unlike Lemma~\ref{approx_card_con_lemma}, where the approximation relied on the result by Wolsey et al.~\citep{wolsey1982analysis} for monotone non-decreasing submodular functions, the setting here requires a different condition due to the supermodular nature of the function. Specifically, Shi et al.\citep{shi2021minimum} provided an approximation guarantee of $\mathcal{O}\left(\frac{1}{1-k^f}+\ln\left(\frac{v^g_{\text{suff}}([n])}{\min_{i \in [n]} v^g_{\text{suff}}(\{i\})}\right)\right)$, where $k^f := 1 - \min_{i \in [n]} \frac{v([n]) - v([n] \setminus {i})}{v({i}) - v(\emptyset)}$. We established in Proposition~\ref{globalmonotonicity_prop} that the value function $v^g_{\text{suff}}$ is monotone non-decreasing, and under the assumption of feature independence, Proposition~\ref{globalsupermod_prop} further shows that it is supermodular. Given that Algorithm~\ref{alg:cardinally-minimal} runs in polynomial time, this leads directly to an approximation guarantee of $\mathcal{O}\left(\frac{1}{1-k^f}+\ln\left(\frac{v^g_{\text{suff}}([n])}{\min_{i \in [n]} v^g_{\text{suff}}(\{i\})}\right)\right)$, where $k^f := 1 - \min_{i \in [n]} \frac{v^g_{\text{suff}}([n]) - v^g_{\text{suff}}(([n] \setminus \{i\})}{v^g_{\text{suff}}(\{i\}) - v^g_{\text{suff}}(\emptyset)}$.



%.We showed in Proposition~\ref{globalmonotonicity_prop} that the value function $v^g_{\text{con}}$ is monotone non-decreasing, and under feature independence, Proposition~\ref{globalsupermod_prop} establishes its supermodularity. Combined with the fact that Algorithm~\ref{alg:cardinally-minimal} runs in polynomial time, this directly yields an approximation guarantee of $\mathcal{O}\left(\ln\left(\frac{v^g_{\text{suff}}([n])}{\min_{i \in [n]} v^g_{\text{suff}}(\{i\})}\right)\right)$.


%We have proven in proposition~\ref{monotonicity_proofs_appendix} that the value function $v^g_{\text{con}}$ is monotone non decreasing, and under feature independence in Proposition~\ref{globalsubmodular_prop} that it is submodular. Hence, the combination of both these results automatically give us that Algorithm~\ref{alg:cardinally-minimal} which we have already proven that operates in polynomial time, gives an approximation factor of $\mathcal{O}\left(\ln\left(\frac{v^g_{\text{con}}([n])}{\min_{i \in [n]} v^g_{\text{con}}(\{i\})}\right)\right)$.

%Regarding the approximation, the classical work of Wolsey et al.~\citep{wolsey1982analysis} has provided a harmonic-series based approximation for integer valued monotone non-decreasing submodular functions $v$ which more generally gives an approximation factor of $\mathcal{O}\left(\ln\left(\frac{v([n])}{\min_{i \in [n]} v(\{i\})}\right)\right)$. 


To show that this expression is both bounded and constant, we follow the same preprocessing step as in Lemma~\ref{approx_card_con_lemma}, where we remove all redundant features from $f$ --- a process that, as previously explained, can be carried out in polynomial time for both neural networks and tree ensembles. As before, a feature $i$ is strictly redundant if and only if $v(\{i\}) - v(\emptyset) = 0$. This preprocessing yields a new function $f'$ that behaves identically to $f$ over the remaining features and ensures that, for every $i$, both $v^g_{\text{suff}}(\{i\})>0$ and $v^g_{\text{suff}}(\{i\}) - v^g_{\text{suff}}(\emptyset)>0$ hold.


In order for us to show that this expression is both bounded and constant, similarly to Lemma~\ref{approx_card_con_lemma} we will perform the exact same preprocessing phase (which, as wel explained there, can be performed in polynomial time both for neural networks as well as tree ensembles) as before where we remove all redundant features from $f$. We note that here too, it holds that a feature $i$ is strictly redundant if and only if $v(\{i\})-v(\emptyset)=0$. Hence, this preprocessing phase will give us a new function $f'$ for which it both holds that it is equivalent to the behaviour of $f$ for all remaining feautres and also satisfies that for any $i$ it holds that both $v^g_{\text{suff}}(\{i\})>0$ and $v^g_{\text{suff}}(\{i\}) - v^g_{\text{suff}}(\emptyset)>0$ hold. 


Now, following the same reasoning as in Lemma~\ref{approx_card_con_lemma}, where the probability term is lower bounded by $\frac{1}{|\text{D}|}$ and its expected value by $\frac{1}{|\text{D}|^2}$, we obtain that $\min_{i \in [n]} v^g_{\text{suff}}({i}) \geq \frac{1}{|\text{D}|^2}$ and $v^g_{\text{suff}}({i}) - v^g_{\text{suff}}(\emptyset) \geq \frac{1}{|\text{D}|^2}$. This implies, as in Lemma~\ref{approx_card_con_lemma}, that the first term is lower bounded by $\mathcal{O}(\ln(|\text{D}|))$, and for the second term we derive the following:

%Now, due to the same reasononging provided in Lemma~\ref{approx_card_con_lemma} that a lower bound to the probability term is $\frac{1}{|\text{D}|}$ and to the expected value over the probabiltiy term is $\frac{1}{|\text{D}|^2}$, we have that $\min_{i \in [n]} v^g_{\text{suff}}(\{i\})\geq \frac{1}{|\text{D}|^2}$ and $v^g_{\text{suff}}(\{i\}) - v^g_{\text{suff}}(\emptyset)\geq \frac{1}{|\text{D}|^2}$ which gives us for the first term, as in Lemma~\ref{approx_card_con_lemma} that it is lower bounded by $\mathcal{O}(\ln(|\text{D}|))$ and for the second term we have the following: 



Since $v^g_{\text{suff}}$ is supermodular (Proposition~\ref{globalsupermod_prop}), we have that
$v^g_{\text{suff}}([n]) - v^g_{\text{suff}}(([n] \setminus \{i\})$. It follows that $1 \leq v^g_{\text{suff}}([n]) - v^g_{\text{suff}}([n] \setminus \{i\}) \leq \frac{1}{|\text{D}|^2}$. This implies: $1 \leq \min_{i \in [n]} \frac{v^g_{\text{suff}}([n]) - v^g_{\text{suff}}([n] \setminus \{i\})}{v^g_{\text{suff}}(\{i\}) - v^g_{\text{suff}}(\emptyset)} \leq |\text{D}|^2$, and therefore, $1 - |\text{D}|^2 \leq k^f \leq 0$. This yields $\frac{1}{1 - k^f} \leq |\text{D}|^2$, demonstrating that the overall approximation bound based on the one established by Shi et al. is constant. While this bound may be significantly smaller in practice, our goal here is simply to show that it remains constant --- unlike in the local setting, where no bounded approximation is achievable.


%From this, we get $\frac{1}{1 - k^f} \leq |\text{D}|^2$, which shows that the total approximation bound established by Shi et al. is constant where we note that this bound in practice can be much lower but our purpose here was just to show that it remains constant in stark contrast to the local configuration where no bounded approximation is possible. This concludes the proof.


%Since $v^g_{\text{suff}}$ is supermodular (Proposition~\ref{globalsupermod_prop}) we know that $v^g_{\text{suff}}([n]) - v^g_{\text{suff}}(([n] \setminus \{i\})$ is greater or equal than $v^g_{\text{suff}}(\{i\}) - v^g_{\text{suff}}(\emptyset)$ and hence we have that: $1\leq v^g_{\text{suff}}([n]) - v^g_{\text{suff}}(([n] \setminus \{i\})\leq \frac{1}{|\text{D}|^2}$. 

%This implies that: $1\leq \min_{i \in [n]} \frac{v^g_{\text{suff}}([n]) - v^g_{\text{suff}}(([n] \setminus \{i\})}{v^g_{\text{suff}}(\{i\}) - v^g_{\text{suff}}(\emptyset)}\leq |\text{D}|^2$ and hence $1-|\text{D}|^2\leq k^f \leq 0$. From here we get that $\frac{1}{1-k^f}\leq \text{|D|^2}$ which gives us that the total bound proposed by Shi et al. is constant, hence concluding the proof.


%However, we must also show that the expression $\mathcal{O}\leftf$ is both finite and bounded by $\mathcal{O}(\ln(|\text{D}|))$, implying that it is effectively constant, due to the assumption of a fixed dataset $|\text{D}|$. To ensure this, we begin by confirming that the expression is well-defined and finite. We then proceed to establish the desired bounds. To do so, we introduce a preprocessing step in which we eliminate ``redundant'' elements --- those that could theoretically cause the denominator $\min_{i \in [n]} v^g_{\text{con}}(\{i\})$ to be zero. We begin by formally defining what we mean by redundancy:


%However, we note that we additionally need to prove that $\mathcal{O}\left(\ln\left(\frac{v^g_{\text{con}}([n])}{\min_{i \in [n]} v^g_{\text{con}}(\{i\})}\right)\right)$ is bounded by $\mathcal{O}\left(ln(|\text{D}|\right)$ and is constant. We first need to make sure that it is even finite, and then will move on to explain the other bounds. We do this using a preprocessing phase, where we detach all ``redundant'' elements which, theoretically can cause the $\min_{i \in [n]} v^g_{\text{con}}(\{i\})$ term to be equal to $0$. In order to do that, we first define redundancy in the following way:

%\begin{definition}
%    Let $\mathcal{D}_p$ denote some empirical distribution over a dataset $\text{D}$. Then we say that some feature $i\in[n]$ is \emph{redundant} with respect to $\mathcal{D}_p$ if for any pair $\x,\z\in\text{D}$ it holds that $f(\x_{[n]\setminus \{i\}};\z_{\{i\}})=f(\x)$.
%\end{definition}


%Here, the notation $f(\x_{[n]\setminus \{i\}};\z_{\{i\}})$ indicates that all features in $[n]\setminus\{i\}$ are fixed to their values in $\x$, while feature $i$ is set to its value in $\z$. Notably, this is equivalent to defining:

%Where the notation $f(\x_{[n]\setminus \{i\}};\z_{\{i\}})$ means that we fix the features $[n]\setminus\{i\}$ to their corresponding features in $\x$ and the feature $i$ to its corresponding feature in $\z$. Particularly, we note that this definition is equivalent to defining that:

%\begin{equation}
%    \forall S\subseteq [n],\z\in\text{D} \quad f(\x_S;\z_{\Bar{S}})=f(\x_{S\setminus \{i\}};\z_{\Bar{S}\cup\{i\}})
%\end{equation}


%As before, the notation $(\x_S; \z_{\Bar{S}})$ indicates that the features in $S$ are fixed to their values in $\x$, while the features in $\overline{S}$ are fixed to their values in $\z$.
%Where, similarly to the previous notation $(\x_S;\z_{\Bar{S}})$ denotes that we fix the features in $S$ to their corresponding values in $\x$ and those in $\overline{S}$ to their corresponding values in $\z$. 
%Thus, in this sense, if we ``remove'' feature $i$ from the input space and define a new function \$f'\$ over the reduced space $[n']:=[n]\setminus \{i\}$, then for any input of size $n'$, the output of $f'$ will exactly match the output of $f$ when applied to the same features $[n] \setminus \{i\}$. This removal can be done in polynomial time for both tree ensembles and neural networks: for neural networks, it involves detaching the corresponding input neuron from the network; for tree ensembles, it involves removing any splits on that feature from all decision trees. Given the empirical nature of the contrastive value function $v^g_{\text{con}}(S)$ --- as previously discussed in this proof and in Proposition~\ref{globalmonotonicity_prop} --- we can compute each $v^g_{\text{con}}(\{i\})$ for all $i \in [n]$ by iterating over all pairs $\x, \z \in \mathcal{D}$ and checking whether:


%Hence, in this sense if we ``remove'' feature $i$ from our input space and define an alternative function $f'$ over an input space $[n']:=[n]\setminus \{i\}$ then the output of any instance of size $n'$ that is passed through $f'$ will correlate exactly to any output of $f$ over the features $[n]\setminus \{i\}$. We note that this process of removing a feature from our associated model $f$ can be performed in polynomial time both for the family of tree ensembles and neural networks. For a neural network this will correspond to detaching the corresponding input neuron from all other layers. In a tree ensemble, this will correspond to removin any corresponding split corresponding to that input in each one of its associated decision trees. Since given the empircial distribution nature of $v^g_{\text{con}}(S)$ due to the exact same reasoning explained earlier in this proof as well as in Proposition~\ref{globalmonotonicity_prop}, we can compute each value $v^g_{\text{con}}(\{i\})$ for every $i\in [n]$ and check whether it is strictly redundant by iterating over all pairs $\x,\z\in$D, and checking whether:

%\begin{equation}    
%f(\x_{[n]\setminus \{i\}};\z_{\{i\}})\neq f(\x)
%\end{equation}

%If this condition holds, it indicates that feature $i$ is not redundant. Conversely, if the condition fails for all pairs considered during iteration, then $i$ is deemed redundant. Once all redundant features have been identified with respect to the empirical distribution $\mathcal{D}_p$, we can remove them and construct an equivalent model $f'$. As discussed above, this transformation can be performed in polynomial time for both neural networks and decision trees.

%if this condition holds, this implies that $i$ is not redundant, and if this does not hold for all pairs that we iterate over, then $i$ is identified as redundant. Finally, after identifying all redundant features with respect to the empirical distribution $\mathcal{D}_p$ --- we can remove them and construct an equivalent model $f'$ --- where, for the reasons explained above, this can be done in polynomial time both for neural networks and decision trees.


%To conclude our proof, we observe that $v^g_{\text{con}}(\{i\})$ equals $0$ for some empirical distribution $\mathcal{D}_p$ if and only if feature $i$ is redundant with respect to $\mathcal{D}_p$. Therefore, once the preprocessing step removes all redundant features from $f$ and we construct the resulting function $f'$, we ensure that $\min_{i \in [n]} v^g_{\text{con}}(\{i\})>0$.


%Finally, to conclude our proof, we note that $v^g_{\text{con}}(\{i\})$ is equal to $0$ for some empirical distribution $\mathcal{D}_p$ if and only if feature $i$ is redundant with respect to $\mathcal{D}_p$. Hence, after performing the preprocessing phase of removing all redundant features from $f$ and constructing $f'$ we know that $\min_{i \in [n]} v^g_{\text{con}}(\{i\})>0$.


%We now present the more precise bounds referenced in our proof. Specifically, we demonstrate that the approximation factor is $\mathcal{O}(\ln(|\text{D}|))$. Given that the empirical distribution $\text{D}$ is fixed, this yields a constant approximation. While tighter bounds may be achievable, our goal here is solely to establish that the bound is \emph{constant} --- a property that will later sharply contrast with the local explainability setting, where no bounded approximation exists.


%We will now explain the more accurate bounds we mention in our proof. We will particularly show the proof of our $\mathcal{O}(\ln(|\text{D}|))$ and since we have assumed that the empirical distribution $\text{D}$ is fixed, this gives us a constant approximation. We do note that tighter bounds can be obtained, but we only have a purpose here to show that this bound is \emph{constant} (which will stand in stark contrast to the local explainability setting we will show later that does not have any bounded approximation). 


%More specifically, we know that the probability function $\textbf{Pr}_{\z\sim\mathcal{D}_p}(f(\x)\neq f(\z) \ | \ \x_{\Bar{S}}=\z_{\Bar{S}})$ is, by the definition of empirical distributions, at least $\frac{1}{|\text{D}|}$. Therefore, the definition of $\mathbb{E}_{\x\sim\mathcal{D}_p}[
%\textbf{Pr}_{\z\sim\mathcal{D}_p}(f(\x)\neq f(\z) \ | \ \x_{\Bar{S}}=\z_{\Bar{S}})]$ (i.e., the value function $v^g_{\text{con}}$) is at least $\frac{1}{|\text{D}|^2}$. In particular, this also implies that:


%More specifically, we know that the probability function $\textbf{Pr}_{\z\sim\mathcal{D}_p}(f(\x)\neq f(\z) \ | \ \x_{\Bar{S}}=\z_{\Bar{S}})$ is, due to the definition of empirical distributions greater or equal to $\frac{1}{|\text{D}|}$. Hence the definition of $\mathbb{E}_{\x\sim\mathcal{D}_p}[
%\textbf{Pr}_{\z\sim\mathcal{D}_p}(f(\x)\neq f(\z) \ | \ \x_{\Bar{S}}=\z_{\Bar{S}})]$ (i.e., the value function $v^g_{\text{con}}$) is greater or equal than $\frac{1}{|\text{D}|^2}$. More specifically, this also shows us that:

%\begin{equation}
%    \min_{i \in [n]} v^g_{\text{con}}(\{i\})\geq \frac{1}{|\text{D}|^2}
%\end{equation}

%Consequently, we obtain that:


%\begin{equation}
%\ln\left(\frac{v^g_{con}([n])}{\min_{i \in [n]} v^g_{con}(\{i\})}\right)\leq 
%\ln(|\text{D}|^2)=\mathcal{O}(\ln(|\text{D}|)
%\end{equation}

%which concludes our proof.

$\qedsymbol{}$



%We indeed perform the same preprocessing phase described in Lemma~\ref{approx_card_con_lemma} where we delete all redundant features in $f$ (given the redundancy notation given in Lemma~\ref{approx_card_con_lemma}) and construct an equivalent model $f'$ defined on all non-redundant features.

%We also note that similarly to the case of non-redun


\begin{lemma}
    Unless PTIME$=$NP, there is \emph{no bounded approximation} for computing a local cardinally minimal $\delta$-sufficient reason for any $\langle f,\x\rangle$, even when $|\text{D}| = 1$.
\end{lemma}

\emph{Proof.} The proof follows a similar approach to Lemma~\ref{no_bounded_approx_con_lemma}. Suppose, for contradiction, that there exists a polynomial-time algorithm that provides a bounded approximation to a cardinally minimal local $\delta$-sufficient reason for some $\langle f, \x \rangle$, where $f$ is either a neural network or a tree ensemble—even when $|\text{D}| = 1$. Yet, Lemma~\ref{subset_minimal_emp_intract_lemma} shows that even in the extreme case where the dataset consists of a single baseline $\z = 0_n$, deciding whether a sufficient reason exists is NP-hard unless PTIME = NP. Hence, the existence of such an approximation algorithm would contradict this hardness result, as it would entail the ability to decide the existence of a sufficient reason.

$\qedsymbol{}$

%\section{Approximating Value Functions and Their Limitations}
%\label{app:approx_val_limit}
%
%In the general case, both $v^{g}$ and $v^{l}$ cannot be computed exactly.
%A common approach is to use Monte Carlo (MC) approximation, but this has several limitations:
%\begin{enumerate}
%\item \textbf{Slow convergence:} achieving a small error requires a large number of samples, making the method computationally expensive. 
%\item \textbf{Sampling overhead:} depending on how the data distribution $\fml{D}_{p}$ is represented, drawing conditional samples $\z\sim\fml{D}_{p}(\cdot\mid \z_S=\x_S)$ can itself be slow or costly.
%\item \textbf{Inefficiency near rare events:} when the event $\{f(\z)=f(\x)\}$ is either very rare or almost always true, most sampled indicators take the same value (almost all 0's or almost all 1's).
%\end{enumerate}


\section{Hardness of Approximating Global Contrastive Reasons for Neural Networks using Empirical Distributions}
\label{section_hardness_of_approximation}

\begin{proposition}
    Assume a neural network $f:\mathbb{F}\to[c]$ with ReLU activations, some dataset $\mathbf{D}$, such that $\mathcal{D}$ is the empirical distribution defined over $\mathbf{D}$, and some threshold $\delta\in\mathbb{Q}$. Then unless $\text{NP}\subseteq \text{DTIME}(n^{\mathcal{O}(log(log(n))})$, there does not exist a polynomial-time $\mathcal{O}(\ln(|\mathbf{D}|)-\epsilon)$ approximation algorithm for computing a cardinally minimal global contrastive reason for $f$ concerning $\mathcal{D}$, for any fixed $\epsilon>0$. 
\end{proposition}

\emph{Proof.} The proof proceeds via a polynomial-time approximation-preserving reduction from a specific variant of Set Cover that is, similarly to regular Set Cover, provably hard to approximate within the factor given in the proposition~\citep{vazirani2001approximation}.

\noindent\fbox{%
	\parbox{\columnwidth}{%
		\mysubsection{Partial Set Cover}:
		
		\textbf{Input}: A universe $U:=\{u_1,\ldots,u_n\}$, a collection of sets $S:=\{S_1,\ldots S_m\}$ such that for all $j$ $S_j\subseteq U$, and a coverage threshold $\delta\in \{0,1,\ldots, n\}$.
		
		\textbf{Output}: 
		A minimum-size subset $C\subseteq S$ such that $\lvert \cup_{S_j \in C} S_j\rvert \geq \delta$.
	}%
}


We reduce this problem to that of computing a \emph{global}, cardinally-minimal contrastive reason $S$ for a neural network $f$ with respect to a distribution $\mathcal{D}$ defined over a dataset $\mathbf{D}$, and some threshold $\delta'$. We begin by recalling the definition of a global contrastive reason and then extend it to the setting where the explanation is evaluated over an empirical dataset $\mathbf{D}$. The definition is as follows:


%We reduce this problem to the problem of obtaining a global cardinally minimal contrastive reason $S$ concerning $f$ and some $\delta$ for neural networks. We dirst recall the definition of a global contrastive reason, and expand it to hold with respect to some empirical dataset $\mathbf{D}$:

\begin{equation}
\label{eq_con_global}
\mathbb{E}_{\x\sim\mathcal{D}}[\textbf{Pr}_{\z\sim \mathcal{D}}(f(\z) \neq f(\x) \ | \  \z_{\Bar{S}}= \x_{\Bar{S}})]\geq \delta.
\end{equation}

More specifically, when we assume an empirical distribution over $\mathbf{D}$, this implies that:


\begin{equation}
\begin{aligned}
\mathbb{E}_{\x\sim\mathcal{D}}[\textbf{Pr}_{\z\sim \mathcal{D}}(f(\z) \neq f(\x) \ | \  \z_{\Bar{S}}= \x_{\Bar{S}})]=\\
\frac{1}{\lvert \textbf{D} \rvert}\sum_{\x\in\textbf{D}}[\textbf{Pr}_{\z\sim \mathcal{D}}(f(\z) \neq f(\x) \ | \  \z_{\Bar{S}}= \x_{\Bar{S}})]=\\
\frac{1}{\lvert \textbf{D} \rvert}\sum_{\x\in\textbf{D}}\Big(\frac{\sum_{\z \in \textbf{D}_{\Bar{S}}(\z,\x)}\mathbf{1}_{\{f(\z)\neq f(\x)\}}}{\textbf{D}_{\Bar{S}}(\z,\x)}\Big)
\end{aligned}
\end{equation}


where we denote $\textbf{D}_{\Bar{S}}(\z,\x):=\{\z\in\mathbf{D}:\z_{\Bar{S}}=\x_{\Bar{S}}\}$. We now construct a ReLU neural network $f$. Before doing so, we specify the functional form that we intend the constructed neural network $f$ to implement. For each set $S_j$, we introduce two vectors $\x^j$ and $\y^j$, both belonging to $\mathbb{Z}^n$ and drawn from the dataset $\mathbf{D}$, such that they agree on all coordinates outside $S_j$; that is, they share identical feature values on the complement $\bar{S_j}$. Formally:


%where we denote $\textbf{D}_{\Bar{S}}(\z,\x):=\{\z\in\mathbf{D}:\z_{\Bar{S}}=\x_{\Bar{S}}\}$. We construct a neural network $f$ with ReLU activations. We first, however, will denote the functional form we would want our constructed function $f$ to take. For each $S_j$ we will define two new vectors $\x^j$ and $\y^j$ both vectors in $\mathbb{Z}^n$, within our dataset $\mathbf{D}$ such that they ``share'' the same features within the $\Bar{S_j}$ features, or more specifically:

\begin{equation}
    \x^j_{\Bar{S_j}} = \y^j_{\Bar{S_j}}
\end{equation}

and we assign strictly different values to the features in $S_j$ by selecting them from two distinct vectors $\z$ and $\z'$ such that:

%and will set strictly different assignments for the features in $S_j$ by choosing them from different vectors $\z$ and $\z'$ such that it holds that:

\begin{equation}
    \x^j=(\x^j_{\Bar{S_j}};\z^j_{S_j}) \neq (\x^j_{\Bar{S_j}};\z'^j_{S_j})= (\y^j_{\Bar{S_j}};\z'^j_{S_j})=\y^j
\end{equation}

We aim to construct $f$ such that it satisfies the following:

%We will want to construct $f$ so that it will satisfy that:

\begin{equation}
f(\x^j)=f(\x^j_{\Bar{S_j}};\z^j_{S_j}) \neq f(\x^j_{\Bar{S_j}};\z'^j_{S_j})= f(\y^j_{\Bar{S_j}};\z'^j_{S_j})=f(\y^j)
\end{equation}

%This can be done for instance by simply choosing the function $f$ such that $f(\x^j)=0$ and $f(\y^j)=1$. More particularly we will choose one specific 
%We in addition to these two features will add all of the $\{0,1\}^{|S|}$ combinations over 

Let us now assume that both $\x^j$ and $\y^j$ occur \emph{uniquely} in the dataset $\mathbf{D}$ across all $[n]$ coordinates --- that is, no other vector in $\mathbf{D}$ matches either of them on any coordinate. Under this assumption, we obtain the following relation:

%Let us assume now that both $\x^j$ and $\y^j$ appear \emph{uniquely} in the dataset $\textbf{D}$ (meaning there is no other vector in the dataset that is equal to them. Then we get the following relation:

\begin{equation}
\begin{aligned}
\label{to_ref_neq_equation}
\mathbb{E}_{\x\sim\mathcal{D}}[\textbf{Pr}_{\z\sim \mathcal{D}}(f(\z) \neq f(\x) \ | \  \z_{\Bar{S_j}}= \x_{\Bar{S_j}})]=\\
\frac{1}{\lvert \textbf{D} \rvert}\sum_{\x\in\textbf{D}}\Big(\frac{\sum_{\z \in \textbf{D}_{\Bar{S_j}}(\z,\x)}\mathbf{1}_{\{f(\z)\neq f(\x)\}}}{\textbf{D}_{\Bar{S_j}}(\z,\x)}\Big)=\frac{2}{2\lvert\textbf{D}\rvert}=\frac{1}{\lvert\textbf{D}\rvert}
\end{aligned}
\end{equation}

Here, the factor of $2$ appearing in both the numerator and denominator arises from the fact that only two vector pairs --- $(\x^j, \z^j)$ and $(\z^j, \x^j)$ --- contribute to the respective sums.

%Where the $2$ obtained in both the nominator and the denominator are the result of only two pairs of vecotrs $(\x^j,\z^j)$, and $(\z^j,\x^j)$ that are counted in the corresponding sums.


We repeat this construction for each $S_i$ in the set-cover instance, each time choosing two \emph{distinct} vectors $\x^i$ and $\y^i$ that differ on \emph{every} coordinate in $[n]$. In total, this yields $2m$ vectors that form the dataset $ \mathbf{D}$ from which we construct $ f$.



%We perform this same exact process for every $S_i$ in our set cover instance, each time choosing two \emph{distinct} vectors $\x^i$ and $\y^i$. Overall, we construct a set of $2m$ vecotrs in the dataset $\textbf{D}$ from which we construct $f$.


We then construct a neural network $f$ that simulates this function. This is achieved using a known result on \emph{neural network memorization}~\citep{vardi2022optimal}, which guarantees that a ReLU network capable of representing any finite set of input–output pairs can be built in polynomial time. Using this result, we construct $f$ accordingly, and set $\delta' := \frac{\delta}{|\mathbf{D}|}$.

We now argue that a subset $S_j$ is a global contrastive reason for $f$ under tolerance $\delta'$ if and only if $C$ is a partial set cover for $U$ under tolerance $\delta$. By Equation~\ref{to_ref_neq_equation}, each constructed pair of vectors $(\x^j, \y^j)$ contributes exactly $\frac{1}{|\mathbf{D}|}$ to the value term. Moreover, because each $S_j$ uses its own distinct pair of vectors in $\mathbb{Z}^n$, there is no overlap between these pairs. Consequently, the coverage achieved by any set cover $C$ corresponds exactly to the accumulated contribution of the selected subsets in the input domain of $f$. This completes the reduction.


%We then construct a neural network $f$ that will simulate this function: this can be done via a known result regarding \emph{neural network memorization}~\citep{vardi2022optimal}, where it is specifically a known result that one can construct a neural network in polynomial time to simulate any finite number of points. We hence construct $f$ in this way, and define via our construction $\delta':=\frac{\delta}{|\mathbf{D}|}$. We now claim that a subset $S_j$ is a global contrastive reason for $f$ under the tolerance $\delta'$ if and only if $C$ is a partial set cover for $U$ with the tolerance $\delta$. This holds true since due to the reason we have shown in Equation~\ref{to_ref_neq_equation}, every constructed pairs of vectors $\x^j,\y^j$ corresponds to an exact addition term of $\frac{1}{|\mathbf{D}|}$. Since for each $S_j$ we construct a distinct set of vectors in $\mathbb{Z}$, then there is no ``overlapp'' between such vectors so the coverageof any set cover $C$ would be equal to the coverage of any subset that is chosen within the input domain of $f$. This concludes the reduction.


$\qedsymbol{}$



%The problem: it may hold that

%\begin{equation}
%    Pr(\x)\neq Pr(\x_1)\cdot Pr(\x_2)\cdot \ldots Pr(\x_n)
%\end{equation}


%?

%which does not show tightness of approximation when assuming feature indpendence.

\section{Tractability Results for Orthogonal DNFs} \label{app:odnf}

\begin{lemma}
If $f$ is an orthogonal DNFs and the probability term $v^{g}_{\text{suff}}$ can be computed in polynomial time given the distribution $\mathcal{D}$ 
(which holds for independent distributions, among others),
then obtaining a subset-minimal \emph{global} $ \delta $-sufficient reason can be obtained in \emph{polynomial time}.
\end{lemma}
\begin{proof}
We begin by defining orthogonal DNFs~\citep{crama2011boolean} before presenting the full proof.
%
\begin{definition}[Orthogonal DNFs]
A Boolean function $\varphi$ is in \emph{orthogonal disjunctive normal form} (orthogonal DNF)~\citep{crama2011boolean}
if it is a disjunction of terms $T_1 \lor T_2 \lor \dots \lor T_m$, such that for every pair of distinct terms $T_i$ and $T_j$ ($i \neq j$):
\[
T_i \land T_j \models \bot
\]
This means that no single variable assignment can satisfy more than one term simultaneously.
\end{definition}
%
The proof mirrors the argument used in Theorem~\ref{subset_minimal_dt_theor}, relying on the fact that $v^{g}_{\text{suff}}$ can be computed in polynomial time for any set $S$.

For a given set $S$, let $\mathcal{C}^{S}{T_i}$ be the constraints that the term $T_i$ imposes on the variables in $S$. We enumerate all pairs of terms $(T_i, T_j)$. Any pair whose constraints on $S$ are \emph{inconsistent}
(i.e., $\mathcal{C}^{S}_{T_i} \land \mathcal{C}^{S}_{T_j} \models \bot$) is ignored, since it contributes zero to $v^{g}{\text{suff}}$.


%The proof follows the same reasoning as Theorem~\ref{subset_minimal_dt_theor} by showing that computing 
%$v^{g}_{\text{suff}}$ takes polynomial time for an arbitrary set $S$.
%
%For an arbitrary set $S$, let $\mathcal{C}^{S}_{T_i}$ denote the constraints on features in $S$ imposed by the term $T_i$.
%We enumerate all pairs of terms $(T_i, T_j)$. Pairs with inconsistent constraints on $S$ 
%(i.e., $\mathcal{C}^{S}_{T_i} \land \mathcal{C}^{S}_{T_j} \models \bot$) are discarded, as they contribute 0 to $v^{g}_{\text{suff}}$.
%

For each \emph{consistent} pair of terms $(T_i, T_j)$, we evaluate
$\mathbf{Pr}(\x_{S} \in \mathcal{C}^{S}_{T_i, T_j})$, where
$\mathcal{C}^{S}_{T_i, T_j} = \mathcal{C}^{S}_{T_i} \cap \mathcal{C}^{S}_{T_j}$.
The contribution of this pair to $v^{g}_{\text{suff}}$ is then
\[
\mathbf{Pr}(\x_{S} \in \mathcal{C}^{S}_{T_i, T_j}) \cdot \mathbf{Pr}(\x_{\bar{S}} \in \mathcal{C}^{\bar{S}}_{T_i}) \cdot \mathbf{Pr}(\x_{\bar{S}} \in \mathcal{C}^{\bar{S}}_{T_j}).
\]

Under feature independence, these probabilities factor as
$\Pr(\x_{S}) = \prod_{i\in S} p(x_i)$.
Since there are at most $m^2$ such term pairs, the value of $v^{g}_{\text{suff}}$
can be computed in polynomial time under an independent distribution $\mathcal{D}$.
%For each consistent pair $(T_i, T_j)$, we compute $\mathbf{Pr}(\x_{S} \in \mathcal{C}^{S}_{T_i, T_j})$,
%where $\mathcal{C}^{S}_{T_i, T_j} = \mathcal{C}^{S}_{T_i} \cap \mathcal{C}^{S}_{T_j}$,
%then the contribution of $(T_i, T_j)$ to $v^{g}_{\text{suff}}$ is:
%\[
%\mathbf{Pr}(\x_{S} \in \mathcal{C}^{S}_{T_i, T_j}) \cdot \mathbf{Pr}(\x_{\bar{S}} \in \mathcal{C}^{\bar{S}}_{T_i}) \cdot \mathbf{Pr}(\x_{\bar{S}} \in \mathcal{C}^{\bar{S}}_{T_j}).
%\]
%
%Since features are independent, $\mbf{Pr}(\x_{S}) = \prod_{i\in S}p(x_i)$.
%There are in total $m^2$ pairs of terms, so the probability term $v^{g}_{\text{suff}}$
%can be computed in polynomial time given the independent distribution $\mathcal{D}$.
\end{proof}



\section{Disclosure: Usage of LLms}
\label{app:llm-disclosure}
An LLM was used exclusively as a writing assistant to refine grammar and typos and improve clarity. It did not contribute to the generation of research ideas, study design, data analysis, or interpretation of results, all of which were carried out solely by the authors.

\end{document}



